\documentclass[11pt]{article}
\usepackage[T1]{fontenc}
\usepackage{lmodern,amsmath,amssymb,amsthm,mathtools,booktabs,array,graphicx}
\usepackage[margin=1.05in]{geometry}
\usepackage{xcolor}
\usepackage{xr-hyper}
\usepackage[colorlinks=true,linkcolor=blue!45!black,citecolor=blue!45!black,urlcolor=blue!45!black,filecolor=blue!45!black,bookmarksnumbered=true]{hyperref}
\externaldocument[M-]{friend-tree-exponent}[friend-tree-exponent.pdf]
\newcommand{\proofsrc}[2]{\href{https://github.com/uzay-eth/friend-tree-exponent/blob/a0972085ad72ad5685ae7b424988648d1965e73c/#1}{#2}}
\newcommand{\ordersrc}[2]{\href{https://github.com/uzay-eth/friend-tree-exponent/blob/62081cad01112c42c2011ac3eed11aa3aa665e09/#1}{#2}}
\newtheorem{theorem}{Theorem}
\newcommand{\resulttag}{?}
\newcommand{\tagnext}[1]{\renewcommand{\resulttag}{#1}}
\renewcommand{\thetheorem}{\resulttag}
\renewcommand{\theHtheorem}{\resulttag}
\newcommand{\sectag}[1]{\renewcommand{\thesection}{#1}\renewcommand{\theHsection}{#1}}
\renewcommand{\thesubsection}{\thesection.\alph{subsection}}
\newtheorem{lemma}{Lemma}[section]
\newtheorem{proposition}[lemma]{Proposition}
\newtheorem*{couplingfact}{Coupling principle}
\newtheorem*{marginalfact}{Marginal site bound}
\theoremstyle{remark}
\newtheorem*{remark}{Remark}
\numberwithin{equation}{section}
\newcommand{\E}{\mathbb E}
\newcommand{\Prob}{\mathbb P}
\newcommand{\ind}{\mathbf 1}
\newcommand{\res}[1]{[#1]_p}
\newcommand{\ftbreak}{\\}% a layout-only line break inside a transcribed display
\newcommand{\reviewnote}[1]{\par\nopagebreak\smallskip\noindent\fbox{\parbox{\dimexpr\linewidth-2\fboxsep-2\fboxrule}{\small\textit{Review note.} #1}}\par\smallskip}
\newsavebox{\ftfitbox}
\newcommand{\fitdisplay}[1]{\[\sbox{\ftfitbox}{$\displaystyle #1$}\ifdim\wd\ftfitbox>\linewidth\typeout{FITDISPLAY page \thepage\space width \the\wd\ftfitbox\space line \the\linewidth}\resizebox{\linewidth}{!}{\usebox{\ftfitbox}}\else\usebox{\ftfitbox}\fi\]}
\newcommand{\trhead}[1]{\par\smallskip\noindent\textbf{#1}\par\nopagebreak\smallskip}
\newenvironment{transcription}{\par\medskip\noindent\rule{\linewidth}{0.4pt}\par\nopagebreak}{\par\nopagebreak\noindent\rule{\linewidth}{0.4pt}\par\medskip}
\makeatletter
\renewcommand*\l@section{\@dottedtocline{1}{0pt}{4.2em}}
\renewcommand*\l@subsection{\@dottedtocline{2}{4.2em}{5.2em}}
\makeatother
\setlength{\emergencystretch}{3em}
\hypersetup{pdftitle={The non-leaf growth exponent of the random friend tree: supplement},pdfauthor={Can Uzay},pdfsubject={Complete proofs for the random friend tree paper: growth, birth order and further material},pdfkeywords={random friend tree, growth exponent, drift certificate, stochastic order, computer-assisted proof}}
\title{The non-leaf growth exponent of the random friend tree\\[4pt]\large Supplement: complete proofs}
\author{Can Uzay\\\small University of Nevada, Reno\\\small\href{mailto:hi@canuzay.com}{hi@canuzay.com}}
\date{25 November 2026}
\begin{document}
\maketitle
\begin{abstract}
This supplement accompanies the paper \emph{The non-leaf growth exponent of the random friend tree} \cite{Main}. Part~I proves the growth results of the paper: the common exponent and its finite-seed construction, the three- and five-vertex reductions behind the numerical enclosure, the marked kernel, the random amplitude and its density, the degree law, and unit hub mass. Part~II is a self-contained module on birth order. It proves that the two founders' limiting degree shares agree in law, that from label two onward the shares decrease strictly in first-order stochastic order, and quantitative forms for hub probabilities and positive moments. Part~III records results with narrower scope or conditional hypotheses, the provenance of every proof, and the verification packages. The theorem that the mean core size is not P-recursive is proved in Appendix~\ref{M-app:nonholonomic} of the paper.
\end{abstract}

\section*{How to use this supplement}\addcontentsline{toc}{section}{How to use this supplement}
\paragraph{Tags.} Every result carries a tag, such as Theorem~\ref{th:gold2}, and keeps it in the paper, here and in the repository. Each section of this supplement is named by the tag of the result it proves, and its lemmas, propositions and equations are numbered within that section, as in Lemma~\ref{lem:classwise} or \eqref{eq:core-kernel}. Theorems are restated here with exactly the wording of the paper.

\paragraph{Transcribed proofs.} Four parts of the argument were written as Markdown notes in the repository and reviewed in that form: the three-vertex reduction (Section~\ref{sec:mot3}), the five-vertex reduction (Section~\ref{sec:enc-reduction}) and the two halves of the complete proof of Theorem~\ref{th:stoch-order} (Section~\ref{sec:ord-proof}). They are reproduced here between horizontal rules. Each one is introduced by its source path, commit and SHA-256. The text has been converted mechanically to \LaTeX{}. Every formula has the same content as in the source: a wide display is either scaled to the line width or broken across lines by a layout-only command, and nothing else in a formula changes. The repository's transcription check compares every formula with the source after removing those documented layout commands, as well as the words of the prose and the list numbers. Inside a transcription, numbers such as ``Lemma~3.3'' or ``\S5.2'' refer to the transcription's own headings, which are printed with their original numbers. Verification records, binary formats and replay logs of the notes are not reproduced; the verification packages are described in Section~\ref{app:replay}.

\paragraph{Computation.} Where a proof reduces every finite tree to finitely many inequalities, the reduction is proved here, and exact programs in the repository check the inequalities. A passing run checks those inequalities, not the reduction. The audits of the five-vertex row verifier and of the two stochastic-order certificates replayed the authored programs and checked selected components independently; neither has been reimplemented in full. None of the proofs has had a human referee.

\tableofcontents

\part{Growth}

\sectag{EXP}
\section{Growth from finite seeds}\label{sec:exist}
\tagnext{EXP}
\begin{theorem}[Existence and the first enclosure]\label{th:exist}
From one edge,
\[
 \lim_{n\to\infty}\frac{\log\E C_n}{\log n}=\mu,\qquad
 \lim_{n\to\infty}\frac{\log C_n}{\log n}=\mu\quad\text{almost surely},
 \qquad 113/200\le\mu\le1133/2000<17/30 .
\]
\end{theorem}
A branch separated by one edge almost grows freely if the two endpoints seldom send across the cut. Planting a fixed branch at many core sites makes its mean growth repeatable. Disjoint sites let one repeat the lower bound on sample paths. Subsections~EXP.a--EXP.d below prove both limits. For the first enclosure, the \proofsrc{proofs/bracket-motifs.md}{three-vertex proof} (Section~\ref{sec:mot3}) constructs two potentials with the one-sided comparisons $\Phi_+\ge c_0 C_n$ and $\Phi_-\le C_0 C_n$ for fixed $c_0,C_0>0$. Both witnesses have positive value at $T_3$; the lower witness need not be nonnegative on every tree. The \proofsrc{verify/motif-bracket/cert_verify_gmp.c}{integer verifier} checks vertex, edge and triple weight rows and bounds the unbounded degree tail. For every tree at size $n\ge3$, it certifies $n\E[\Delta\Phi_-\mid T_n]\ge(113/200)\Phi_-$ and $n\E[\Delta\Phi_+\mid T_n]\le(1133/2000)\Phi_+$. Taking expectations and iterating from $T_3$ gives
\[
 \frac{\Phi_-(T_3)}{C_0}\prod_{j=3}^{n-1}\left(1+\frac{113}{200j}\right)
 \le \E C_n \le
 \frac{\Phi_+(T_3)}{c_0}\prod_{j=3}^{n-1}\left(1+\frac{1133}{2000j}\right).
\]
These products are constant multiples of their respective powers of $n$ up to factors tending to one. Taking the logarithmic mean limit gives the enclosure. The upper witness alone also gives the almost-sure upper exponent: for each $\epsilon>0$, apply Markov's inequality and Borel--Cantelli at $n=2^j$ with threshold $n^{1133/2000+\epsilon}$, then use monotonicity between dyadic sizes and send $\epsilon\downarrow0$. This needs neither the lower witness nor existence of the limit. The gap is $17/30-1133/2000=1/6000$. The finite-size prefactors and an independent degree-band upper bound $5657/9983$ appear in Theorem~1.2 of the author's earlier preprint \cite{UzayFirst}.
Give each vertex a rate-one Poisson clock. At a ring the source chooses a uniform neighbour and a new leaf attaches there. Each newborn starts an independent clock. Conditional on the tree at size $j$, the next waiting time is exponential of rate $j$, independent of the choice of source and of its neighbour. The jump chain is therefore the prescribed discrete chain, independent of the sequence of waiting times. The sum of the waiting times diverges almost surely: their means have divergent sum and their centred partial sums converge in $L^2$ and almost surely by the summability of $j^{-2}$.

\begin{couplingfact}
A construction with the prescribed conditional law of each waiting time, source
and neighbour given the revealed past has the law of the friend-tree process.
At a processed ring or deterministic time, remaining clocks may be sampled afresh.
\end{couplingfact}
\begin{proof}
The independent exponential clocks have independent, mean-one residual waits at such times by the strong Markov property. A fresh uniform mark gives the stated choice of neighbour. Induction on the jump index gives the law through each finite number of births, and non-explosion gives the law on bounded time intervals.
\end{proof}
A neighbour ordering chosen from the past may be used with a fresh
uniform mark. To choose among $d+1$ neighbours one may also make a
uniform choice among the first $d$, accepting it with independent probability
$d/(d+1)$ and otherwise taking the last one.

Use $N(t),C(t)$ for the population and core count in continuous time. The population from $k$ individuals is the sum of $k$ independent Yule families, by assigning each birth to its ringing source. In one family the probabilities $p_j(t)=e^{-t}(1-e^{-t})^{j-1}$ solve $p_j'=(j-1)p_{j-1}-jp_j$ with initial size one. Thus $\E N(t)=ke^t$ and $\E N(t)^2=k(k+1)e^{2t}-ke^t$. Put $F(t)=\E C(t)$ starting from the edge. From a specified seed $T$ write $C_T(t)$ for the core count.

\subsection{Sites and preparation}
At a core vertex $u$ let $c_u$ be its number of core neighbours, $L_u$ its number of leaf neighbours and $D_u=c_u+L_u$. Call a vertex with $(c,L)=(2,0)$ bare. An eligible core vertex has either $c\le2,L\ge1$, or is bare and adjacent to another bare vertex. Here and below $C$ denotes the current core count.
\begin{lemma}\label{lem:eligible} At least $C/4$ core vertices are eligible.\end{lemma}
\begin{proof}
For $C=0$ the assertion is empty. For $C=1$ this holds at the centre of the star. For $C\ge2$, the $I$ bare vertices without bare neighbours have $2I\le C-1$ distinct incident core edges. The number of core vertices of core degree at least three is at most the number of ends of the core tree, by the degree sum. Each end has a leaf. If $E$ vertices are eligible, $C\le E+I+E$. Since $I\le(C-1)/2$, this gives $E\ge C/4$.\end{proof}

\begin{lemma}\label{lem:plant} Fix integers $M,k\ge2$, a rooted $k$-vertex seed $(T,v)$ and any old eligible $u$. There are $h=k+1$ and $p=p(M,k)>0$ such that at the end of $h$ units a branch attached to $u$ has shape $(T,v)$, with $v$ adjacent to $u$ and $D_u\ge M$, with conditional probability at least $p$ in every initial tree. This branch contains no vertex that was core at the start of the interval.
\end{lemma}
\begin{proof} Use independent rate-one clocks and fresh uniform marks, as above.

First keep $u$ silent for one unit. If $u$ has a leaf, tag one; it stays a leaf. Otherwise take a bare neighbour $w$ and its other neighbour $z$. Keep $z$ silent too and require $w$'s first ring in the unit to choose $u$. Before this ring $w$ still has degree two, since only $u,z$ could cause it to gain. Tag the new leaf at $u$. The probability of this first stage is at least $e^{-2}(1-e^{-1})/2$. This lower bound also applies in the first case, where silence alone has probability $e^{-1}$.

For one further unit keep $u$ silent and require at least $M$ rings of its tagged leaf. The conditional probability of this stage is $e^{-1}\Prob(\mathrm{Pois}(1)\ge M)$. These rings give $u$ at least $M+1$ leaves, counting the tagged one, and still $c_u\le2$. Only a ring of $u$ could promote its leaf.

During the third unit require exactly one ring of $u$. Until that ring $u$ has at least $M+1$ leaves and at most two core neighbours. Hence $L_u/D_u\ge(M+1)/(M+3)>M/(M+2)$, so with leaves ordered first a mark below $M/(M+2)$ forces a promotion at $u$. Call the promoted vertex $v$ and require silence of $v$ and its new leaf until the end of the unit. This last conditional probability is at least $e^{-2}$. Additional leaf arrivals at $u$ cannot spoil these requirements.

Next silence $u$ for $k-2$ units. Beginning with a prescribed edge incident to $v$ in $T$, adjoin the remaining vertices in parent-before-child order. In each unit give an internal neighbour of the desired recipient exactly one ring and require silence of the other current branch vertices and of the newborn after the ring. With at most $k-1$ branch vertices before this birth, the ringing source has degree at most $k-1$ including the exterior edge. The weaker bound $1/(k+1)$ suffices for its choice probability. There are at most $k-1$ current branch vertices, together with $u$ and the newborn. Including the silence of $u$, these requirements have conditional probability at least $e^{-(k+1)}/(k+1)$ per unit. Thus we may take
\[
p=\frac{e^{-2}(1-e^{-1})}{2}\,e^{-1}\Prob(\mathrm{Pois}(1)\ge M)
   \, e^{-3}\frac{M}{M+2}
    \left(\frac{e^{-(k+1)}}{k+1}\right)^{k-2}.
\] All effects of the exterior during the last stage must pass through the silent vertex $u$. The promoted leaf was a leaf or unborn at the beginning of the interval, and every other vertex of the branch was born later.\end{proof}

When preparations succeed at two distinct old core vertices, the two resulting branches stay disjoint. Indeed, the connected subtree on the old cores persists. Orient edges away from any old core. Each cut edge points out of this subtree; neither cut can lie below the other since its parent is an old core. Subsequent additions preserve these orientations.

\subsection{Comparison across one cut edge}
Compare the branch over the next $s$ units with a freestanding copy of $T$ driven by fresh independent clocks and internal choice marks. Before the first discrepancy use the rings and internal choices of the free tree in the branch.

At a ring of its root with internal degree $d$, accept its internal choice with probability $d/(d+1)$ using an independent mark. Otherwise redirect to $u$. In exterior rings of $u$, order the branch root first. Since $D_u\ge M$ throughout the comparison, the absence of a mark in $[0,1/M]$ at $u$ suffices to prevent inward redirection. At the first discrepancy make the prescribed choice in the actual tree (including the birth), then give the branch fresh drivers. The free reference keeps its drivers throughout. A choice mark at $u$ in the test interval need not produce a discrepancy: the actual probability of choosing the root is $1/D_u$. The construction gives uniform actual choices: at the branch root each internal neighbour and the exterior neighbour has probability $1/(d+1)$.

Marks in $[0,1/M]$ at $u$ form a Poisson process of rate $1/M$ independent of the free reference. Its silence and acceptance of every root choice imply agreement. Given the free ring times, sources and chosen neighbours (the free history), the unused acceptance marks remain independent uniforms. The two tests then have joint probability
\[
e^{-s/M}\prod_{\text{rings at }v}d/(d+1).
\]
On agreement all reference core vertices are actual ones on the far side of the cut. In particular, writing $X$ for the actual number there at the end and $d_j$ for the internal degree just before the $j$th root ring, we have, conditional on the whole tree at the end of a successful preparation,
\begin{equation}\label{eq:cut}
 \E[X\mid \mathcal F_{t+h}]
 \;\ge\; e^{-s/M}\,\E\!\left[C_T(s)\prod_j\frac{d_j}{d_j+1}\right].
\end{equation}
Here $t$ is the start of preparation, and $\mathcal F_r$ records the tree, processed rings and marks, and auxiliary tests revealed through time $r$. Future increments of a reserved source are revealed only as time advances. The expectation on the right uses only the free process. The product retains the dependence between growth and the number and timing of rings at the root.

Take $A=8$, $0<\epsilon<1$ and $b=7/\log(1/\epsilon)$. We have $b>\epsilon$ since $\epsilon\log(1/\epsilon)\le1/e$, and
\[
 b\log(b/\epsilon)-b+\epsilon=7+b\log b-b+\epsilon\ge6 .
\] Also $A\log A-A+1>5$. Take $k\ge6M/\epsilon$. A good reference history has each original clock ringing at most $bs$ times up to $\epsilon s$ and at most $As$ times up to $s$, and some original vertex having degree at least $M$ by $\epsilon s$. The Poisson upper-tail bound $\Prob(\mathrm{Pois}(as)\ge xs)\le\exp[-s(x\log(x/a)-x+a)]$ follows by the exponential Markov inequality. It handles all $2k$ clock bounds.

In the free process, order one original neighbour first at each original vertex (the seed is connected and $k\ge2$). Until some original degree reaches $M$, a mark below $1/M$ increases an original degree. If none reaches $M$, at most $kM$ such marks have occurred. Their number without stopping has Poisson law of mean $k\epsilon s/M\ge6s$. Its lower tail at the fixed $kM$ is $O(e^{-6s}s^{kM})$, at most $e^{-4s}$ for large $s$.

Thus the exceptional probability is $O_k(e^{-4s})$. By Cauchy--Schwarz and the Yule moment formula its contribution to $\E C_T(s)$ is $O_k(e^{-s})$. All bounds hold for $s\ge s_0(k,M,\epsilon)$ uniformly over the finitely many $k$-vertex seeds. Here $M,\epsilon,k$ are fixed before $s$ grows.

Assign each good history to the first numbered original vertex whose degree is at least $M$ at $\epsilon s$. For a history assigned to $v$, before $\epsilon s$ each factor in the comparison product is at least $1/2$; subsequently it is at least $e^{-1/M}$ by $\log(1+1/d)\le1/d$. The comparison weight is at least $e^{-\eta s}$, where $\eta=b\log2+(A+1)/M$. Some fixed original $v$ accounts for at least $(\E C_T(s)-O_k(e^{-s}))/k$ of the expected good core count.

\subsection{Iteration and the mean exponent}
Stop a process from the edge at its first visit to size $k$. Before that visit there are at most $k$ core vertices. After it, the strong Markov property and monotonicity in elapsed time give $F(s)\le k+\max_{|T|=k}\E C_T(s)$. Fix a maximizing seed and a maximizing root as just described, before the preparations.

\begin{marginalfact}
For all sufficiently large $s$ and every $t\ge0$, the following bound holds:
\[
F(t+h+s)\ge F(t)\frac{p}{4k}e^{-\eta s}
                    \big(F(s)-k-O_k(e^{-s})\big).
\]
\end{marginalfact}
\begin{proof}
For each old core vertex $u$, define $X_u$ using the actual tree path alone. At time $t+h$, take every edge $uw$ whose component on the $w$ side contains no vertex that was core at $t$. Count all vertices on those sides that are core at time $t+h+s$. Components used for different edges are disjoint: they are the outward components attached to the connected subtree of old cores, and later leaf additions cannot join them. Hence $\sum_u X_u\le C(t+h+s)$ on each path.

For a particular eligible $u$ we may evaluate $\E[X_u\mid\mathcal F_t]$ using the preparation and coupling for that site. On success its planted branch is one of the components counted by $X_u$. Success is determined at $t+h$, so the weighted lower bound \eqref{eq:cut} applies conditionally there. The marginal continuation preserves the tree law for each $u$. We may therefore sum the resulting bounds for these path functionals even when the clock requirements at two sites are incompatible. Lemma~\ref{lem:eligible} supplies at least $C(t)/4$ eligible sites, and taking expectations proves the bound.
\end{proof}

After the first birth the promotion rate $\sum_u L_u/D_u$ is at least one: a single core contributes one, and otherwise the two ends of the core each contribute at least one half. Writing $\tau$ for the first birth time and integrating the intensity gives $F(t)\ge\E[(t-\tau)^+]\to\infty$; the intensity is bounded on $[0,t]$ by $N(t)$, which is integrable.

For fixed large $s$ the last parenthesis is at least $F(s)/2$, so the multiplier $B_s$ of $F(t)$ is positive. Iterate at $t=1+j(h+s)$ and interpolate by monotonicity to obtain $\liminf t^{-1}\log F(t)\ge\log B_s/(h+s)$. Take $s$ along a sequence attaining the upper limit. The two limiting rates differ by at most $\eta$. Sending $M\to\infty$ and $\epsilon\downarrow0$ gives a common limit $\mu\in[0,1]$, with the upper bound from $F(t)\le2e^t$. The block length $s$ tends to infinity for fixed $M,\epsilon,k$ before the loss $\eta$ is sent to zero.

Write $a_n=\E C_n$. Independence gives $F(t)=\E a_{N(t)}$ and $a_n$ is nondecreasing. At $t=\log n$ a single original Yule family has probability $(1-1/n)^{n-1}\ge e^{-1}$ of size at least $n$. Thus $a_n\le eF(\log n)$.

Conversely let $t=(1-\alpha)\log n$ for $0<\alpha<1$. Since $N$ is a sum of two geometric variables, $\Prob(N>n)\le2(1-e^{-t})^{\lfloor n/2\rfloor}$. With the second moment formula this gives $\E[N;N>n]=o(1)$. Hence $F(t)\le a_n+o(1)$ and $\liminf\log a_n/\log n\ge(1-\alpha)\mu$. Let $\alpha\downarrow0$.

\subsection{Independent contributions and the sample-path exponent}
\begin{lemma}[Independent contributions in one block]\label{lem:block}
Fix the preparation parameters and a sufficiently large $s$. From any tree with $C\ge1$ core vertices there is a coupling in which the core count after $h+s$ time units is at least a sum of $m\ge C/6$ independent copies of a nonnegative variable $Y$. Its law depends only on the block parameters, and
\[
 \E Y\ge(p/k)e^{-\eta s}\big(F(s)-k-O_k(e^{-s})\big),\qquad
 \E Y^2\le k(k+1)e^{2s}.
\]
\end{lemma}
\begin{proof}
Select disjoint clock sets. With $C\ge2$, let $A_0$ count sites with $c\le2,L\ge1$. The number of non-bare core vertices is at most $2A_0$, by the core degree sum. Suppressing maximal paths of bare vertices gives a tree on those non-bare vertices. There are at most $2A_0-1$ such paths. Split each into disjoint triples leaving at most two vertices per path. For $J$ triples, $C\le6A_0+3J$. Take all $A_0$ sites and one end of each triple. This gives at least $C/6$ sites. For $C=1$ take the centre.

At a site with a leaf reserve its clock and that of the tagged leaf; at a triple reserve its three clocks, assigning $(u,w,z)$ in path order. All sets are disjoint. Use the four stages above: signals from outside may add leaves but cannot decrease $L_u$ while $u$ is silent. The first two stages can be decided from independent site-specific clocks, including new clocks for tagged leaves at triples.

In the third stage take the source ring and its mark below the fixed threshold from the site's own random inputs. Which leaf is selected may depend on the surroundings. A leaf of $u$ has no other neighbour, so it belongs to no other site's reserved vertices. Even when it is the original tagged leaf, the second stage has finished using that clock. At selection time its future clock can be sampled afresh from the site's independent randomness, as can that of the newborn. Given the ring time within the unit, $\tau\in[0,1]$, the two-clock silence probability is $e^{-2(1-\tau)}$; thin it by an independent coin of probability $e^{-2\tau}$.

For the final stage use fresh independent branch clocks and choices under the silent exterior parent. To obtain a common flag, proceed as follows. At a leaf site thin stage-one silence by a coin of probability $e^{-1}(1-e^{-1})/2$, giving the same probability as at a triple. In each final-stage unit with $j$ branch vertices, rank the required recipient first and test the source mark against $1/(k+1)$. Require exactly one ring of that source and silence of the other $j-1$ branch clocks and of $u$. For a newborn at time $\tau$ within the unit, require silence up to the end and an independent coin of probability $e^{-\tau}$. The newborn test then has probability $e^{-1}$ for every birth time. The probability of this unit's flag is $e^{-(j+2)}/(k+1)$; thin by $e^{-(k-1-j)}$ to obtain $e^{-(k+1)}/(k+1)$, since $2\le j\le k-1$. All coins and scripted clocks can be drawn from the given site's own randomness. This realizes the displayed preparation probability $p$ at every site, even when outside rings affect the identity of the promoted leaf. Equivalently, the post-selection clocks can be assigned beforehand by site and revealed only at selection, preserving the exponential waiting times. Flags are evaluated from the reserved inputs. The actual process uses each prescribed ranking and fresh driver only when its stage is reached, and processes a failed test with the ordinary transition law before continuing.

Independently of preparation, generate at every site a free reference history, acceptance marks and the exterior test on the site's own $u$ clock. Use its reference core count as $Y$ if preparation succeeds, the history is good and assigned to the prescribed root and both tests succeed; otherwise put $Y=0$. The old vertices whose neighbour orderings are prescribed are distinct across sites.

Drive each prepared branch by its reference up to the first discrepancy at its own cut. Since the retained flags all have probability $p$, $Y$ has the same law at a site with a leaf and at a triple. The marginal comparison \eqref{eq:cut} alone gives no independence between sites. Here it is the retained flags and reference histories that use disjoint randomness, not the actual branch paths: conditional on the block-start history the flags and reference histories at distinct selected sites are independent. On each retained agreement the reference cores occupy that site's actual branch, even though the branches share the actual exterior. Thus, given the whole past at the start of the block, the endpoint core count dominates the sum of independent copies of $Y$, with the asserted mean by the good-history estimate. Since $Y$ is bounded by its free population, the Yule second moment gives the second bound.
\end{proof}

For $m\ge C/6$ disjoint sites their sum has conditional mean $m\E Y$ and variance at most $m\E Y^2$. Chebyshev's inequality therefore bounds the probability of fewer than $C\E Y/12$ endpoint cores by $4\E Y^2/(m(\E Y)^2)$, whenever $\E Y>0$. The estimate is conditional on the tree at the start of the block and uniform in trees of core size $C$.

After size three, the conditional probability of a promotion is at least $1/n$ at size $n$. Hence from any size $n_0\ge3$ the conditional probability of no promotion through the step to $n_1$ is at most
\[
  \prod_{n=n_0}^{n_1-1}(1-1/n)=\frac{n_0-1}{n_1-1}.
\]
Sending $n_1\to\infty$ and taking the union over $n_0$ gives $C(t)\to\infty$ almost surely.

Fix a block for which $g=\E Y/12>1$, write $t_j=j(h+s)$, and let $I_j$ indicate $C(t_j)<g C(t_{j-1})$. Since the bound above applies starting from every tree with a positive core count, the Markov property gives
\[
 q_j=\Prob(I_j=1\mid\mathcal F_{t_{j-1}})
 \le\frac{24\E Y^2}{C(t_{j-1})(\E Y)^2}
 \longrightarrow 0\quad\text{almost surely}
\]
once $C(t_{j-1})\ge1$.

The partial sums of $\sum_{j\ge1}(I_j-q_j)/j$ form a martingale bounded in $L^2$ since $\sum j^{-2}<\infty$. They converge almost surely. Kronecker's lemma gives $J^{-1}\sum_{j\le J}(I_j-q_j)\to0$ and the Cesàro mean of the $q_j$ also tends to zero. After the first grid point with a core vertex, each successful block increases $\log C$ by at least $\log g$, and each failed block has nonnegative increment. Monotonicity between grid points now gives
\[
 \liminf_{t\to\infty}\frac{\log C(t)}t\ge\frac{\log g}{h+s}
 \quad\text{almost surely}.
\]
If $\mu>0$, choose $M,\epsilon$ with $\eta<\mu$. The mean bound in Lemma~\ref{lem:block} and the existence of the mean rate imply $\liminf_{s\to\infty}s^{-1}\log(\E Y/12)\ge\mu-\eta>0$. Take $s\to\infty$ through the integers and then use a countable choice of parameters with $\eta\downarrow0$. The probability-one bounds thus give the lower limit $\mu$. For $\mu=0$ the lower bound follows from $C(t)\ge1$ eventually.

For each $\delta>0$, Markov's inequality gives $\Prob(C(j)>e^{(\mu+\delta)j})\le e^{-j(\delta+o(1))}$ at integer times. Borel--Cantelli and monotonicity, followed by $\delta\downarrow0$ along a countable sequence, give the upper limit.

Finally the hitting time of size $n$ differs from $\sum_{j=2}^{n-1}1/j$ by an almost surely convergent centred sum. Its ratio to $\log n$ tends to one. Substituting those hitting times proves Theorem~\ref{th:exist}.

\sectag{MOT3}
\section{The three-vertex motif bounds}\label{sec:mot3}
This section proves the numerical bounds $113/200\le\mu\le1133/2000$ of Theorem~\ref{th:exist}, as finite-size bounds on $\E C_n$, and the almost-sure upper bound. It reproduces Sections~1--7 of the note below. The note's Sections~8 and~9, a verification record and the binary certificate format, stay in the repository, together with a \proofsrc{proofs/bracket-motifs-code-map.md}{line-by-line map to the checker}. The weights are checked by the exact integer verifier \proofsrc{verify/motif-bracket/cert_verify_gmp.c}{\texttt{cert\_verify\_gmp.c}}; the builder is not trusted. The note's ``Theorem~1'' is the finite-size bound and its ``Corollary~2'' the almost-sure upper bound. Existence of the exponent is proved in Section~\ref{sec:exist} and is not assumed here.
\begin{center}\footnotesize\begin{tabular}{@{}ll@{}}\toprule
Source & \proofsrc{proofs/bracket-motifs.md}{three-vertex motif note}, lines 5--362\\
Commit & \texttt{a0972085ad72ad5685ae7b424988648d1965e73c}\\
SHA-256 & \texttt{a0126464e9e5d8b6fd13e65d31c92dad20258d3fca257f14224d38f816d038f2}\\
\bottomrule\end{tabular}\end{center}
% BEGIN TRANSCRIPTION proofs/bracket-motifs.md lines 5-362 sha256 a0126464e9e5d8b6fd13e65d31c92dad20258d3fca257f14224d38f816d038f2 commit a0972085ad72ad5685ae7b424988648d1965e73c
\begin{transcription}
\subsection*{1. Model and Main Theorem}

Consider the \textbf{random friend tree} (Krapivsky--Redner isotropic redirection model with $r = 1$, arXiv:1706.01514 and arXiv:2403.20185):
\begin{itemize}
\item Start at $n = 2$ from a single edge $T_2$ on two vertices.
\item Given $T_n$, pick a vertex $V_n \in V(T_n)$ uniformly at random, pick a neighbour $W_n$ of $V_n$ uniformly at random, and attach a new leaf $n+1$ to $W_n$ to form $T_{n+1}$.
\item Let $C_n := \#\{v \in V(T_n) : \deg_{T_n}(v) \ge 2\}$ be the number of non-leaves (\textbf{core} vertices) of $T_n$.
\end{itemize}

Krapivsky and Redner measured a growth exponent near $0.566$. Strogatz later asked whether $17/30$ could be its exact value. The following bounds on $\mathbb{E}[C_n]$ place the upper exponent below that fraction. \href{https://arxiv.org/pdf/1706.01514\#page=12}{Krapivsky and Redner, \S{}5.1} \ensuremath{\cdot} \href{https://x.com/stevenstrogatz/status/1572287647715934210}{Strogatz's post}

\subsubsection*{Theorem 1 (Finite-size bounds for the mean core size)}
For every integer $n \ge 3$,
\fitdisplay{%
\frac{7{,}304{,}352{,}308}{8{,}763{,}936{,}965} \prod_{j=3}^{n-1}\left(1 + \frac{113}{200\, j}\right) \;\le\; \mathbb{E}[C_n] \;\le\; \frac{14{,}606{,}753{,}031}{2{,}989{,}767{,}809} \prod_{j=3}^{n-1}\left(1 + \frac{1133}{2000\, j}\right).
}
Consequently,
\fitdisplay{%
0.5650 \;=\; \frac{113}{200} \;\le\; \liminf_{n\to\infty}\frac{\log \mathbb{E}[C_n]}{\log n} \;\le\; \limsup_{n\to\infty}\frac{\log \mathbb{E}[C_n]}{\log n} \;\le\; \frac{1133}{2000} \;=\; 0.5665 \;<\; \frac{17}{30} = 0.5666\overline{6}.
}
The upper exponent differs from $17/30$ by exactly $17/30-1133/2000=1/6000$.

\subsubsection*{Corollary 2 (Almost-Sure Upper Bound Strictly Below $17/30$)}
Almost surely,
\fitdisplay{%
\limsup_{n\to\infty} \frac{\log C_n}{\log n} \;\le\; \frac{1133}{2000} \;=\; 0.5665 \;<\; \frac{17}{30}.
}
\emph{Proof.} Fix any $\varepsilon > 0$ and put $n_k = 2^k$. By Theorem 1, $\mathbb{E}[C_{n_k}] = O(n_k^{1133/2000})$ since $\prod_{j=3}^{n-1}(1 + \alpha/j) = \frac{2}{\Gamma(3+\alpha)}\frac{\Gamma(n+\alpha)}{\Gamma(n)} \sim \frac{2}{\Gamma(3+\alpha)} n^\alpha$. By Markov's inequality,
\fitdisplay{%
\mathbb{P}\!\left(C_{n_k} > n_k^{1133/2000 + \varepsilon}\right) \le \frac{\mathbb{E}[C_{n_k}]}{n_k^{1133/2000 + \varepsilon}} = O(2^{-k\varepsilon}),
}
which is summable over $k \ge 2$. By Borel--Cantelli, a.s. $C_{n_k} \le n_k^{1133/2000 + \varepsilon}$ for all sufficiently large $k$. Because $C_n$ is nondecreasing in $n$, for $n_k \le n < n_{k+1} = 2 n_k$ we have $C_n \le C_{n_{k+1}} \le (2n)^{1133/2000+\varepsilon}$, so $\limsup_{n\to\infty} \frac{\log C_n}{\log n} \le \frac{1133}{2000} + \varepsilon$ a.s. Intersecting over rational $\varepsilon \downarrow 0$ yields the claim. $\square$

\emph{(Remark: Theorem 1 and Corollary 2 are completely self-contained and do not assume existence of $\lim \frac{\log \mathbb{E}[C_n]}{\log n}$. Combined with the independent existence and almost-sure convergence theorems of \proofsrc{proofs/existence.md}{\texttt{existence.md}} and \proofsrc{proofs/almost-sure-exponent.md}{\texttt{almost-sure-exponent.md}}, which show that a deterministic $\mu$ exists with $\mathbb{E}[C_n] = n^{\mu+o(1)}$ and $\frac{\log C_n}{\log n} \to \mu$ a.s., Theorem 1 yields $\mu \in \left[\frac{113}{200}, \frac{1133}{2000}\right] = [0.5650, 0.5665]$.)}


\subsection*{2. Dynamics in Core Coordinates}

\subsubsection*{2.1 Rooting, core and local states}
At $n = 3$ the tree $T_3$ is the star $K_{1,2}$. Its centre $r$ is the unique vertex of degree $2$. For every $n \ge 3$ we root $T_n$ at $r$.

\textbf{Lemma 2.1 (structure, $n \ge 3$).}
\begin{enumerate}
\item $\deg(r) \ge 2$, so $r \in K(T_n) := \{v : \deg_{T_n}(v) \ge 2\}$.
\item Every non-root core vertex $u$ has a core parent $p_u$. Hence $K(T_n)$ is a connected rooted subtree containing $r$.
\item Every leaf has exactly one neighbour, its parent, and that parent is a core vertex.
\end{enumerate}

\emph{Proof.} Degrees never decrease and $\deg_{T_3}(r) = 2$, which gives (1). If $u \ne r$, its parent $p_u$ is either $r$ or has a parent of its own as well as the child $u$, so $\deg(p_u) \ge 2$. That gives (2). A leaf has degree $1$. It is not $r$, so its only neighbour is its parent, which is core by (2). $\square$

For $u \in K := K(T)$ write $h_u := \mathbf{1}_{u \ne r}$, $L_u :=$ the number of leaf children of $u$, $k_u :=$ the number of core children of $u$, and $D_u := \deg(u) = L_u + k_u + h_u \ge 2$. The triple $s_u := (L_u, k_u, h_u)$ is the \textbf{state} of $u$. Write $\mathrm{ch}_K(u)$ for the core children of $u$, $\iota_u := 1/D_u$, and
\fitdisplay{%
\rho_u := \sum_{v \in N_K(u)} \iota_v = h_u\,\iota_{p_u} + \sum_{c \in \mathrm{ch}_K(u)} \iota_c, \qquad r_u := L_u + \rho_u, \qquad b_u := \frac{L_u}{D_u}.
}
Here $N_K(u)$ is the set of core neighbours of $u$; it has $h_u + k_u$ elements. Since every core degree is at least $2$,
\fitdisplay{%
0 \le \rho_u \le \frac{h_u + k_u}{2}, \qquad L_u \le r_u \le D_u - \frac{h_u + k_u}{2} < D_u + 1 .
}
For a state $s = (L, k, h)$ put $s^+ := (L+1, k, h)$ and, when $L \ge 1$, $s^{\mathrm{bud}} := (L-1, k+1, h)$. Let $\mathrm{nb} := (1, 0, 1)$.

\subsubsection*{2.2 One step of the process}
\textbf{Lemma 2.2 (exact one-step law).} Let $T_n = T$ with $n \ge 3$. Exactly one of the following happens at step $n \to n+1$.
\begin{enumerate}
\item \textbf{Leaf gain at a core vertex $u$}, with probability $r_u/n$. The new leaf becomes a child of $u$, so $s_u \mapsto s_u^+$ and no other state changes.
\item \textbf{Budding at a core vertex $u$ with $L_u \ge 1$}, with probability $b_u/n$. The new leaf is attached to one of the $L_u$ leaf children $\ell$ of $u$, each with probability $1/(nD_u)$. Then $\ell$ becomes a new core vertex in state $\mathrm{nb}$, $s_u \mapsto s_u^{\mathrm{bud}}$ (so $D_u$ is unchanged), and no other state changes.
\end{enumerate}

The probabilities add up to one: $\sum_{u \in K}(r_u + b_u) = n$.

\emph{Proof.} The new leaf is attached to $w$ with probability $\frac{1}{n}\sum_{v \sim w} \iota_v$. For core $w$, its leaf neighbours contribute $1$ each and its core neighbours contribute $\rho_w$, which gives $r_w$. A leaf $\ell$ has the single neighbour $p_\ell$ (Lemma 2.1(3)), which gives $\iota_{p_\ell}$. The state changes are read off directly. The total is $\frac{1}{n}\sum_w \sum_{v \sim w}\iota_v = \frac{1}{n}\sum_v \deg(v)\,\iota_v = 1$. $\square$

\textbf{Lemma 2.3 (monotonicity).} Along the process, $h_u$ is constant, $D_u$ and $k_u$ never decrease, and core vertices stay core. Hence $C_n$ is nondecreasing. For the sets $H_u$ and $C_u$ of \S{}3 (each is defined by $D < \text{const}$, $k \le \text{const}$ and a lower bound on $k$, and $H_u$ also fixes $h = 1$):
\begin{itemize}
\item a vertex that leaves $H_u$ or $C_u$ never re-enters;
\item a vertex enters $H_u$ only by budding from $k = 0$ to $k = 1$, and enters $C_u$ only by budding from $k = 1$ to $k = 2$.
\end{itemize}

\emph{Proof.} Leaf gain raises $D$ by one and keeps $k$. Budding keeps $D$ and raises $k$ by one. Membership in $H_u$ needs $k \ge 1$, and membership in $C_u$ needs $k \ge 2$. The upper bounds can only be broken, never restored, and the lower bound on $k$ can only be crossed by budding, one unit at a time. $\square$

\subsubsection*{2.3 Scaled generator}
For any functional $\Phi$ of rooted trees, Lemma 2.2 gives
\fitdisplay{%
(\mathcal{L}\Phi)(T) := n\,\mathbb{E}\big[\Phi(T_{n+1}) - \Phi(T_n) \,\big|\, T_n = T\big] = \sum_{u \in K(T)} \Big[r_u\,\Delta^{L}_u \Phi(T) + b_u\,\Delta^{B}_u \Phi(T)\Big],
}
where $\Delta^L_u\Phi$ and $\Delta^B_u\Phi$ are the exact changes of $\Phi$ under leaf gain and under budding at $u$. The budding change $\Delta^B_u\Phi$ includes the terms of $\Phi$ that the newborn vertex creates.


\subsection*{3. The Motif Potential $\Phi$}

\subsubsection*{3.1 Parameters}
The potential uses eight integer parameters. The conditions below are the ones the argument uses, each tagged with where it is used.

\begin{center}\footnotesize\begin{tabular}{@{}>{\raggedright\arraybackslash}p{0.230\linewidth}>{\raggedright\arraybackslash}p{0.230\linewidth}>{\raggedright\arraybackslash}p{0.230\linewidth}>{\raggedright\arraybackslash}p{0.230\linewidth}@{}}\toprule
Parameter & Shipped value & Condition & Used in \\\midrule
$L_F$ & 800 & $D_0 \le L_F$ & \S{}5.2, f-cases (C) and (D): $L = L_F$ implies a tail state \\
$K_F$ & 150 & $D_0 \le K_F + 1$ & \S{}5.2, f-case (E), and \S{}6: an overflow state is a tail state \\
$D_0$ & 24 & $3 \le D_0$ & $\mathrm{nb}$ and $T_3$'s root are inner \\
$K_T$ & 12 & $K_T + 1 \le K_F$, $K_T \ge 0$ & clamp $\kappa = \min(k, K_T+1)$ \\
$D_{hu}, K_{hu}$ & 16, 8 & $2 \le D_{hu} \le D_0$, $K_{hu} \ge 1$ & middle vertices of $h_V$ are inner \\
$D_{cu}, K_{cu}$ & 12, 6 & $2 \le D_{cu} \le D_0$, $K_{cu} \ge 2$ & apex vertices of $h_C$ are inner \\
\bottomrule\end{tabular}\end{center}

The verifier tests only some of these conditions on the header. The shipped headers satisfy all of them, for example $D_0 = 24 \le K_F + 1 = 151$. See \proofsrc{proofs/bracket-motifs-code-map.md}{\texttt{bracket-motifs-code-map.md}}, \S{}3.

\subsubsection*{3.2 State sets}
\begin{itemize}
\item A state $(L, k, h)$ is \textbf{inner} if $D = L + k + h < D_0$ and \textbf{tail} otherwise. $\mathcal{I}$ is the set of inner states with $D \ge 2$.
\item \textbf{Tail endpoints:} $\mathrm{INF}$ (standing for $D = \infty$) and $T_1(\kappa, h)$ for $\kappa \in \{0, \dots, K_T+1\}$, $h \in \{0, 1\}$ (standing for $D = D_0$ with clamped child count $\kappa$).
\item $\mathcal{G} := \mathcal{I} \cup \{\mathrm{INF}\} \cup \{T_1(\kappa, h)\}$. The \textbf{parent states} $\mathcal{P} \subset \mathcal{G}$ are $\mathrm{INF}$ and the inner or $T_1$ states with $k \ge 1$ (resp. $\kappa \ge 1$). The \textbf{child states} $\mathcal{W} \subset \mathcal{G}$ are $\mathrm{INF}$ and the inner or $T_1$ states with $h = 1$.
\item \textbf{Middle and apex domains:}
\fitdisplay{%
H_u := \{(L, k, 1) : 1 \le k \le K_{hu},\ L + k + 1 < D_{hu}\}, \qquad C_u := \{(L, k, h) : 2 \le k \le K_{cu},\ L + k + h < D_{cu}\}.
}
  Both are subsets of $\mathcal{I}$, because $D_{hu}, D_{cu} \le D_0$.
\item \textbf{Vertex projection:} $\pi_f(L, k, h) := \mathrm{VSTAR}_h$ if $k > K_F$, and $(\min(L, L_F), k, h)$ otherwise.
\item For $v \in \mathcal{G}$ put $\mathrm{inv}(v) := 1/D_v$ if $v \in \mathcal{I}$, $\mathrm{inv}(\mathrm{INF}) := 0$ and $\mathrm{inv}(T_1(\kappa, h)) := 1/D_0$.
\end{itemize}

\subsubsection*{3.3 Tail interpolation}
For a core vertex $v$ define a probability vector $\mu_v$ on $\mathcal{G}$:
\begin{itemize}
\item if $D_v < D_0$, then $\mu_v := \delta_{s_v}$;
\item if $D_v \ge D_0$, then with $t_v := D_0/D_v \in (0, 1]$ and $\kappa_v := \min(k_v, K_T+1)$,
\fitdisplay{%
\mu_v := (1 - t_v)\,\delta_{\mathrm{INF}} + t_v\,\delta_{T_1(\kappa_v, h_v)} .
}
\end{itemize}

For a tail vertex, $\iota_v = t_v/D_0$. That is exactly the $\mu_v$-average of $\mathrm{inv}$. For an inner vertex, $\iota_v = \mathrm{inv}(s_v)$. So in all cases
\fitdisplay{%
\iota_v = \sum_{a} \mu_v(a)\,\mathrm{inv}(a).
}
We refer to this identity as (3.1).

\subsubsection*{3.4 Tables and the potential}
The certificate supplies rational tables $f$ on $\mathrm{dom}(\pi_f)$, $g$ on $\mathcal{P} \times \mathcal{W}$ with $g[a, b] = 0$ whenever $a, b \notin \mathcal{I}$, $h_V$ on $\mathcal{P} \times H_u \times \mathcal{W}$, and $h_C$ on $C_u \times \mathcal{W} \times \mathcal{W}$, symmetric in its last two arguments. These are evaluated on core vertices by
\fitdisplay{%
g(p \to w) := \sum_{a, b} \mu_p(a)\mu_w(b)\,g[a, b], \qquad h_V(p \to u \to w) := \mathbf{1}_{s_u \in H_u}\sum_{a, b} \mu_p(a)\mu_w(b)\,h_V[a, s_u, b],
}
\fitdisplay{%
h_C(u; \{w_1, w_2\}) := \mathbf{1}_{s_u \in C_u}\sum_{a, b} \mu_{w_1}(a)\mu_{w_2}(b)\,h_C[s_u, a, b],
}
and
\fitdisplay{%
\Phi(T) := \sum_{u \in K} f(\pi_f(s_u)) + \sum_{p \to w} g(p \to w) + \sum_{p \to u \to w} h_V(p \to u \to w) + \sum_{u \in K}\ \sum_{\{w_1, w_2\} \subset \mathrm{ch}_K(u)} h_C(u; \{w_1, w_2\}).
}
The second sum runs over core edges and the third over vertical core paths of length two. Every summand is a \textbf{term}. A term $\tau$ has a \textbf{support} $S(\tau)$ of one, two or three core vertices, and its value depends only on the states of $S(\tau)$. If both ends of an edge are tail vertices, every pair $(a, b)$ in the sum has $a, b \notin \mathcal{I}$, so the term is $0$.

The certificate also supplies centring constants $r_{\mathrm{par}}, r_{\mathrm{ch}}, x_{\mathrm{par}}, x_{\mathrm{ch}}, x_{cc}$ on $\mathcal{I}$, and knot values for three piecewise-linear functions $A_{pc}[s], A_{cp}[s], A_{cc}[s]$ on $[0, 1/2]$ for each $s \in \mathcal{I}$. All three use the knots $0 = \xi_0 < \xi_1 = 1/D_0 < \dots < \xi_{m-1} = 1/2$.


\subsection*{4. The Exact Drift Identity}

This section has no inequalities. It rewrites $(\mathcal{L}\Phi)(T)$ exactly as a sum of one \textbf{local expression} per term.

\subsubsection*{4.1 Event-wise form}
For a term $\tau$ and $x \in S(\tau)$, let $\Delta^L_x\tau$ be the change of $\tau$ when $s_x$ is replaced by $s_x^+$. Let $\Delta^B_x\tau$ be the change when $s_x$ is replaced by $s_x^{\mathrm{bud}}$, with $\Delta^B_x\tau := 0$ if $L_x = 0$. These changes include the indicators $\mathbf{1}_{H_u}$ and $\mathbf{1}_{C_u}$: a path or cherry whose middle or apex leaves its domain jumps to $0$. For a core vertex $x$ with $L_x \ge 1$ let
\fitdisplay{%
N_x := f(\mathrm{nb}) + g(s_x^{\mathrm{bud}} \to \mathrm{nb}) + h_x\,h_V(p_x \to s_x^{\mathrm{bud}} \to \mathrm{nb}) + \sum_{c \in \mathrm{ch}_K(x)} h_C(s_x^{\mathrm{bud}}; \{c, \mathrm{nb}\}),
}
where each motif is evaluated with the current states of $p_x$ and $c$.

\textbf{Identity (I0).}
\fitdisplay{%
(\mathcal{L}\Phi)(T) = \sum_{\tau} \sum_{x \in S(\tau)} \Big[r_x\,\Delta^L_x\tau + b_x\,\Delta^B_x\tau\Big] + \sum_{x \in K} b_x\,N_x .
}

\emph{Proof.} By Lemma 2.2, an event at $x$ changes only the state of $x$, plus the newborn in the case of budding. A term changes only if $x$ is in its support, and then it changes by $\Delta^L_x\tau$ or $\Delta^B_x\tau$. Leaf gain creates no new core vertex, so it creates no new term. Budding at $x$ creates the core vertex $\mathrm{nb}$, which has no core child. The new terms are therefore exactly the terms whose support contains $\mathrm{nb}$:
\begin{itemize}
\item $f(\mathrm{nb})$;
\item the edge $x \to \mathrm{nb}$;
\item the path $p_x \to x \to \mathrm{nb}$ when $h_x = 1$ ($\mathrm{nb}$ can only be the bottom of a path);
\item the cherries $\{c, \mathrm{nb}\}$ at apex $x$, for $c \in \mathrm{ch}_K(x)$.
\end{itemize}

These are the four parts of $N_x$. Could a path or cherry with an \emph{old} support become nonzero because its middle or apex enters $H_u$ or $C_u$? By Lemma 2.3 that happens only when $x$ itself buds from $k = 0$ (no earlier child, so no old path through $x$) or from $k = 1$ (one earlier child, so no old cherry at $x$). Multiply by the probabilities of Lemma 2.2 and by $n$. $\square$

\subsubsection*{4.2 Allocation of neighbour credit}
In (I0) each rate is $r_x = L_x + \sum_{v \in N_K(x)} \iota_v$. Each product $\iota_v\,\Delta^L_x\tau$ is kept in the local expression of $\tau$, or moved to a neighbouring term, according to this table. "Kept" means it stays in $\tau$. A centring constant $c$ splits a product as $\iota_v\Delta = c\Delta + (\iota_v - c)\Delta$. A Fenchel--Young (FY) split with function $A$ writes it as $\iota_v\Delta = A(\iota_v) + R$ with $R := \iota_v\Delta - A(\iota_v)$.

\begin{center}\footnotesize\begin{tabular}{@{}>{\raggedright\arraybackslash}p{0.230\linewidth}>{\raggedright\arraybackslash}p{0.230\linewidth}>{\raggedright\arraybackslash}p{0.230\linewidth}>{\raggedright\arraybackslash}p{0.230\linewidth}@{}}\toprule
Term $\tau$ & Vertex $x$ & Neighbour $v$ of $x$ & Allocation of $\iota_v\,\Delta^L_x\tau$ \\\midrule
$f(u)$ & $u$ tail & any & kept \\
$f(u)$ & $u$ inner & $p_u$ & $r_{\mathrm{par}}[s_u]\Delta$ kept; $(\iota_{p_u} - r_{\mathrm{par}}[s_u])\Delta$ to $E(p_u \to u)$ \\
$f(u)$ & $u$ inner & $c \in \mathrm{ch}_K(u)$ & $r_{\mathrm{ch}}[s_u]\Delta$ kept; $(\iota_c - r_{\mathrm{ch}}[s_u])\Delta$ to $E(u \to c)$ \\
$g(p \to w)$ & $x \in \{p, w\}$ tail & any & kept \\
$g(p \to w)$ & $p$ inner & $w$ & kept \\
$g(p \to w)$ & $p$ inner & $p_p$, $s_p \in H_u$ & $x_{\mathrm{par}}[s_p]\Delta$ kept; $(\iota_{p_p} - x_{\mathrm{par}}[s_p])\Delta$ to $P(p_p \to p \to w)$ \\
$g(p \to w)$ & $p$ inner & $p_p$, $s_p \notin H_u$ & FY with $A_{pc}[s_p]$: $R$ kept; $A_{pc}[s_p](\iota_{p_p})$ to $E(p_p \to p)$ \\
$g(p \to w)$ & $p$ inner & sibling $c \ne w$, $s_p \in C_u$ & $x_{cc}[s_p]\Delta$ kept; $(\iota_c - x_{cc}[s_p])\Delta$ to $Ch(p; \{w, c\})$ \\
$g(p \to w)$ & $p$ inner & sibling $c \ne w$, $s_p \notin C_u$ & FY with $A_{cc}[s_p]$: $R$ kept; $A_{cc}[s_p](\iota_c)$ to $E(p \to c)$ \\
$g(p \to w)$ & $w$ inner & $p$ & kept \\
$g(p \to w)$ & $w$ inner & $c \in \mathrm{ch}_K(w)$, $s_w \in H_u$ & $x_{\mathrm{ch}}[s_w]\Delta$ kept; $(\iota_c - x_{\mathrm{ch}}[s_w])\Delta$ to $P(p \to w \to c)$ \\
$g(p \to w)$ & $w$ inner & $c \in \mathrm{ch}_K(w)$, $s_w \notin H_u$ & FY with $A_{cp}[s_w]$: $R$ kept; $A_{cp}[s_w](\iota_c)$ to $E(w \to c)$ \\
$h_V$, $h_C$ & any & any & kept \\
\bottomrule\end{tabular}\end{center}

Conditions such as "$p_p$ exists" ($h_p = 1$), "a sibling exists" ($k_p \ge 2$) and "a child exists" ($k_w \ge 1$) are implicit. When they fail, the row is empty.

\subsubsection*{4.3 The local expressions}
With this allocation the local expressions are as follows. Write $[\cdot]$ for an indicator, and $\Delta^L_x g$, $\Delta^B_x g$ for the jumps of the term being described.

\textbf{Vertex term, $u \in K$.} Let $\theta_u := h_u r_{\mathrm{par}}[s_u] + k_u r_{\mathrm{ch}}[s_u]$ if $u$ is inner, and $\theta_u := \rho_u$ if $u$ is tail. Then
\fitdisplay{%
V(u) := (L_u + \theta_u)\,\Delta^L_u f + b_u\Big(f(\pi_f(s_u^{\mathrm{bud}})) - f(\pi_f(s_u)) + f(\mathrm{nb}) + g(s_u^{\mathrm{bud}} \to \mathrm{nb})\Big).
}

\textbf{Edge term, $p \to w$.} Let $\sigma_p$ and $\sigma_w$ be the kept rates. If $p$ is inner, $\sigma_p := L_p + \iota_w + [s_p \in H_u]\,x_{\mathrm{par}}[s_p] + [s_p \in C_u](k_p - 1)\,x_{cc}[s_p]$. If $w$ is inner, $\sigma_w := L_w + \iota_p + [s_w \in H_u]\,k_w x_{\mathrm{ch}}[s_w]$. For a tail end, the kept rate is the full rate ($\sigma_p := r_p$, resp. $\sigma_w := r_w$). Then
\fitdisplay{%
\begin{aligned}
E(p \to w) :=\;& [p\ \mathrm{inner}](\iota_w - r_{\mathrm{ch}}[s_p])\Delta^L_p f + [w\ \mathrm{inner}](\iota_p - r_{\mathrm{par}}[s_w])\Delta^L_w f \\
&+ \sigma_p\,\Delta^L_p g + b_p\big(\Delta^B_p g + [s_p^{\mathrm{bud}} \in C_u]\,h_C(s_p^{\mathrm{bud}}; \{w, \mathrm{nb}\})\big) \\
&+ \sigma_w\,\Delta^L_w g + b_w\big(\Delta^B_w g + [s_w^{\mathrm{bud}} \in H_u]\,h_V(p \to s_w^{\mathrm{bud}} \to \mathrm{nb})\big) \\
&+ [p\ \mathrm{inner}, h_p = 1, s_p \notin H_u]\big(\iota_{p_p}\Delta^L_p g - A_{pc}[s_p](\iota_{p_p})\big) \\
&+ [p\ \mathrm{inner}, s_p \notin C_u] \sum_{c \in \mathrm{ch}_K(p) \setminus \{w\}} \big(\iota_c\,\Delta^L_p g - A_{cc}[s_p](\iota_c)\big) \\
&+ [w\ \mathrm{inner}, s_w \notin H_u] \sum_{c \in \mathrm{ch}_K(w)} \big(\iota_c\,\Delta^L_w g - A_{cp}[s_w](\iota_c)\big) \\
&+ [w\ \mathrm{inner}, s_w \notin H_u]\,k_w\,A_{pc}[s_w](\iota_p) + [p\ \mathrm{inner}, h_p = 1, s_p \notin H_u]\,A_{cp}[s_p](\iota_w) \\
&+ [p\ \mathrm{inner}, s_p \notin C_u](k_p - 1)\,A_{cc}[s_p](\iota_w).
\end{aligned}
}
The last two lines are the FY pieces this edge \textbf{receives}. It gets one $A_{pc}[s_w](\iota_p)$ from each path $p \to w \to c$ ($k_w$ of them), one $A_{cp}[s_p](\iota_w)$ from the path $p_p \to p \to w$, and one $A_{cc}[s_p](\iota_w)$ from each sibling $c \ne w$ ($k_p - 1$ of them). The two budding terms in the second and third lines are the newborn cherries and paths of $N_p$ and $N_w$, booked on this edge.

\textbf{Vertical path, $p \to u \to w$ with $s_u \in H_u$.}
\fitdisplay{%
P(p \to u \to w) := (\iota_p - x_{\mathrm{par}}[s_u])\,\Delta^L_u g(u \to w) + (\iota_w - x_{\mathrm{ch}}[s_u])\,\Delta^L_u g(p \to u) + \sum_{x \in \{p, u, w\}} \big(r_x\,\Delta^L_x h_V + b_x\,\Delta^B_x h_V\big).
}

\textbf{Cherry, $u$ with $s_u \in C_u$, $\{w_1, w_2\} \subset \mathrm{ch}_K(u)$.}
\fitdisplay{%
Ch(u; \{w_1, w_2\}) := (\iota_{w_1} - x_{cc}[s_u])\,\Delta^L_u g(u \to w_2) + (\iota_{w_2} - x_{cc}[s_u])\,\Delta^L_u g(u \to w_1) + \sum_{x \in \{u, w_1, w_2\}} \big(r_x\,\Delta^L_x h_C + b_x\,\Delta^B_x h_C\big).
}

\textbf{Identity (I).} For every rooted tree $T$ arising at some $n \ge 3$,
\fitdisplay{%
(\mathcal{L}\Phi)(T) = \sum_{u \in K} V(u) + \sum_{p \to w} E(p \to w) + \sum_{p \to u \to w,\ s_u \in H_u} P(p \to u \to w) + \sum_{u:\ s_u \in C_u}\ \sum_{\{w_1, w_2\}} Ch(u; \{w_1, w_2\}).
}

\emph{Proof.} Start from (I0) and expand each $r_x$ as $L_x$ plus its neighbour credits. Every product $\iota_v\,\Delta^L_x\tau$ in (I0) is listed in exactly one row of the table in \S{}4.2. Every row replaces it by pieces that add back to it: $c\Delta + (\iota_v - c)\Delta = \iota_v\Delta$, and $A(\iota_v) + R = \iota_v\Delta$. The pieces are placed as follows.
\begin{itemize}
\item $V(u)$ keeps $L_u + h_u r_{\mathrm{par}} + k_u r_{\mathrm{ch}}$ for inner $u$. The parent edge and each of the $k_u$ child edges receive one residual.
\item The kept parts of $g(p \to w)$ are $\sigma_p$, $\sigma_w$ and the three $R$ sums.
\item The residuals of the centring splits $x_{\mathrm{par}}, x_{\mathrm{ch}}$ (resp. $x_{cc}$) are the two source terms of $P$ (resp. $Ch$). Each path $p \to u \to w$ with $s_u \in H_u$ receives exactly one residual of each kind, one from edge $u \to w$ (neighbour $p$) and one from edge $p \to u$ (neighbour $w$). Each cherry receives one from each of its two edges.
\item The $A(\iota_v)$ pieces go to the edges listed in the last two lines of $E$, with the multiplicities stated there.
\item The budding parts $b_x\Delta^B_x\tau$ and the four parts of $b_x N_x$ are placed as follows: $f(\mathrm{nb}) + g(s_x^{\mathrm{bud}} \to \mathrm{nb})$ in $V(x)$, the newborn path in $E(p_x \to x)$, and each newborn cherry $\{c, \mathrm{nb}\}$ in $E(x \to c)$.
\end{itemize}

Nothing is counted twice or left out. $\square$

\textbf{Check against the certificate (independent of the verifier).} The \proofsrc{verify/tree-drift/motif\_identity.py}{identity test} implements this section independently of the production checker. It computes $(\mathcal{L}\Phi)(T)$ by brute force over all $n$ attachment targets, computes all local expressions exactly, and compares. Identity (I) held exactly on every tree tested (\S{}8).


\subsection*{5. Bounding the Local Expressions in Both Directions}

Fix a direction. \textbf{UPPER} means $\lambda = 1133/2000$, $\mathop{\mathrm{ext}} = \max$ and $\lessgtr$ is $\le$. \textbf{LOWER} means $\lambda = 113/200$, $\mathop{\mathrm{ext}} = \min$ and $\lessgtr$ is $\ge$. The goal of this section is
\fitdisplay{%
V(u) \lessgtr \lambda f(\pi_f(s_u)), \quad E(p \to w) \lessgtr \lambda g(p \to w), \quad P \lessgtr \lambda h_V, \quad Ch \lessgtr \lambda h_C \qquad \text{for every term of every tree.}
}
We refer to these four inequalities together as (5.0).
By Identity (I), (5.0) implies $(\mathcal{L}\Phi)(T) \lessgtr \lambda\Phi(T)$.

\subsubsection*{5.1 Tools}
\textbf{Lemma 5.1 (piecewise-linear Fenchel--Young).} Let $A$ interpolate $(\xi_j, a_j)_{j<m}$ with $\xi_0 = 0$, $\xi_1 = 1/D_0$, $\xi_{m-1} = 1/2$. Put $A^{*}_{\max}(y) := \max_j(\xi_j y - a_j)$ and $A^{*}_{\min}(y) := \min_j(\xi_j y - a_j)$. Then:
\begin{enumerate}
\item For $x \in [0, 1/2]$ and all $y$: $A^*_{\min}(y) \le xy - A(x) \le A^*_{\max}(y)$.
\item $A^*_{\max}$ is convex and $A^*_{\min}$ is concave.
\item $A(t/D_0) = (1-t)A(0) + tA(1/D_0)$ for $t \in [0, 1]$.
\end{enumerate}

\emph{Proof.} Write $x = (1-\alpha)\xi_i + \alpha\xi_{i+1}$. Then $xy - A(x) = (1-\alpha)(\xi_i y - a_i) + \alpha(\xi_{i+1}y - a_{i+1})$, which proves (1). (2) holds because a max (min) of affine functions is convex (concave). (3) holds because $[0, 1/D_0]$ is one knot interval. $\square$

So for the FY remainders of \S{}4.3, $R = \iota_v\Delta - A(\iota_v)$ satisfies $R \lessgtr A^*_{\mathop{\mathrm{ext}}}(\Delta)$, with $\Delta$ the relevant jump. $A^*_{\mathop{\mathrm{ext}}}$ is convex for UPPER and concave for LOWER.

\textbf{Lemma 5.2 (tail vertices).} Let $v$ be a tail vertex with $\kappa := \kappa_v$, $h := h_v$. For any $F[\cdot]$ on $\mathcal{G}$ write $F(\mu_v) := \sum_a \mu_v(a)F[a]$.
\begin{enumerate}
\item $r_v\,\Delta^L_v F(\mu_v) = t_v\,\gamma_v\,(F[\mathrm{INF}] - F[T_1(\kappa, h)])$ with $\gamma_v := r_v/(D_v+1) \in [\gamma_{\min}(\kappa, h), 1]$. Here $\gamma_{\min}(\kappa, h) := (D_0 - \kappa - h)/(D_0 + 1)$ for $\kappa \le K_T$, and $\gamma_{\min}(K_T + 1, h) := 0$.
\item If $k_v \le K_T$, then $b_v\,\Delta^B_v F(\mu_v) = t_v\,b_v\,(F[T_1(\kappa+1, h)] - F[T_1(\kappa, h)])$ with $b_v \in [b_{\min}(\kappa, h), 1]$, where $b_{\min}(\kappa, h) := (D_0 - \kappa - h)/D_0$. If $k_v > K_T$, then $\Delta^B_v F(\mu_v) = 0$.
\end{enumerate}

\emph{Proof.} (1) Leaf gain keeps $\kappa$ and $h$ and sends $t_v \mapsto D_0/(D_v+1)$, a change of $-t_v/(D_v+1)$. The bound $\gamma_v < 1$ comes from $r_v < D_v + 1$ (\S{}2.1). For the lower bound, $\gamma_v \ge L_v/(D_v+1) = (D_v - k_v - h)/(D_v + 1)$. If $k_v = \kappa \le K_T$, this is nondecreasing in $D_v$ because $k_v + h + 1 \ge 0$, so it is at least its value at $D_v = D_0$. Otherwise $0$ is a valid bound. (2) Budding keeps $D_v$, hence $t_v$, and changes $\kappa$ only if $k_v \le K_T$. Also $b_v = 1 - (k_v + h)/D_v \ge 1 - (\kappa + h)/D_0$. $\square$

\textbf{Lemma 5.3 (quadratic credit).} For $t \in [0, 1]$ and reals $d_0, d_1$:
\fitdisplay{%
\frac{t}{D_0}\min(d_0, d_1) \le \frac{t}{D_0}\big((1-t)d_0 + t d_1\big) \le \frac{t}{D_0}\max(d_0, d_1).
}
With a second parameter $t'$ and bi-affine $d_{ij}$, apply this for each fixed $t'$.

\textbf{Lemma 5.4 (corner reduction).} Let the tail vertices of a support carry parameters $t = (t_1, \dots, t_q)$, $q \le 2$. Let $\omega_c(t) \ge 0$, $\sum_c \omega_c = 1$, be the product weights $\prod_i \mu_{v_i}$ over the corners $c \in \{\mathrm{INF}, T_1\}^q$. Suppose a local expression satisfies $X(t) \lessgtr \sum_c \omega_c(t)\,\widehat X(c)$, while the term is exactly $\tau(t) = \sum_c \omega_c(t)\,\tau[c]$. If $\widehat X(c) \lessgtr \lambda\,\tau[c]$ at every corner, then $X(t) \lessgtr \lambda\tau(t)$. $\square$

The bound $X(t) \lessgtr \sum_c \omega_c(t)\widehat X(c)$ comes from the following facts. Each one holds in the UPPER form; the LOWER form is the same with $\max \leftrightarrow \min$ and $\le \leftrightarrow \ge$.
\begin{itemize}
\item (a) Multi-affine pieces equal their corner interpolation. The source terms $(\iota_p - x)\Delta^L_u g(u \to w)$ are of this kind: by (3.1), $\iota_p$ is affine in $t_p$, and $\Delta^L_u g(u \to w)$ is affine in $t_w$. So are the transfers $(\iota_w - r_{\mathrm{ch}})\Delta^L_p f$ and the received pieces $A(\iota_v)$ for tail $v$, by Lemma 5.1(3).
\item (b) A rate $r$ that is independent of $t$, lies in a known interval, and multiplies a multi-affine jump $Y(t)$ satisfies $rY(t) = \sum_c \omega_c\,rY(c) \le \sum_c \omega_c \max_{r' \in \{r_{\mathrm{lo}}, r_{\mathrm{hi}}\}} r'Y(c)$.
\item (c) $A^*_{\max}(Y(t)) \le \sum_c \omega_c A^*_{\max}(Y(c))$, by convexity (Jensen).
\item (d) Tail rates (Lemma 5.2): $t_v\gamma_v Y(t') \le t_v\sum_{c'}\omega_{c'}(t')\max_{\gamma \in \{\gamma_{\min}, 1\}}\gamma Y(c')$. This puts weight at the $T_1$ corner of $v$ and $0$ at its $\mathrm{INF}$ corner. Budding at tail vertices works the same way with $b$.
\item (e) Products of a tail credit $\iota_v = t_v/D_0$ with a jump that depends on $t_v$ (the only non-multi-affine pieces) are bounded by Lemma 5.3. That again puts weight only at the $T_1$ corner of $v$.
\end{itemize}

\subsubsection*{5.2 All cases, family by family}
Each line below gives a case of the local expression, the bound used, and the resulting finite inequality. Each finite inequality is checked by the verifier (map in \proofsrc{proofs/bracket-motifs-code-map.md}{\texttt{bracket-motifs-code-map.md}}). "Corner" means a table state in $\mathcal{G}$; for tail endpoints, $\mathrm{inv}$ replaces $\iota$ as in (3.1).

\textbf{Vertex family $V(u) \lessgtr \lambda f$.} Cases by $\pi_f(s_u)$:
\begin{itemize}
\item (A) \textbf{Inner} ($D < D_0$, so $L < L_F$). $V(u)$ is a fixed rational:
\fitdisplay{%
(L + h r_{\mathrm{par}} + k r_{\mathrm{ch}})\Delta^L f + B(s) \lessgtr \lambda f(s), \qquad B(s) := [L \ge 1]\tfrac{L}{D}\big(f(s^{\mathrm{bud}}) - f(s) + f(\mathrm{nb}) + g(s^{\mathrm{bud}} \to \mathrm{nb})\big).
}
  When $s^{\mathrm{bud}}$ is inner, $g(s^{\mathrm{bud}} \to \mathrm{nb})$ is a table entry, since $s^{\mathrm{bud}}$ has $D$ unchanged.
\item (B) \textbf{Tail with $L < L_F$ and $k \le K_F$.} Here $\theta_u = \rho_u \in [0, (h+k)/2]$, and every other quantity is a fixed rational. $g(s^{\mathrm{bud}} \to \mathrm{nb})$ is evaluated exactly at $t = D_0/D$, by the verifier's \texttt{eval\_g\_q}. The two inequalities are those with $r \in \{L, L + (h+k)/2\}$ in place of $L + \theta_u$.
\item (C) \textbf{$L = L_F$ exactly, $k \le K_F$.} This is a tail state because $D \ge L_F \ge D_0$. $\Delta^L f = 0$, since $\pi_f$ clamps $L_F + 1 \mapsto L_F$. So $V = \frac{L_F}{D}(f(\pi_f(L_F - 1, k+1, h)) - f + f(\mathrm{nb}) + g(\cdot \to \mathrm{nb}))$, a single exact inequality.
\item (D) \textbf{$L \ge L_F + 1$, $k \le K_F$.} $\Delta^L f = 0$, and $f(\pi_f(s^{\mathrm{bud}})) = f(\pi_f(L_F, k+1, h))$, which is $\mathrm{VSTAR}_h$ if $k = K_F$. The free quantities are $b = L/D \in [\frac{L_F+1}{L_F+1+k+h}, 1]$ and $t = D_0/D \in (0, \frac{D_0}{L_F+1+k+h}]$. The expression $b\cdot(\text{const} + \text{affine}(t))$ is bilinear in $(b, t)$. It is checked at the four corners of the rectangle, which contains every attainable $(b, t)$.
\item (E) \textbf{$k > K_F$ ($\mathrm{VSTAR}_h$).} $\Delta^L f = \Delta^B f = 0$, and $V = b(f(\mathrm{nb}) + g(s^{\mathrm{bud}} \to \mathrm{nb}))$. Since $D \ge K_F + 1 + h \ge D_0$ (this is where the condition $D_0 \le K_F + 1$ is needed), $s^{\mathrm{bud}}$ is tail with $\kappa = K_T + 1$, and $t \le D_0/(K_F + 1 + h)$. The checked rectangle is $b \in [0, 1]$, $t \in [0, D_0/(K_F+1+h)]$.
\end{itemize}

In (A) and (B) the vertex $\mathrm{nb}$ is inner ($D = 2 < D_0$). In every case $s^{\mathrm{bud}}$ has the same degree as $s$.

\textbf{Edge family $E(p \to w) \lessgtr \lambda g$.} Corners are $(a, b) \in \mathcal{P} \times \mathcal{W}$.
\begin{itemize}
\item \textbf{Both ends tail} ($a, b \notin \mathcal{I}$). Then $g \equiv 0$. Events keep both ends tail, so every jump of $g$ is $0$. Tail ends transfer nothing and receive nothing, and the newborn indicators vanish because tail states lie outside $H_u$ and $C_u$. Hence $E \equiv 0 = \lambda g$. The verifier checks $g[a, b] = 0$ for these pairs and skips their rows. This case needs no inequality.
\item \textbf{At least one inner end.} The following summands are evaluated at the corner $(a, b)$; each item names its source in \S{}5.1.
\begin{itemize}
\item Transfers $(\mathrm{inv}(b) - r_{\mathrm{ch}}[a])\Delta^L f(a)$ for inner $a$, and $(\mathrm{inv}(a) - r_{\mathrm{par}}[b])\Delta^L f(b)$ for inner $b$ (fact (a)).
\item Received pieces $A_{\mathrm{tot}}[a, b] := [b \in \mathcal{I}\setminus H_u, k_b \ge 1]\,k_b A_{pc}[b](\mathrm{inv}(a)) + [a \in \mathcal{I}\setminus H_u, h_a = 1]\,A_{cp}[a](\mathrm{inv}(b)) + [a \in \mathcal{I} \setminus C_u, k_a \ge 2](k_a - 1)A_{cc}[a](\mathrm{inv}(b))$ (fact (a) with Lemma 5.1(3)).
\item Kept $p$-rate for inner $a$: $(L_a + [b \in \mathcal{I}]\mathrm{inv}(b) + [a \in H_u]x_{\mathrm{par}}[a] + [a \in C_u](k_a - 1)x_{cc}[a])\,\Delta_p g$, where $\Delta_p g := g[a^+, b] - g[a, b]$. If $D_a = D_0 - 1$, then $a^+ = T_1(\kappa(k_a), h_a)$.
\item Budding at inner $a$: $\frac{L_a}{D_a}(g[a^{\mathrm{bud}}, b] - g[a, b] + [a^{\mathrm{bud}} \in C_u]h_C[a^{\mathrm{bud}}, b, \mathrm{nb}])$ (exact).
\item The same three items for inner $b$ with $\Delta_w g := g[a, b^+] - g[a, b]$, kept rate $L_b + [a \in \mathcal{I}]\mathrm{inv}(a) + [b \in H_u]k_b x_{\mathrm{ch}}[b]$, and newborn path $[b^{\mathrm{bud}} \in H_u]h_V[a, b^{\mathrm{bud}}, \mathrm{nb}]$.
\item FY remainders: $A^*_{pc}[a](\Delta_p g)$ if $a \in \mathcal{I} \setminus H_u$, $h_a = 1$; $(k_a - 1)A^*_{cc}[a](\Delta_p g)$ if $a \in \mathcal{I}\setminus C_u$, $k_a \ge 2$; $k_b A^*_{cp}[b](\Delta_w g)$ if $b \in \mathcal{I}\setminus H_u$, $k_b \ge 1$ (Lemma 5.1(1), fact (c)). The $k_a - 1$ sibling remainders all have the same argument $\Delta^L_p g$, so each is bounded by the same $A^*$.
\item Tail end $a = T_1(\kappa, h)$: $\mathop{\mathrm{ext}}_{\gamma \in \{\gamma_{\min}(\kappa, h), 1\}}\gamma(g[\mathrm{INF}, b] - g[a, b])$; if $\kappa \le K_T$, also $\mathop{\mathrm{ext}}_{\beta \in \{b_{\min}, 1\}}\beta(g[T_1(\kappa+1, h), b] - g[a, b])$; and the quadratic term $\frac{1}{D_0}\mathop{\mathrm{ext}}(g[\mathrm{INF}, b^+] - g[\mathrm{INF}, b],\ g[a, b^+] - g[a, b])$ for the credit $\iota_p\,\Delta^L_w g$ (facts (d), (e)). At $a = \mathrm{INF}$ all these summands are $0$.
\item Tail end $b = T_1(\kappa, 1)$: the mirror image, with quadratic term $\frac{1}{D_0}\mathop{\mathrm{ext}}(g[a^+, \mathrm{INF}] - g[a, \mathrm{INF}],\ g[a^+, b] - g[a, b])$.

  The check is $\mathrm{LHS}_g(a, b) \lessgtr \lambda g[a, b]$ for every pair with at least one inner end.
\end{itemize}
\end{itemize}

\textbf{Path family $P \lessgtr \lambda h_V$.} Corners are $(a, s, b) \in \mathcal{P} \times H_u \times \mathcal{W}$. The middle $s$ is inner with $D_s + 1 \le D_{hu} \le D_0$, so $s^+$ is a table state. If $s^+ \notin H_u$ or $s^{\mathrm{bud}} \notin H_u$, the corresponding value is $0$ (exit, Lemma 2.3).
\begin{itemize}
\item Sources: $(\mathrm{inv}(a) - x_{\mathrm{par}}[s])(g[s^+, b] - g[s, b]) + (\mathrm{inv}(b) - x_{\mathrm{ch}}[s])(g[a, s^+] - g[a, s])$ (fact (a)).
\item Top $a$ inner: rate in $[L_a + \mathrm{inv}(s),\ L_a + \mathrm{inv}(s) + (h_a + k_a - 1)/2]$, since $s$ is one of its $h_a + k_a$ core neighbours (fact (b)); budding exact. Top $a$ tail: Lemma 5.2 together with the quadratic term $\frac{1}{D_0}\mathop{\mathrm{ext}}(h_V[\mathrm{INF}, s^+, b] - h_V[\mathrm{INF}, s, b],\ h_V[a, s^+, b] - h_V[a, s, b])$ for $\iota_p\,\Delta^L_u h_V$.
\item Bottom $b$ inner: rate in $[L_b + \mathrm{inv}(s),\ L_b + \mathrm{inv}(s) + k_b/2]$. Bottom $b$ tail: as for the top, with quadratic term $\frac{1}{D_0}\mathop{\mathrm{ext}}(h_V[a, s^+, \mathrm{INF}] - h_V[a, s, \mathrm{INF}],\ \cdot)$.
\item Middle: rate in $[r_0, r_0 + (k_s - 1)/2]$ with $r_0 = L_s + [a \in \mathcal{I}]\mathrm{inv}(a) + [b \in \mathcal{I}]\mathrm{inv}(b)$, since $s$ has $k_s - 1$ other core neighbours. Budding is exact.
\item When both $a$ and $b$ are tail, the middle's credit $(t_p + t_w)/D_0$ multiplies a jump that is bi-affine in $(t_p, t_w)$. Lemma 5.3 is applied to each summand separately. The verifier adds the two quadratic terms independently at each corner, which is the correct bound.
\end{itemize}

\textbf{Cherry family $Ch \lessgtr \lambda h_C$.} Corners are $(s, b_1, b_2) \in C_u \times \mathcal{W} \times \mathcal{W}$, unordered, including $b_1 = b_2$. Two distinct children in the same table state give the corner $b_1 = b_2$, and the two child ends are still treated as separate vertices.
\begin{itemize}
\item Sources: $(\mathrm{inv}(b_1) - x_{cc}[s])(g[s^+, b_2] - g[s, b_2]) + (\mathrm{inv}(b_2) - x_{cc}[s])(g[s^+, b_1] - g[s, b_1])$.
\item Each child end $w_i$, inner: rate in $[L_i + \mathrm{inv}(s),\ L_i + \mathrm{inv}(s) + k_i/2]$. Tail: Lemma 5.2 and the quadratic term in $\iota_{w_i}\Delta^L_u h_C$.
\item Apex: rate in $[r_0, r_0 + (h_s + k_s - 2)/2]$ with $r_0 = L_s + \sum_i [b_i \in \mathcal{I}]\mathrm{inv}(b_i)$. $s$ may be the root ($h_s = 0$).
\end{itemize}

Applying Lemma 5.4 to each family gives (5.0) in both directions, given that every finite inequality holds. That is what the verifier checks, with 0 failures (\S{}8).


\subsection*{6. Comparison of $\Phi$ with $C$ in Both Directions}

Assign each edge term and each path term to its bottom vertex, and each cherry term to its apex. For a core vertex $u$, let $W(u)$ be $f(\pi_f(s_u))$ plus the terms assigned to $u$. Then $\Phi(T) = \sum_{u \in K} W(u)$. Vertex $u$ receives at most one edge term (none for the root), at most one path term, and $\binom{k_u}{2}[s_u \in C_u] \le \binom{K_{cu}}{2}$ cherry terms.

For $w \in \mathcal{W}$ and $s \in C_u$ put (UPPER with $\mathop{\mathrm{ext}} = \min$, LOWER with $\max$; the $0$ covers an absent term)
\fitdisplay{%
g_{\mathrm{ext}}(w) := \mathop{\mathrm{ext}}\big(0, \mathop{\mathrm{ext}}_{a \in \mathcal{P}} g[a, w]\big), \quad h_{V,\mathrm{ext}}(w) := \mathop{\mathrm{ext}}\big(0, \mathop{\mathrm{ext}}_{a \in \mathcal{P},\, s' \in H_u} h_V[a, s', w]\big), \quad h_{C,\mathrm{ext}}(s) := \mathop{\mathrm{ext}}\big(0, \mathop{\mathrm{ext}}_{b_1, b_2} h_C[s, b_1, b_2]\big).
}
For a core vertex $u$ define its \textbf{candidate endpoints} $E(u) \subset \mathcal{G}$: $E(u) := \{s_u\}$ if $D_u < D_0$, and $E(u) := \{\mathrm{INF}, T_1(\kappa_u, h_u)\}$ if $D_u \ge D_0$. Then $\mathrm{supp}\,\mu_u \subseteq E(u)$. The inclusion can be strict: at $D_u = D_0$ we have $t_u = 1$, so $\mathrm{INF}$ has weight $0$, but it stays in $E(u)$. An evaluated term is a convex combination of table entries whose parent (resp. middle) argument lies in $\mathcal{P}$ (resp. $H_u$) and whose bottom argument lies in $\mathrm{supp}\,\mu_u \subseteq E(u)$. An average over a subset of $E(u)$ is at least the minimum over all of $E(u)$. So in the UPPER direction
\fitdisplay{%
W(u) \ge f(\pi_f(s_u)) + h_u\big(\mathop{\mathrm{ext}}_{b \in E(u)} g_{\mathrm{ext}}(b) + \mathop{\mathrm{ext}}_{b \in E(u)} h_{V,\mathrm{ext}}(b)\big) + [s_u \in C_u]\tbinom{k_u}{2}h_{C,\mathrm{ext}}(s_u) =: \underline W(u),
}
and in the LOWER direction the same holds with $\le$ and $\mathop{\mathrm{ext}} = \max$; call that right side $\overline W(u)$. The verifier takes the extremum over both candidate endpoints for every tail state, including the corner of weight $0$. This gives a bound that is never sharper than the one over $\mathrm{supp}\,\mu_u$, and it is still valid. The right side depends only on $\pi_f(s_u)$. For $L < L_F$ and $k \le K_F$ this is the state itself. For $L \ge L_F$ and $k \le K_F$ the state is tail because $D \ge L_F \ge D_0$, and $E(u) = \{\mathrm{INF}, T_1(\min(k, K_T+1), h)\}$ does not depend on $L$. For $\mathrm{VSTAR}_h$ the state is tail because $D \ge K_F + 1 \ge D_0$, $\kappa = K_T + 1$, and $s_u \notin C_u$. So $\underline W$ and $\overline W$ range over finitely many values. Let $c_0$ be the minimum of $\underline W$ (UPPER) and $C_0$ the maximum of $\overline W$ (LOWER), over all values of $\pi_f(s_u)$. The verifier finds
\fitdisplay{%
c_0 = \tfrac{2{,}989{,}767{,}809}{10^{10}} > 0, \qquad C_0 = \tfrac{17{,}527{,}873{,}930}{10^{10}} > 0,
}
and hence $\Phi_{\mathrm{up}}(T) \ge c_0\,C(T)$ and $\Phi_{\mathrm{lo}}(T) \le C_0\,C(T)$ for every tree with $n \ge 3$.


\subsection*{7. Completion of the Proof of Theorem 1}

$T_3$ has a single core vertex, the root in state $(2, 0, 0)$, and no core edge. So $\Phi(T_3) = f(2, 0, 0)$. This equals $14{,}606{,}753{,}031/10^{10}$ for UPPER and $14{,}608{,}704{,}616/10^{10}$ for LOWER. Both are positive.

\textbf{UPPER.} By \S{}5, $(\mathcal{L}\Phi_{\mathrm{up}})(T) \le \lambda\Phi_{\mathrm{up}}(T)$ for every $T$. Hence $\mathbb{E}\Phi_{\mathrm{up}}(T_{n+1}) = \mathbb{E}\Phi_{\mathrm{up}}(T_n) + \frac{1}{n}\mathbb{E}(\mathcal{L}\Phi_{\mathrm{up}})(T_n) \le (1 + \frac{\lambda}{n})\mathbb{E}\Phi_{\mathrm{up}}(T_n)$. Induction from $n = 3$ and \S{}6 give
\fitdisplay{%
\mathbb{E}[C_n] \le \frac{\mathbb{E}\Phi_{\mathrm{up}}(T_n)}{c_0} \le \frac{\Phi_{\mathrm{up}}(T_3)}{c_0}\prod_{j=3}^{n-1}\Big(1 + \frac{1133}{2000\,j}\Big), \qquad \frac{\Phi_{\mathrm{up}}(T_3)}{c_0} = \frac{14{,}606{,}753{,}031}{2{,}989{,}767{,}809}.
}
\textbf{LOWER.} By \S{}5, $(\mathcal{L}\Phi_{\mathrm{lo}})(T) \ge \lambda\Phi_{\mathrm{lo}}(T)$. Hence $\mathbb{E}\Phi_{\mathrm{lo}}(T_{n+1}) \ge (1 + \frac{\lambda}{n})\mathbb{E}\Phi_{\mathrm{lo}}(T_n)$, and
\fitdisplay{%
\mathbb{E}[C_n] \ge \frac{\mathbb{E}\Phi_{\mathrm{lo}}(T_n)}{C_0} \ge \frac{\Phi_{\mathrm{lo}}(T_3)}{C_0}\prod_{j=3}^{n-1}\Big(1 + \frac{113}{200\,j}\Big), \qquad \frac{\Phi_{\mathrm{lo}}(T_3)}{C_0} = \frac{14{,}608{,}704{,}616}{17{,}527{,}873{,}930} = \frac{7{,}304{,}352{,}308}{8{,}763{,}936{,}965}.
}
Neither induction step needs $\Phi \ge 0$: each step only takes expectations of a pointwise linear inequality. The exponent statements follow from $\prod_{j=3}^{n-1}(1 + \alpha/j) = \frac{2\,\Gamma(n+\alpha)}{\Gamma(3+\alpha)\Gamma(n)} \sim \frac{2}{\Gamma(3+\alpha)}n^{\alpha}$. $\square$
\end{transcription}
% END TRANSCRIPTION proofs/bracket-motifs.md

\sectag{ENC}
\section{The five-vertex enclosure}\label{sec:enc}
\tagnext{ENC}
\begin{theorem}[Five-vertex enclosure]\label{th:gold2}
The exponent of Theorem~\ref{th:exist} satisfies
\[
 a_-:=\frac{3843726496355222669}{6794465353700000000}\ \le\ \mu\ \le\
 \frac{127993535256830043}{226242120050000000}=:a_+ .
\]
The outward decimal enclosure is $[0.565714341933,0.565736986679]$. Consequently $99/175<\mu<142/251$, and if $\mu$ is rational its reduced denominator is at least $426$.
\end{theorem}
The upper endpoint and the weaker lower endpoint
\[
 a_-^{(5)}:=\frac{35083855339718511}{62017003637500000}
\]
come from one pair of five-vertex tables, which the repository calls Gold2; its outward decimal enclosure is $[0.565713486333,0.565736986679]$. The lower endpoint $a_-$ comes from a larger lower table of the same format, combined class by class (Lemma~\ref{lem:classwise}). The enclosure does not exclude $241/426$, the mediant of $99/175$ and $142/251$, and nothing here says on which side of it $\mu$ lies.

\subsection{Proof from the drift tables}\label{sec:enc-proof}
\begin{proof}
In a given core let $c_v$ count the core neighbours of $v$ and $L_v=D_v-c_v$ its leaf neighbours. A core vertex gains a leaf at rate $L_v+\sum_{u\sim v,\,u\ {\rm core}}1/D_u$. It promotes one of its leaves to degree two at rate $L_v/D_v$. These are all the changes in the marked core. For a potential $\Phi$ (a real-valued score on trees), $n$ times its conditional mean increment is its clock generator $\mathcal L\Phi$. For the core count itself, only promotions matter: $\mathcal L C=\sum_v L_v/D_v$, summed over core vertices. Trees with the same core count can have different promotion rates. The certificate adds weights for the local configurations responsible for the difference.

The two potentials sum weights over connected sets of one to five core vertices. A promotion attaches a new leaf to an old leaf; the old leaf becomes the new core vertex and its singleton is charged to the core that selected it. In a newly created supported motif with more than one core vertex, that promoted vertex is a pendant. Removing it leaves the old connected configuration to which the charge is assigned. The nested support conditions ensure that this smaller configuration has an accounting row. For example, the singleton drift at $v$ contains $\sum_{u\sim v,\,u\ {\rm core}}D_u^{-1}\Delta_v^L f$, where $\Delta_v^L f$ is its weight change on gaining a leaf. At an eligible vertex the table substitutes a reference reciprocal $r(v)$ for each donor. The edge row for $uv$ receives $(D_u^{-1}-r(v))\Delta_v^L f$. Summing the singleton and edge rows restores the original drift. Higher transfers use the same cancellation, with integer envelopes for the interpolated terms. This reduces the bound on every finite tree to local inequalities. Weights are in units of $1/S$ with $S=2^{40}$.

For each signed table the checks give $m C_n\le S\Phi\le M C_n$ and
\[
 \sigma S(\mathcal L\Phi-\lambda\Phi)\le e C_n
 \quad\text{on every tree with }n\ge3.
\]
Here $M<\infty$ and $m>0$ depend only on the table. For the Gold2 pair the entries that determine the endpoints are
\[
\begin{array}{c|r|r|r}
 \sigma &\lambda&e&m\\ \hline
 -1&7071421/12500000&19217&99227205820\\
 +1&28286843/50000000&20062&158369484035
\end{array}
\]
Thus $\mathcal L\Phi_-\ge a_-^{(5)}\Phi_-$ and $\mathcal L\Phi_+\le a_+\Phi_+$ with $a_-^{(5)}=\lambda_--e_-/m_-$ and $a_+=\lambda_++e_+/m_+$. In discrete time the bounds multiply at each step by $1+a/n$ for the respective rate $a$. The logarithm of their product from size three to size $n$ is $a\log n+O(1)$ since $\sum_j j^{-2}<\infty$. Comparability with the core size and Theorem~\ref{th:exist} give $a_-^{(5)}\le\mu\le a_+$.

For $a_-$, a larger lower table of the same format has nominal rate $\lambda=56571437/10^8$ and is accepted by the same unmodified finite verifier, tail checker and combiner. The combiner above, which divides the largest charge per core by the smallest mass per core, certifies only about $0.565714206646$ from it. Lemma~\ref{lem:classwise} instead compares charge and mass within classes of core vertices and gives $\mathcal L\Phi\ge(\lambda-\kappa)\Phi$ on every finite tree with
\[
 \kappa=\frac{3814}{135889307074}=\frac{1907}{67944653537},\qquad
 \lambda-\kappa=a_- .
\]
The same iteration and comparability give $\mu\ge a_-$.

Finally $175\cdot142-99\cdot251=1$, so $99/175$ and $142/251$ are Farey neighbours, and exact arithmetic gives $99/175<a_-$ and $a_+<142/251$. If $a/b$ in lowest terms lies strictly between them, then $175a-99b$ and $142b-251a$ are positive integers, and
\[
 b=251(175a-99b)+175(142b-251a)\ge 426 .\qedhere
\]
\end{proof}

\subsection{What the programs check}\label{sec:enc-check}
\paragraph{What the machine checks.}
At cutoff $d$ an endpoint of degree $D_v\ge d$ is represented by two table entries. Weights $d/D_v$ and $1-d/D_v$ on reciprocal degrees $1/d$ and zero give mean $1/D_v$. The changes of these weights under a birth also enter the row bounds. The \proofsrc{round2/proofs/neon2/quintet-cert/verify_quintet.cpp}{finite verifier} checks the transition at the cutoff, every finite drift row and every interpolation envelope. For each Gold2 table there are 767,588,844 rows and 106,953,150 envelope inequalities. At degree at least 4000 the singleton drift is bounded on rectangles in $(L_v/D_v,1/D_v)$. The \proofsrc{round2/proofs/neon2/quintet-cert/verify_tail_bernstein.py}{tail verifier} bounds each polynomial by rational Bernstein coefficients. The \proofsrc{round2/proofs/neon2/quintet-cert/cert_check.py}{combiner} computes the displayed fractions from the signed residuals and comparability constants. The \proofsrc{round2/proofs/neon2/quintet-cert/PROOF.md}{full reduction}, \proofsrc{round2/proofs/neon2/quintet-cert/certificate_gold2.json}{receipt} and \proofsrc{round2/verify/gold/RESULTS.md}{timed rebuild} identify the inputs. The \proofsrc{round2/proofs/neon2/quintet-cert/audit_independent.py}{separate sampler} tests selected rows, projections, the tail and physical trees. The exhaustive scan uses the original verifier. The smaller historical tables use the same proof. Their \proofsrc{round2/proofs/neon2/quintet-cert/certificate_gold1.json}{first five-vertex receipt} gives $[0.565702177913,0.565748407438]$ rounded outwards. The larger lower table, its per-class data and their replay are described in Section~\ref{app:refined}. For comparison with the arithmetic bounds discussed in Section~\ref{sec:further}, the Gold2 pair alone excludes $142/251$ but not $99/175$; the larger lower table excludes $99/175$.

To replay the Gold2 enclosure, rebuild the lower and upper witnesses for the \texttt{gold2} preset, then run the \proofsrc{round2/verify/neon2/run.sh}{strict checker} with \texttt{gold2-check} and the directory containing the pair. Allow 20 GiB of RAM for the check with four workers. The small receipt check in the README recomputes the last step from the stored integers. It checks neither the drift rows nor the reduction from trees to rows.

\subsection{The reduction from trees to rows}\label{sec:enc-reduction}
This subsection proves the two-sided drift bounds and the comparability on every finite tree, given the finite inequalities that the verifier checks. It reproduces Sections~1--4 of the note below, which is all of the note except its header. The note's \texttt{gold1} and \texttt{gold2} are presets of the same table format: \texttt{gold2} gives the upper endpoint of Theorem~\ref{th:gold2}, and Section~\ref{app:refined} uses the same reduction for the larger lower table. Three conventions of the note help in reading it. Its nesting display uses supports such as $G_j^{\mathrm{ex}}$, $M_{4P}$, $M_{4v}$, $M_{5P}$ and $M_{5B}$ that it does not define separately: they are the exact (non-tail) corner states and the internal supports of the shapes listed just before. Their membership is constructed in \texttt{Q5Layout::init} of the verifier's \proofsrc{round2/proofs/neon2/quintet-cert/quintet_format.hpp}{table layout}, including the degree-sum tests, after \texttt{q5\_validate\_header} has checked the inequalities among the cutoffs. Its citation ``\texttt{proofs/existence.md} Theorem~1, Astra'' is the mean-exponent part of Theorem~\ref{th:exist} (Section~\ref{sec:exist}); model names such as Astra identify who wrote a cited proof (Section~\ref{sec:prov}). Its \texttt{UPPER} and \texttt{LOWER} are the two signs of a table.
\begin{center}\footnotesize\begin{tabular}{@{}ll@{}}\toprule
Source & \proofsrc{round2/proofs/neon2/quintet-cert/PROOF.md}{five-vertex reduction note}, lines 14--152\\
Commit & \texttt{a0972085ad72ad5685ae7b424988648d1965e73c}\\
SHA-256 & \texttt{cc7b13244957ded8aa99a9cc7159246369547de656a30fa8994ed3c6f8ab2f45}\\
\bottomrule\end{tabular}\end{center}
% BEGIN TRANSCRIPTION round2/proofs/neon2/quintet-cert/PROOF.md lines 14-152 sha256 cc7b13244957ded8aa99a9cc7159246369547de656a30fa8994ed3c6f8ab2f45 commit a0972085ad72ad5685ae7b424988648d1965e73c
\begin{transcription}
\subsubsection*{1. Unoriented Core Tree and Exact Physical Generator}

For any finite tree $T$ with $|V(T)| \ge 3$, let $K(T) := \{u \in V(T) : \deg_T(u) \ge 2\}$ be its core set, with $C(T) := |K(T)| \ge 1$. Because $T$ is a tree, the induced subgraph $T_K := T[K(T)]$ is a connected tree on $C(T)$ vertices with $C(T) - 1$ edges. Each core vertex $u \in K(T)$ carries its \textbf{unoriented} physical state $(D_u, c_u)$ where:
\begin{itemize}
\item $D_u := \deg_T(u) \ge 2$ is its total degree in $T$,
\item $c_u := \deg_{T_K}(u) \in \{0, 1, \dots, D_u\}$ is its core degree ($c_u = 0 \iff C(T) = 1$, and $c_u \ge 1$ whenever $C(T) \ge 2$),
\item $L_u := D_u - c_u \ge 0$ is its leaf count, and $\iota_u := 1/D_u$.
\end{itemize}

In the continuous-time friend-tree embedding (each vertex independently chooses a uniformly random neighbor at rate $1$ and attaches a fresh leaf to that chosen neighbor), the \textbf{exact} core-tree generator on any function $\Phi$ of the unoriented core tree $(T_K, (D_u)_{u \in K(T)})$ is:
\fitdisplay{%
(\mathcal L \Phi)(T) \;=\; \sum_{u \in K(T)} \Big(r_u \, \Delta_u^L \Phi(T) \;+\; b_u \, \Delta_u^B \Phi(T)\Big),
}
where:
\begin{itemize}
\item $r_u := L_u + \sum_{v \sim_{T_K} u} \iota_v$ is the leaf-gain rate at $u$ (each of $u$'s $L_u$ leaves chooses $u$ at rate $1$, and each core neighbor $v \sim_{T_K} u$ chooses $u$ at rate $1/D_v = \iota_v$), changing $(D_u, c_u) \mapsto (D_u + 1, c_u)$ and leaving $T_K$ unchanged;
\item $b_u := L_u / D_u$ is the leaf-promotion rate at $u$ ($u$ chooses one of its $L_u$ leaves at rate $L_u / D_u$, attaching a fresh leaf to that leaf so it becomes a new core leaf $J$ with state $(2, 1)$ adjacent to $u$, while $u$ changes from $(D_u, c_u)$ to $(D_u, c_u + 1)$; zero if $L_u = 0$).
\end{itemize}

Notice that this identity holds \textbf{globally at every core vertex of every finite tree $T$ with $C(T) \ge 1$} (including $C(T) = 1$, where $r_u = D_u$ and $b_u = 1$, and $C(T) = 2$), with no root/founder exception and no dummy vertex.


\subsubsection*{2. Eight Unoriented Motif Shapes and $c$-Capped Endpoint Interpolation}

\trhead{2.1. Endpoint spaces with $c$-capped inner overflow}
At each endpoint level $j \in \{2, 3, 4, 5\}$ with integer cutoffs $(d_j, K_j)$ ($2 \le K_j < d_j$), define the finite corner set $G_j$ consisting of four disjoint kinds:
\begin{enumerate}
\item \textbf{Exact inner} (\texttt{kind = 0}): pairs $(D, c)$ with $2 \le D < d_j$ and $1 \le c \le \min(D, K_j)$, reciprocal $\mathrm{inv} = 1/D$.
\item \textbf{Exact-degree inner overflow} (\texttt{kind = 3}): one state $\mathrm{OVF}_j(D)$ for each $K_j + 1 \le D < d_j$ representing all states $(D, c)$ with $c \ge K_j + 1$, exact reciprocal $\mathrm{inv} = 1/D$.
\item \textbf{One-knot degree tail} (\texttt{kind = 1}): classes $T_j(\kappa)$ for $\kappa = 1, \dots, K_j + 1$, representing $D \ge d_j$ with $\min(c, K_j + 1) = \kappa$, reciprocal $\mathrm{inv} = 1/d_j$.
\item \textbf{Infinity corner} (\texttt{kind = 2}): $\mathrm{INF}_j$, reciprocal $\mathrm{inv} = 0$.
\end{enumerate}

For any physical core vertex $u$ with $c_u \ge 1$, its deterministic level-$j$ interpolation distribution $\mu_u^{(j)}$ on $G_j$ is:
\fitdisplay{%
\mu_u^{(j)} := \begin{cases}
\delta_{(D_u, c_u)} & \text{if } D_u < d_j \text{ and } c_u \le K_j, \\ \delta_{\mathrm{OVF}_j(D_u)} & \text{if } D_u < d_j \text{ and } c_u > K_j, \\ \frac{d_j}{D_u}\,\delta_{T_j(\min(c_u, K_j+1))} + \big(1 - \frac{d_j}{D_u}\big)\,\delta_{\mathrm{INF}_j} & \text{if } D_u \ge d_j.
\end{cases}
}
Observe that $\mathbb E_{\mu_u^{(j)}}[\mathrm{inv}] = 1/D_u = \iota_u$ \textbf{identically} for every $(D_u, c_u)$.

Whenever $d_{j+1} \le d_j$ and $K_{j+1} \le K_j$, define the exact stochastic projection $\Pi_{j \to j+1} : G_j \to \mathcal P(G_{j+1})$ by applying the exact same rule to the corner label $(D, \kappa)$ of $G_j$ (mapping $\mathrm{INF}_j \mapsto \delta_{\mathrm{INF}_{j+1}}$). Because degree scaling composes ($\frac{d_j}{D}\cdot\frac{d_{j+1}}{d_j} = \frac{d_{j+1}}{D}$) and clamping composes ($\min(\min(c, K_j+1), K_{j+1}+1) = \min(c, K_{j+1}+1)$), we have the exact identity $\mu_u^{(j)} \Pi_{j \to j+1} = \mu_u^{(j+1)}$ for every $(D_u, c_u)$.

\trhead{2.2. Eight unoriented motifs and nested downward-closed supports}
A quantized potential $\Phi(T)$ in integer ticks $S = 2^{40}$ is the sum over all unoriented embeddings into $T_K$ of \textbf{only eight shapes}, covering every connected tree up to 5 vertices (diameter 4):
\begin{enumerate}
\item \textbf{Singleton} $f(D_u, c_u)$: exact table for $2 \le D_u < D_F$, analytic polynomial in $(q, x) = (L_u/D_u, 1/D_u)$ for $D_u \ge D_F$.
\item \textbf{Edge} $g(\{u, v\})$: symmetric bilinear expectation over $\mu_u^{(2)} \otimes \mu_v^{(2)}$, with $g(\mathrm{INF}_2, \mathrm{INF}_2) = 0$.
\item \textbf{3-path} $H_3(m; \{a, b\})$: internal middle $m \in M_3$ ($D_m < d_{3m}, 2 \le c_m \le K_{3m}$); endpoints $\{a, b\}$ evaluated over $\mu_a^{(3)} \otimes \mu_b^{(3)}$.
\item \textbf{4-path} $P_4(a - b - c - d)$: palindrome-symmetric, internal edge $(b, c)$ with $D_b, D_c < d_{4i}$, $2 \le c_b, c_c \le K_{4i}$, $D_b + D_c \le S_{P4}$; endpoints $(a, d)$ over $\mu_a^{(4)} \otimes \mu_d^{(4)}$.
\item \textbf{3-star} $S_3(o; \{a, b, c\})$: hub $o \in M_{S3}$ ($D_o < d_{S3}, 3 \le c_o \le K_{S3}$); endpoints $\{a, b, c\}$ over $(\mu^{(4)})^{\otimes 3}$.
\item \textbf{5-path} $P_5(a - b - c - d - e)$: palindrome-symmetric, internal 3-path $(b, c, d)$ with $D_b, D_c, D_d < d_{5i}$, $2 \le c_b, c_c, c_d \le K_{5i}$, $D_b + D_c \le S_{P4}$, $D_c + D_d \le S_{P4}$, $D_b + D_c + D_d \le S_{P5}$; endpoints $(a, e)$ over $\mu_a^{(5)} \otimes \mu_e^{(5)}$.
\item \textbf{5-broom} $B_5(o; \{a, b\}; c - d)$: hub $o$ ($D_o < d_{5o}, 3 \le c_o \le K_{5o}$) adjacent to $\{a, b\}$ and arm $c$ ($D_c < d_{5c}, 2 \le c_c \le K_{5c}$) with $D_o + D_c \le S_{B5}$; endpoints $\{a, b\}$ and $d$ over $(\mu^{(5)})^{\otimes 3}$.
\item \textbf{4-star} $S_4(o; \{a, b, c, d\})$: hub $o \in M_{S4}$ ($D_o < d_{S4}, 4 \le c_o \le K_{S4}$); endpoints $\{a, b, c, d\}$ over $(\mu^{(5)})^{\otimes 4}$.
\end{enumerate}

The header validation (\texttt{q5\_validate\_header}) enforces strict nesting of all internal and endpoint cutoffs:
\fitdisplay{\begin{multlined}[\linewidth]%
M_3 \subseteq G_2^{\mathrm{ex}}, \quad M_{4P} \subseteq (M_3 \cap G_3^{\mathrm{ex}})^2, \quad M_{S3} \subseteq M_3, \ftbreak \quad M_{5P} \subseteq (M_{4v} \cap G_4^{\mathrm{ex}})^3, \quad M_{5B} \subseteq (M_{4v} \cap M_{S3}) \times (M_{4v} \cap G_4^{\mathrm{ex}}), \quad M_{S4} \subseteq M_{S3}.
\end{multlined}}
\textbf{Downward closure \& exact newborn bijection:} Every core event increases either $D_v$ (gain) or $c_v$ (promotion), so an existing motif embedding can never enter a higher support (it can only stay in support or drop to $0$). Conversely, when a vertex $u$ promotes and spawns $J = (2, 1)$, deleting $J$ from any newly created supported $k$-vertex motif ($2 \le k \le 5$) yields a \textbf{unique} pre-existing supported $(k-1)$-vertex motif (because of the strict nesting above), and its newborn weight is allocated cleanly to that parent motif's promotion term.


\subsubsection*{3. Exact Transfers, Envelopes, and Local Corner Inequalities}

Let $\sigma = +1$ (UPPER) or $\sigma = -1$ (LOWER). Every local inequality bounds $\sigma S (\text{row drift} - \lambda \cdot \text{row weight})$ from above in integer ticks.

1. \textbf{Paired inter-level transfers (exact cancellation):}
\begin{itemize}
\item $f \to g$: At any exact $G_2$ vertex $u$ ($D_u < d_2, 1 \le c_u \le K_2$), replace $u$'s neighbor gain rate $\sum_{v \sim u} \iota_v$ in the $f$-row by $c_u r_1(u)$, transferring $(\iota_v - r_1(u))\Delta_u^L f$ to the incident edge $g(u, v)$.
\end{itemize}

\ensuremath{\circ} $g \to H_3$ \& Fenchel--Young (FY): At an exact $G_2$ vertex $u$ with $c_u \ge 2$:
\begin{itemize}
\item If $u \in M_3$: replace the external donor rate $\sum_{w \sim u, w \ne v} \iota_w$ in $g(u, v)$ by $(c_u - 1) r_1(u)$, transferring $(\iota_w - r_1(u))\Delta_u^L g(u, v)$ to $H_3(u; \{v, w\})$.
\end{itemize}

\ensuremath{\blacksquare} If $u \notin M_3$: transfer a piecewise-linear function $A_u(\iota_w)$ (on rational knot grid $0 = \xi_0 < \xi_1 = 1/d_2 < \dots < \xi_{m-1} = 1/2$, denominator $Q_K$) from edge $g(u, v)$ to edge $g(u, w)$ for every $w \ne v$. Summing over all ordered pairs of incident edges at $u$ cancels identically. Thus $g(u, v)$ receives $(c_u - 1) A_u(\iota_v)$ (since $A_u$ is affine on $[\xi_0, \xi_1] = [0, 1/d_2]$, its expectation over $\mu_v^{(2)}$ is \textbf{exact} even when $D_v \ge d_2$) and keeps, in signed units,

\fitdisplay{%
\sum_{w \ne v} \sigma\big(\iota_w \Delta_u^L g(u, v) - A_u(\iota_w)\big) \;\le\; (c_u - 1) \max_{0 \le k < m} \sigma\big(\xi_k \Delta_u^L g(u, v) - A_u(\xi_k)\big),
}

because each segment $[\xi_k, \xi_{k+1}]$ is affine in $\iota_w$, and $z \mapsto \max_k \sigma(\xi_k z - A_u(\xi_k))$ is convex in $z = \Delta_u^L g(u, v) = \mathbb E_{V \sim \mu_v^{(2)}}[\Delta_u^L g(u, V)]$ so Jensen's inequality passes directly to the corners $V \in G_2$.
\begin{itemize}
\item $H_3 \to P_4, S_3$ and $P_4, S_3 \to P_5, B_5, S_4$: Whenever adjoining an external core neighbor $w$ of an exact-inner motif vertex $x$ produces a supported motif of the next size, retain $r_1(x)$ per external donor in the lower-level row and transfer $(\iota_w - r_1(x))\Delta_x^L h$ to the unique higher-level motif. Because every higher-level motif has full $G_{j+1}$ tail endpoints, the support condition depends \textbf{only} on the internal vertices, never on the degree of $w$ or the other endpoints. Uncovered external donors at exact-inner vertices have $\iota_w \in [0, 1/2]$, bounded by extremizing over $[0, \frac{1}{2} n_{\mathrm{ext}}]$.
All transfers occur in exact cancelling pairs, so summing over all embeddings on any tree $T$ preserves $(\mathcal L \Phi)(T)$ identically.
\end{itemize}

2. \textbf{Exact 1D, 2D, and 3D projection envelopes:}
Each cross-level source change $d_{\mathrm{fine}}$ on $G_j^k$ ($k \in \{1, 2, 3\}$) is enclosed by stored integer arrays $(m, r)$ on $G_{j+1}^k$ with $r \ge 0$ verified over \textbf{every} fine state/pair/triple in exact \texttt{\_\_int128} arithmetic (clearing integer denominators):

\fitdisplay{%
\big| d_{\mathrm{fine}}(s) - (\Pi^{\otimes k} m)(s) \big| \;\le\; (\Pi^{\otimes k} r)(s).
}

By independence of the donor coordinate from the remaining coordinates and convexity of $|\cdot|$,

\fitdisplay{%
\sigma(\iota_w - r_1)\,\mathbb E[d_{\mathrm{fine}}] \;\le\; \mathbb E\big[\sigma(\mathrm{inv}_{j+1}(w) - r_1)\,m + |\mathrm{inv}_{j+1}(w) - r_1|\,r\big].
}

\begin{enumerate}\setcounter{enumi}{2}
\item \textbf{Exact-degree overflow inner endpoints (\texttt{kind = 3}), one-knot tail endpoints (\texttt{kind = 1}), and double-tail / \texttt{INF} edges:}
\begin{itemize}
\item For \texttt{kind = 3} ($K_j < D_u < d_j$, $c_u \ge K_j + 1$): $\mathrm{inv} = 1/D_u$ is exact; promotion leaves the state unchanged ($\Delta_u^B = 0$); its adjacent motif neighbor contributes exact rate $\iota_{\mathrm{nbr}}$ (when inner) or tail quadratic credit (charged once at the tail when the adjacent vertex is a tail); and its remaining gain rate $L_u + \sum_{z \ne \mathrm{nbr}} \iota_z \le (D_u - c_u) + \frac{c_u - 1}{2} = D_u - \frac{c_u + 1}{2} \le D_u - \frac{K_j + 2}{2}$ lies in $[0, (2 D_u - K_j - 2)/2]$.
\item For \texttt{kind = 1} ($D_u \ge d_j$, $t = d_j / D_u$): writing $H_1, H_0$ for $T_j(\kappa)$ and $\mathrm{INF}_j$, leaf gain contributes $t \gamma (H_0 - H_1)$ with $\gamma = (L_u + \sum \iota_z)/(D_u + 1) \in [\gamma_{\min}, 1]$, where $\gamma_{\min} = \max(0, (d_j - \kappa + \iota_{\mathrm{nbr,lo}})/(d_j + 1))$ for exact $\kappa \le K_j$ (and $0$ for overflow $\kappa = K_j + 1$); promotion contributes $t b (H_1^{\mathrm{bud}} - H_1)$ with $b \in [\max(0, d_j - \kappa)/d_j, 1]$ for $\kappa \le K_j$ (and zero weight change for overflow $\kappa = K_j + 1$, since promotion leaves the clamped class unchanged, contributing $0$ even when $b > 0$); and whenever the adjacent vertex $y$ is inner (\texttt{kind = 0} or \texttt{kind = 3}), $y$'s leaf gain from $u$ satisfies $\sigma \frac{t}{d_j}((1-t)d_0 + t d_1) \le \frac{t}{d_j}\max(\sigma d_0, \sigma d_1)$ (absorbing both the $t$ and $1-t$ parts of $\Delta_y^L$ into the $t$-row). When several distinct tail endpoints attach to the same internal hub, the same linear inequality applies independently to each coordinate after taking expectation over the others.
\item \textbf{Double-tail and \texttt{INF} corners in $g(u, v)$:} Only $g(u, v)$ has two adjacent endpoints that can simultaneously be physical tails ($D_u, D_v \ge d_2$, with $t_u = d_2/D_u, t_v = d_2/D_v$). Expanding $g(u, v)$ over $\mu_u^{(2)} \otimes \mu_v^{(2)}$ and using $g(\mathrm{INF}_2, \mathrm{INF}_2) = 0$: the $(\mathrm{INF}_2, \mathrm{INF}_2)$ corner carries zero value, zero incoming transfers, and zero residual; on $(\mathrm{INF}_2, y)$ with $y$ inner, $y$ still evolves under its own leaves and external neighbors while $u$'s credit $\iota_u$ is already accounted for in $(T_u, y)$. When both $u$ and $v$ have $D_u, D_v \ge d_2$, the cross-donor rates $\iota_v$ and $\iota_u$ are \textbf{not} dropped: they are fully contained inside $r_u$ and $r_v$ (hence inside $\gamma_u, \gamma_v \in [(d_2 - \kappa)/(d_2 + 1), 1]$ for exact classes $\kappa \le K_2$ --- taking $\iota_{\mathrm{nbr,lo}} = 0$ since $\iota > 0$ --- and inside $[0, 1]$ for overflow class $\kappa = K_2 + 1$), and averaging $t_u \gamma_u (g(\mathrm{INF}_2, \cdot) - g(T_u, \cdot))$ over $\mu_v^{(2)}$ assigns exact weights $t_u t_v$ to row $(T_u, T_v)$ and $t_u(1 - t_v)$ to row $(T_u, \mathrm{INF}_2)$ (and symmetrically for $v$).

\end{itemize}
\item \textbf{Analytic singleton tail ($D \ge D_F$) and seam ($D = D_F - 1 \to D_F$):}
   Writing $q = L/D \in [0, 1]$ and $x = 1/D \in [0, 1/D_F]$, the unsigned analytic potential in ticks is $w(q, x) = f_{p0} + f_{p1} q$ (UPPER) or $w(q, x) = f_{p0} q - f_{p1} q^2 - f_{p2} q x$ (LOWER). At the seam $D = D_F - 1$, \texttt{verify\_quintet.cpp} evaluates the exact rational jump $w((D_F - c)/D_F, 1/D_F) - f(D_F - 1, c)$ directly in exact 128-bit arithmetic. For $D \ge D_F$, no supported higher motif can contain the promoted vertex as the internal neighbor of $J$, so promotion adds only $f(2,1) + g((D, c+1), J)$ with post-promotion class $\min(c+1, K_2+1) \in \{1, \dots, K_2+1\}$. With $C_{\mathrm{nb}} = f(2,1) + g(\mathrm{INF}_2, J)$ and $d_g = d_2(\mathrm{ext}_{\kappa \in \{1,\dots,K_2+1\}} g(T_2(\kappa), J) - g(\mathrm{INF}_2, J))$ ($\max$ for UPPER, $\min$ for LOWER), multiplying the signed residual at either normalized gain-rate endpoint $\rho \in \{q, (1+q)/2\}$ by $(1+x)^2$ yields an exact polynomial $P_\rho(q, x)$ enclosed on $[0, 1] \times [0, 1/D_F]$ by exact rational 2D tensor Bernstein coefficients across 16 equal $q$-panels (\texttt{verify\_tail\_bernstein.py}). Dividing a panel's Bernstein upper bound $M_j$ by $(1+x)^2 \in [1, (1 + 1/D_F)^2]$ gives the rigorous upper bound $M_j$ if $M_j \ge 0$ and $M_j / (1 + 1/D_F)^2$ if $M_j < 0$.
\end{enumerate}


\subsubsection*{4. Per-Core Aggregation, Global Comparability, and Exponent Deduction}

Let $e_f := \max(\mathrm{cap}_{f,\mathrm{in}}, \mathrm{cap}_{f,\mathrm{tab}}, \mathrm{cap}_{f,\mathrm{ana}})$, $e_g := \max(0, \mathrm{cap}_g)$, and let $e_{\mathrm{high}} \ge 0$ be the maximum over all finite states $(D, c)$ ($D < d_{3m}$) of the positive higher-motif caps allocated to a host core vertex ($H_3 \to m$; $\frac{1}{2} P_4 \to b, c$; $S_3 \to o$; $P_5 \to c_{\mathrm{mid}}$; $B_5 \to o$; $S_4 \to o$). Since $T_K$ has $C(T)$ vertices and $C(T) - 1$ edges, setting $E_{\mathrm{tot}} := \max(0, e_f + e_g + e_{\mathrm{high}})$ gives
\fitdisplay{%
\sigma S \big((\mathcal L \Phi)(T) - \lambda \Phi(T)\big) \;\le\; E_{\mathrm{tot}}\, C(T) \qquad \text{for every finite tree } T \text{ with } C(T) \ge 1.
}

For comparability, choose any spanning rooting of $T_K$ (purely as an accounting bijection, leaving $\Phi$ unchanged): allocate $f(u)$ to $u$, $g(p_u, u)$ to child $u \ne r$ ($0$ to $r$), and each higher motif to its finite-degree host vertex as above. On any core tree $T_K$, $\sum_{u \in K(T)} (c_u - 2) = -2$, so for any $\alpha \ge 0$ (\texttt{alpha = 0} for UPPER; any validated $\alpha \in [0, 1]$ for LOWER, with $\alpha = 1/4$ in \texttt{gold1}):
\fitdisplay{%
S\Phi(T) \;=\; 2\alpha S + \sum_{u \in K(T)} \big(W(u) + \alpha S (c_u - 2)\big) \;\ge\; m_{\min}\, C(T), \qquad S\Phi(T) \;\le\; m_{\max}\, C(T),
}
where $m_{\min} > 0$ and $m_{\max} > 0$ are verified over all $D < D_F$ (\texttt{verify\_quintet.cpp}) and analytically over $D \ge D_F$ (\texttt{verify\_tail\_bernstein.py}, where no higher motif is hosted; any incident higher motif having a $D \ge D_F$ vertex as an outer endpoint is hosted by its finite-degree internal vertex and already accounted for in that host's finite comparability check).

Consequently, for every finite tree $T$ with $C(T) \ge 1$, setting $\lambda_{\mathrm{up}}^* := \lambda_{\mathrm{up}} + E_{\mathrm{tot}}/m_{\min}$ and $\lambda_{\mathrm{lo}}^* := \lambda_{\mathrm{lo}} - E_{\mathrm{tot}}/m_{\min}$:
\fitdisplay{%
(\mathcal L \Phi_{\mathrm{up}})(T) \;\le\; \lambda_{\mathrm{up}}^*\,\Phi_{\mathrm{up}}(T), \qquad (\mathcal L \Phi_{\mathrm{lo}})(T) \;\ge\; \lambda_{\mathrm{lo}}^*\,\Phi_{\mathrm{lo}}(T).
}

\textbf{Non-explosion, stopped Dynkin, and Gr\"onwall deduction:} Starting from the 3-vertex tree $T_3$ ($C(T_3) = 1$), the total vertex count $N(t) = |V(T_t)|$ increases by $1$ at total rate $N(t)$, so $N(t)$ is a non-explosive rate-$1$ Yule process with $\mathbb E[N(t)^2] < \infty$ for every $t \ge 0$, and $C(T_t) \le N(t)$. Let $\tau_M := \inf\{t \ge 0 : N(t) \ge M\}$. For either potential $\Phi \in \{\Phi_{\mathrm{lo}}, \Phi_{\mathrm{up}}\}$, comparability gives $0 < (m_{\min}/S) C(T) \le \Phi(T) \le (m_{\max}/S) N(T)$, and every single jump increases $N$ by $1$ (total jump rate $N(T)$), so
\fitdisplay{%
|(\mathcal L \Phi)(T)| \;\le\; N(T)\big(\tfrac{m_{\max}}{S}(N(T) + 1) + \tfrac{m_{\max}}{S} N(T)\big) \;=\; \tfrac{m_{\max}}{S} N(T)(2N(T) + 1).
}
On $[0, t \wedge \tau_M]$, the chain takes values in the finite set $\{N \le M\}$, so Dynkin's formula gives
\fitdisplay{%
\mathbb E[\Phi(T_{t \wedge \tau_M})] \;=\; \Phi(T_3) + \mathbb E\left[\int_0^{t \wedge \tau_M} (\mathcal L \Phi)(T_s)\,ds\right].
}
Since $t \wedge \tau_M \uparrow t$ a.s., $0 \le \Phi(T_{t \wedge \tau_M}) \le (m_{\max}/S) N(t) \in L^1$, and $\int_0^{t \wedge \tau_M} |(\mathcal L\Phi)(T_s)|\,ds \le (m_{\max}/S)\,t\,N(t)(2N(t)+1) \in L^1$, dominated convergence as $M \to \infty$ yields $\mathbb E[\Phi(T_t)] = \Phi(T_3) + \int_0^t \mathbb E[(\mathcal L\Phi)(T_s)]\,ds$ for all $t \ge 0$. Hence $t \mapsto \mathbb E[\Phi(T_t)]$ is absolutely continuous and satisfies the differential Gr\"onwall inequalities, giving for all $t \ge 0$:
\fitdisplay{%
\Phi_{\mathrm{lo}}(T_3)\,e^{\lambda_{\mathrm{lo}}^* t} \;\le\; \mathbb E[\Phi_{\mathrm{lo}}(T_t)] \;\le\; \frac{m_{\max,\mathrm{lo}}}{S}\,\mathbb E[C(T_t)], \qquad \frac{m_{\min,\mathrm{up}}}{S}\,\mathbb E[C(T_t)] \;\le\; \mathbb E[\Phi_{\mathrm{up}}(T_t)] \;\le\; \Phi_{\mathrm{up}}(T_3)\,e^{\lambda_{\mathrm{up}}^* t}.
}
Taking $t^{-1}\log(\cdot)$ unconditionally proves, from the 3-vertex seed $T_3$,
\fitdisplay{%
\lambda_{\mathrm{lo}}^* \;\le\; \liminf_{t\to\infty} \frac{1}{t}\log \mathbb E_{T_3}[C(T_t)] \;\le\; \limsup_{t\to\infty} \frac{1}{t}\log \mathbb E_{T_3}[C(T_t)] \;\le\; \lambda_{\mathrm{up}}^*.
}
Passing to the original single-edge seed $T_2$ follows directly by conditioning on the first jump time $\tau_1 \sim \mathrm{Exp}(2)$ (which transitions deterministically from $T_2$ to $T_3$, with $C(T_s) = 0$ for $s < \tau_1$):
\fitdisplay{%
\mathbb E_{T_2}[C(T_t)] \;=\; \int_0^t 2 e^{-2s}\,\mathbb E_{T_3}[C(T_{t - s})]\,ds,
}
so the upper bound multiplies by $2/(2 + \lambda_{\mathrm{up}}^*)$ and restricting the lower integral to $s \in [0, 1]$ (for $t \ge 1$) gives a positive constant multiple of $e^{\lambda_{\mathrm{lo}}^* t}$, preserving the identical $\liminf/\limsup$ bounds from $T_2$. (In discrete size $n \ge 3$, $(\mathcal L\Phi)(T_n) = n\,\mathbb E[\Phi(T_{n+1}) - \Phi(T_n) \mid T_n]$ directly gives $\Phi_{\mathrm{lo}}(T_3)\prod_{j=3}^{n-1}(1 + \lambda_{\mathrm{lo}}^*/j) \le \mathbb E[\Phi_{\mathrm{lo}}(T_n)]$ when $\lambda_{\mathrm{lo}}^* > 0$ --- and when $\lambda_{\mathrm{lo}}^* \le 0$, $C(T_n) \ge 1$ trivially yields $\liminf_{n\to\infty}(\log n)^{-1}\log\mathbb E[C(T_n)] \ge 0 \ge \lambda_{\mathrm{lo}}^*$ --- together with $\mathbb E[\Phi_{\mathrm{up}}(T_n)] \le \Phi_{\mathrm{up}}(T_3)\prod_{j=3}^{n-1}(1 + \lambda_{\mathrm{up}}^*/j)$, yielding the identical unconditional $\liminf/\limsup$ bounds for $(\log n)^{-1}\log\mathbb E[C(T_n)]$ as well.) Combining these unconditional drift bounds with the real-variable existence of the mean growth exponent $\mu = \lim_{t\to\infty} t^{-1}\log \mathbb E[C(T_t)] = \lim_{n\to\infty} (\log n)^{-1}\log \mathbb E[C(T_n)]$ (\texttt{proofs/existence.md} Theorem 1, Astra) yields the certified enclosure $\lambda_{\mathrm{lo}}^* \le \mu \le \lambda_{\mathrm{up}}^*$.
\end{transcription}
% END TRANSCRIPTION round2/proofs/neon2/quintet-cert/PROOF.md

\subsection{The refined lower certificate}\label{app:refined}
This subsection proves the lower endpoint $a_-$ of Theorem~\ref{th:gold2}. It reuses the five-vertex reduction of Section~\ref{sec:enc} and its verifier unchanged. The one new step is Lemma~\ref{lem:classwise}, which combines the same verified per-vertex bounds class by class instead of through one global ratio.

\paragraph{The imported framework.}
A lower table of the five-vertex format fixes a nominal rate $\lambda=P/Q$, a scale $S=2^{40}$, a comparability shift $\alpha\in[0,1]$ stored in its header, and a degree cutoff $D_F$. For a finite tree $T$ with $C(T)\ge1$ core vertices, choose any core vertex as an accounting root. The potential is $\Phi(T)=S^{-1}\sum_uW(u)$, summed over core vertices, where $W(u)$ is the table value allocated to $u$: its singleton term; the term of the edge to its parent (none at the root); and its share of the supported higher motifs on three to five core vertices. A three-vertex path is allocated to its middle, three- and four-vertex stars and the five-vertex broom to their centres, and a five-vertex path to its middle. A four-vertex path is split, one half of its value and of its residual to each internal vertex. Vertices of degree at least $D_F$ host no higher motif. For a lower table accepted by the unmodified finite verifier, tail checker and combiner, the five-vertex reduction behind Theorem~\ref{th:gold2} \proofsrc{round2/proofs/neon2/quintet-cert/PROOF.md}{(full reduction, Sections 1--4)} gives the following on every finite tree.
\begin{itemize}\setlength{\itemsep}{0pt}
\item[(R1)] If $C(T)\ge2$, then $-S(\mathcal L\Phi-\lambda\Phi)(T)\le\sum_ue(u)$, where
\[
 e(u)=\max(0,e_f)+\max(0,e_g)+h(D_u,c_u).
\]
Here $e_f$ is the largest singleton row cap, including the analytic tail cap; $e_g$ is the largest edge row cap; and $h(D,c)\ge0$ bounds the total residual of all higher motifs hosted at a vertex of degree $D$ and core degree $c$, with $h=0$ for $D\ge D_F$. If $C(T)=1$, the residual is at most $\max(0,e_f)$.
\item[(R2)] With $m(u)=W(u)+\alpha S(c_u-2)$, one has $S\Phi(T)=2\alpha S+\sum_um(u)$, because $\sum_u(c_u-2)=2(C-1)-2C=-2$ on the core tree.
\item[(R3)] If $\Phi>0$, $mC\le S\Phi\le MC$ for constants $0<m\le M<\infty$, and $\mathcal L\Phi\ge\lambda'\Phi$ on every finite tree with a core vertex, then $\mu\ge\lambda'$. If $\lambda'\le0$ this follows from $C_n\ge1$ for $n\ge3$. If $\lambda'>0$, every factor below is positive, and this is the iteration in the proof of Theorem~\ref{th:gold2}: in discrete time $\mathcal L\Phi=n\,\E[\Delta\Phi\mid T_n]$, so $\E\Phi(T_n)\ge\Phi(T_3)\prod_{j=3}^{n-1}(1+\lambda'/j)$.
\end{itemize}
The bound (R1) uses the termwise estimates underlying the original combiner, with the singleton and edge caps replaced by their nonnegative parts: the residual of each supported motif is bounded once per embedding by its row cap, the transfers between rows cancel in pairs, and the caps are then allocated to host vertices as above. The verifier computes $h(D,c)$ from the row caps and from the number of motifs of each type that a vertex of core degree $c$ can host, for example at most $\binom c2$, $\binom c3$ and $\binom c4$ three-vertex paths, three-stars and four-stars, with the core degrees of the neighbours entering for the longer motifs. The original combiner uses $E=\max\bigl(0,e_f+\max(0,e_g)+\max_{D,c}h(D,c)\bigr)$ and the smallest mass $m_{\min}$, and certifies $\lambda-E/m_{\min}$. For the refined table $e_f=0$, so $E$ equals the largest per-core charge in (R1).

\begin{lemma}[Classwise aggregation]\label{lem:classwise}
Partition the core vertices into the following classes: the single core vertex of a tree with $C=1$; for each $1\le c<D_F$, vertices of degree $D<D_F$ and core degree $c$; and vertices of degree $D\ge D_F$ with core degree $1$, $2$ or at least $3$. In the single-core class take $e_0\ge\max(0,e_f)$ and $0<m_0\le\inf W(u)$ over single-core trees. For every other class $j$ take $e_j\ge\sup e(u)$ and $0<m_j\le\inf m(u)$ over the class. Put $\kappa=\max_je_j/m_j$, including the single-core class. If $\alpha\ge0$, then $\mathcal L\Phi\ge(\lambda-\kappa)\Phi$ on every finite tree with $C(T)\ge1$, and hence $\mu\ge\lambda-\kappa$.
\end{lemma}
\begin{proof}
If $C(T)\ge2$, then by (R1), the termwise bound $e(u)\le\kappa m(u)$, (R2) and $\alpha,\kappa\ge0$,
\[
 -S(\mathcal L\Phi-\lambda\Phi)\le\sum_ue(u)\le\kappa\sum_um(u)=\kappa(S\Phi-2\alpha S)\le\kappa S\Phi .
\]
If $C(T)=1$, then $S\Phi=W(u)\ge m_0$ for the single core $u$, and the residual is at most $\max(0,e_f)\le e_0\le\kappa m_0\le\kappa S\Phi$. Dividing by $S$ gives the drift bound. For $C\ge2$, (R2) and $\alpha\ge0$ give $S\Phi\ge\sum_um(u)\ge(\min_jm_j)C>0$; for $C=1$, $S\Phi\ge m_0>0$. The upper comparison $S\Phi\le MC$ of the table is unchanged. Now (R3) applies with $\lambda'=\lambda-\kappa$.
\end{proof}

\paragraph{The numbers.}
The refined lower table has nominal rate $\lambda=56571437/10^8$, cutoff $D_F=4096$ and $\alpha S=310153455616$. The unmodified checker accepts it in lower mode. Its receipt records $e_f=\max(-4077,0,0)$, so $\max(0,e_f)=0$, and $e_g=1873$. It also records $\max_{D,c}h(D,c)=20325$, so $E=22198$, and $m_{\min}=135889307074$. The original combiner therefore certifies only
\[
 \lambda-\frac{22198}{135889307074}=\frac{3843725577155222669}{6794465353700000000}<\frac{99}{175}.
\]
The per-class data come from an instrumented copy of the same verifier, with four added lines in three places that record, for each core degree $c<D_F$, the minimum over $D$ of the unshifted value lower bound and the maximum over $D$ of $h(D,c)$. These are values the unmodified code already computes. The added lines change no check and no output field, and the instrumented run reproduces the receipt's verification record exactly. With $m=\min_D(\text{value lower bound})+\alpha S(c-2)$ on the finite classes, and the tail and single-core masses from the receipt, the classes are:
\begin{center}\small
\begin{tabular}{@{}lrr@{}}
\toprule
class & charge $e_j$ & mass $m_j$\\
\midrule
single core ($C=1$), mass $W$ & $0$ & $1577747427683$\\
$D\ge D_F$, $c=1$ & $1873$ & $1324075825814$\\
$D\ge D_F$, $c=2$ & $1873$ & $1634069402427$\\
$D\ge D_F$, $c\ge3$ & $1873$ & $135889307332$\\
$D<D_F$, core degree $c$ & $1873+\max_Dh(D,c)$ & $\min_D(\text{value})+\alpha S(c-2)$\\
\bottomrule
\end{tabular}
\end{center}
Every mass is positive. The largest ratio occurs at $D<D_F$, $c=2$, where $\max_Dh(D,2)=1941$ and the mass is $135889307074$. Hence
\[
 \kappa=\frac{1873+1941}{135889307074}=\frac{1907}{67944653537},\qquad
 \lambda-\kappa=\frac{3843726496355222669}{6794465353700000000}=a_- .
\]
The gain over the original combiner comes from where the charges sit: the large hosted caps sit at vertices of high core degree $c$, whose mass carries the extra $\alpha S(c-2)$, while the smallest masses sit at $c\le2$, which host little.

\paragraph{What this step assumes.}
The lower endpoint rests on the five-vertex reduction and on the verifier, tail checker and combiner as its implementation, exactly as for the Gold2 pair; on Theorem~\ref{th:exist}, which names the exponent; on Lemma~\ref{lem:classwise}; and on the witness bytes identified in Section~\ref{app:replay}. The upper endpoint is the unchanged Gold2 upper endpoint; no new upper table was built. The choice of $\lambda$ and $\alpha$ is not claimed to be optimal.

\sectag{CORE}
\section{The marked kernel, the amplitude and the degree law}\label{sec:core}\label{sec:core-details}
We use the exponent and first enclosure of Theorem~\ref{th:exist}. The first attachment from the founding edge creates the first core vertex, which has two leaf neighbours. From size three onward a promotion occurs when a core vertex selects an old leaf as the attachment recipient. The old leaf becomes a core vertex with one core neighbour (the selecting vertex) and one leaf neighbour (the newly attached vertex). Write $b_x=L_x/D_x$ for the promotion rate at $x$, $F(t)$ for the expected core count from the single edge at time $t$, and $C(T)$ for the number of core vertices in a tree $T$.

At size $j$ the holding time is $\operatorname{Exp}(j)$, independent of the next arrow. The holding times can thus be sampled independently of the entire discrete chain. At time $t$ the size from $m$ vertices is a sum of $m$ independent geometric variables on $\{1,2,\ldots\}$ with parameter $e^{-t}$, by the pure-birth equations and superposition. If $\tau_n$ is the size-$n$ epoch, then $\tau_n-\log n$ converges almost surely by the summable variances of the independent waits and the harmonic sum asymptotic. Hence $e^{-t}|T(t)|\to W_{\rm size}>0$, whose law is $\operatorname{Gamma}(m,1)$ (shape, rate), by the geometric Laplace transforms. We often reset time at $T_3$, the path at the first birth, and orient from its centre. By the independent first wait and monotonicity, $F(t)\le F_3(t)\le F(t+1)/(1-e^{-2})$, where $F_3(t)=\mathbb E_{T_3}C(T(t))$. In particular $t^{-1}\log(F_3(t)-1)\to\mu$.

\subsection{The marked kernel}
\tagnext{KER}
\begin{theorem}[Compact reproduction kernel]\label{th:core-kernel}
Set $\mathcal A=\{0\}\cup\{1/j:j\ge2\}$. The positive marked-fringe reproduction operator $K_s$ of \eqref{eq:core-kernel}, on $C(\mathcal A)$, is compact for $s>51/100$. Its spectral radius crosses one exactly once on that interval, at $s=\mu$. In particular $\mathbb E C_n\asymp n^\mu$.
\end{theorem}
\begin{proof}
Lemma~\ref{lem:core-renewal} identifies the critical discount using the logarithmic exponent of Theorem~\ref{th:exist}; it does not reprove existence of that exponent. The comparison with the physical mean uses $\mu\ge113/200>51/100$, supplied by the computer-checked three-vertex lower witness of Theorem~\ref{th:exist}, not by the five-vertex enclosure. Long's earlier characterization uses the generation growth rate of discounted birth sums from an inert-parent newborn \cite{LongKernel}. Here a generation counts direct births in the auxiliary marked families; the kernel integrates over intrinsic evolution and copy events. The two characterizations identify the same discount, not necessarily the same function of the discount.

Long first proved the constant-factor mean bounds \cite{LongMean}. Here is the route from the marked kernel: Lemma~\ref{lem:core-uniform} gives $\E_{T_3}C(T(t))\asymp e^{\mu t}$, since $d_2\Psi(T_3)>0$. The independent rate-two first wait gives $F(t)\asymp e^{\mu t}$. For $a_n=\E C_n$ and $m\ge n\ge3$, the promotion probability is at most $C_j/j$ at size $j$, so $a_m\le (m/n)a_n$ by iteration. Also $a_n$ is nondecreasing, and $F(t)=\E a_{|T(t)|}$ by clock independence. At $t=\log n$ the probability of size at least $n$ is at least $e^{-1}$ by one initial Yule family, and the mean size is $2n$. Thus $e^{-1}a_n\le F(\log n)\le3a_n$. Theorem~\ref{th:core-means} sharpens this to fixed-degree mean equivalents. Source: \proofsrc{round2/proofs/astra/source/round2/fringe/renewal.md}{marked-fringe renewal}.
\end{proof}

\subsection{The amplitude}\label{sec:core-amp}
\tagnext{AMP}
\begin{theorem}[Core amplitude]\label{th:core-amplitude}
From one edge, $C_n/n^\mu\to X$ almost surely, where $X$ is finite, strictly positive almost surely, and square integrable. For every $\varepsilon,M>0$, $\mathbb P(X<\varepsilon)>0$ and $\mathbb P(X>M)>0$. Simultaneously for all fixed integers $k\ge2$,
\[
 \frac{N_{\ge k}(n)}{C_n}\longrightarrow\frac{c_k}{c_2},\qquad
 \frac{N_{=k}(n)}{C_n}\longrightarrow\frac{c_k-c_{k+1}}{c_2}
 \quad\text{a.s.},\qquad \mathbb E X=c_2,
\]
where the $c_k$ are the mean coefficients of Theorem~\ref{th:core-means}.
\end{theorem}
\begin{proof}
Use the finite-mark uniform mean expansion of Lemma~\ref{lem:core-uniform}: for all finite seeds with a core, $\mathbb E_T N_{\ge k}(T(t))=e^{\mu t}d_k\Psi(T)+O(C(T)e^{\sigma t})$, with $\sigma<\mu$, $d_2>0$, and $C\asymp\Psi$ uniformly. The normalized mean at $k=2$ identifies $\Psi$ intrinsically. Applying the expansion after a fixed time $u$ and sending the remaining time to infinity (the expected size at $u$ is finite) gives $\mathbb E_T\Psi(T(u))=e^{\mu u}\Psi(T)$. Accordingly $M_u=e^{-\mu u}\Psi(T(u))$ is a nonnegative c\`adl\`ag martingale.

We give the concentration argument, since a first-moment expansion alone does not give the asserted limit. Compare a seed $T$ with a core to $T^+$, obtained by adding one leaf. If the recipient is a core, mark its two copies dirty. If the recipient was a leaf, mark both copies of its old parent and the newly promoted core in $T^+$ dirty. There are at most three initially dirty cores, counting both trees. A clean pair has matched leaf stocks and matched neighbour slots, including a corresponding slot for each adjacent dirty core. Share the clocks and slot choices of clean pairs, and share the clocks of their paired leaves. A clean ring choosing a leaf promotes a matched pair; choosing a core sends into the corresponding slot in both trees, so no clean pair becomes dirty. Use separate clocks for dirty cores and their leaves. A dirty core can promote one private leaf or send into a clean core, in which case mark both copies of that recipient dirty. A leaf of a dirty core can send only into its parent and does not spread dirt. These updates preserve the matching on all remaining clean pairs.

If $J_t$ is the combined dirty count, its upward drift is at most $2J_t$: only the $J_t$ rate-one dirty core clocks can increase it, by at most two per ring. Stop at a bound on the combined size, apply compensation and Gr\"onwall, then remove the stop by nonexplosion and Fatou to obtain $\E J_t\le3e^{2t}$. For $k\ge2$ every clean pair contributes identically to the two cumulative counts, so the absolute difference of their means is bounded by $\E J_t$.

For a fixed observation time $b$, the backward martingale of $N_{\ge k}(T(b))$ jumps at a birth at time $u$ by this one-leaf difference at lag $b-u$. Between births its mean is differentiable by the first-event equation. Compensating squared jumps (localize by size and use the finite-time Yule moments) gives
\[
 \operatorname{Var}_T N_{\ge k}(T(b))\le9|T|\int_0^b e^{u+4(b-u)}\,du.
\]
Write $\eta=\mu-\sigma$ and $g_h(T)=d_2^{-1}e^{-\mu h}\mathbb E_T C(T(h))$. The same jump bound on $[0,1]$ for the backward martingale of $g_h(T(1))$ uses lag $h+1-u$, giving variance $O(e^{4h}|T|)$. The expansion bounds $|g_h-\Psi|$ by $O(e^{-\eta h}C)$, and $\mathbb E_T C(T(1))^2\le2e^2 C(T)^2$ by rate-one per-core Yule domination. Thus
\[
 \operatorname{Var}_T\Psi(T(1))\le K\big(e^{4h}|T|+e^{-2\eta h}C(T)^2\big).
\]
At integer time $j$ take $h=\eta^{-1}\log(j+2)$. By conditional variance, $C\lesssim\Psi$ and $\mathbb E_T|T(j)|=|T|e^j$,
\[
 \mathbb E_T M_{j+1}^2\le(1+K'(j+2)^{-2})\mathbb E_T M_j^2
 +K' |T|(j+2)^{4/\eta}e^{-(2\mu-1)j}.
\]
The first enclosure gives $2\mu>1$, so both error series are summable. Jensen between integers and martingale convergence give $M_t\to M_\infty$ a.s. and in $L^2$, with $\mathbb E_T M_\infty=\Psi(T)$ and $\mathbb E_T M_\infty^2\lesssim\Psi(T)^2+|T|$. For the edge chain with time reset at $T_3$, the logarithmic path exponent implies $|T(j)|/\Psi(T(j))^2\to0$ a.s. Conditional Cauchy--Schwarz gives the lower bound
\[
 \mathbb P(M_\infty>0\mid\mathcal F_j)\ge
 \frac{\Psi(T(j))^2}{K''(\Psi(T(j))^2+|T(j)|)}.
\]
Its lower limit is positive; the conditional probabilities converge to the event indicator, proving positivity.

On a fixed mesh of spacing $1/m$, use lag $h=\eta^{-1}\log(t+2)$ and condition at $u=t-h$, for mesh times large enough that $u\ge0$. The variance bound at lag $h$ is at most $3|T(u)|e^{4h}$, since $\int_0^h e^{-3v}\,dv\le1/3$. Taking expectations using $\E |T(u)|=3e^u$ and multiplying by $e^{-2\mu t}$ bounds the squared centered error by $O(e^{-(2\mu-1)t+3h})$. This is summable on the mesh: the lag contributes only a power of $t+2$, and the three-vertex lower witness gives $2\mu>1$. The scaled conditional mean differs from $d_k M_u$ by $O(e^{-\eta h}e^{-\mu u}C(T(u)))=o(1)$ a.s. Chebyshev and Borel--Cantelli give $e^{-\mu t}N_{\ge k}(T(t))\to d_k M_\infty$ on the mesh. Intersect over all $k,m$ and interpolate the monotone cumulative counts with factors $e^{\pm\mu/m}$, then let $m\to\infty$. At the size epochs the limits are $d_k W_{\rm size}^{-\mu}M_\infty$, simultaneously. Exact counts follow by subtraction.

To pass to means and show the support assertions, the harmonic identity and the first holding time at size $n$ give $\mathbb E[\Psi(T_{n+1})\mid T_n]=(1+\mu/n)\Psi(T_n)$ (equivalently, differentiate the first-event identity at zero). Thus $Y_n=\Gamma(n)\Psi(T_n)/\Gamma(n+\mu)$ is a discrete martingale. At $\tau_n$, conditional Jensen for the $L^2$-closed continuous martingale gives $\mathbb E M_{\tau_n}^2\le\mathbb E M_\infty^2$. Independence of the size epochs and the discrete chain gives $\mathbb E M_{\tau_n}^2=\mathbb E\Psi(T_n)^2\prod_{j=3}^{n-1}j/(j+2\mu)$, with product $\asymp n^{-2\mu}$. This bounds $Y_n$ in $L^2$. To identify its limit, use the same clock reset at $T_3$ in $\tau_n$, $M_t$ and $W_{\rm size}$:
\[
 Y_n=\frac{\Gamma(n)}{\Gamma(n+\mu)}e^{\mu\tau_n}M_{\tau_n}
 \longrightarrow W_{\rm size}^{-\mu}M_\infty=X/d_2
 \qquad\text{a.s.}
\]
Here $ne^{-\tau_n}\to W_{\rm size}$ and $\Gamma(n)/\Gamma(n+\mu)\sim n^{-\mu}$. The clock is independent of the discrete chain, not of $M_\infty$. Discrete $L^2$ martingale convergence now also gives convergence in $L^2$. In particular $\mathbb E[X\mid\mathcal F_n^{\rm disc}]=d_2Y_n$. Since $C\lesssim\Psi$, the scaled cumulative counts are uniformly integrable, and
\begin{equation}\label{eq:core-ck}
 c_k=d_k\,\frac{\Gamma(3)}{\Gamma(3+\mu)}\Psi(T_3).
\end{equation}
The ratio claims follow. At every fixed size $n\ge3$, both the star and path have positive probability: preserving the star has probability $(j-1)/j$ at size $j$, preserving the path has probability $1/j$. Their core counts are $1$ and $n-2$. The conditional mean of $X$ is $O(n^{-\mu})$ on the star and bounded below by a positive multiple of $n^{1-\mu}$ on the path. Comparing these means with $\varepsilon$ and $M$ proves the two support assertions. Source: \proofsrc{round2/proofs/astra/source/round2/standalone/claim-6.md}{uniform renewal and strong limit}.
\end{proof}

\subsection{Fixed-degree means and the degree law}\label{sec:core-means}
\tagnext{DEG}
\begin{theorem}[Fixed-degree means and the non-leaf degree law]\label{th:core-means}
For each fixed integer $k\ge2$, $\mathbb E N_{\ge k}(n)\sim c_k n^\mu$, where $c_k>c_{k+1}>0$ and $c_k\downarrow0$. With $p_k=(c_k-c_{k+1})/c_2$, one has $\sum_{k\ge2}p_k=1$ and
\[
 \lim_{n\to\infty}\frac{\mathbb E N_{=k}(n)}{\mathbb E C_n}=p_k,
 \qquad \sum_{j\ge k}p_j\sim A k^{-\mu},\qquad p_k\sim\mu A k^{-1-\mu}
 \quad(k\to\infty),
\]
for some $0<A<\infty$.
\end{theorem}
\begin{proof}
The mean equivalent and coefficient formula were proved in Theorem~\ref{th:core-amplitude}; Lemma~\ref{lem:core-uniform} gives the strict inequalities and decay of $d_k$, hence of $c_k$. Subtracting cumulative equivalents at fixed $k$ and telescoping yields the ratio and normalization. Lemma~\ref{lem:core-spatial} constructs a probability $\pi$ for the core-cell evolution with an exponential moment of root core degree $c_o$. Only these unconditional statements of Lemma~\ref{lem:core-spatial} are used; the hypotheses of Lemma~\ref{lem:core-ports} are not. Under its mean tilts, bounded time intervals have vanishing weight and the fixed-$k$ normalized mean tends to $d_k/d_2$, so $\pi(D_o\ge k)=c_k/c_2$. We now compute its tails; this does not assume a unique law of the rooted environment.

Write $S_a$ for the evolution tracking just the same old root, and $\mathcal Bf$ for $b_o$ times the score at a newborn rooted at the leaf promoted by $o$, in the post-promotion environment. Each cell member other than the old root has one promotion time. Compensating these births and using the conditional law after promotion gives $Q_t f=S_t f+\int_0^t Q_u\mathcal B S_{t-u}f\,du$ for bounded nonnegative $f$, where $Q_t$ is the cell sum from Lemma~\ref{lem:core-spatial}. The Yule bound justifies the integral. Apply $\pi Q_u=e^{\mu u}\pi$, divide by $e^{\mu t}$ and let $t\to\infty$. The old-root term vanishes. At $f=1$ we obtain $\pi b_o=\mu$, and for general $f$,
\[
 \pi f=\int_0^\infty \mu e^{-\mu a}\,\beta S_a f\,da,
 \qquad \beta=\mu^{-1}\pi\mathcal B.
\]
Thus the tracked root starts at a genuine promotion (degree two, one leaf, real parent) and is observed at independent exponential age of rate $\mu$.

The degree $D(t)$ of this root jumps by one at rate at most $D(t)$. Thus $e^{-t}D(t)$ is a nonnegative supermartingale with a limit $W$. By adding jumps it is dominated before normalization by a Yule chain from two; Doob's $L^2$ inequality and the geometric second moment bound give an integrable square for the supremum after normalization (at most $24$ in expectation). For positivity, silence the newborn's own promotion driver for one unit and let its tagged leaf ring at least twice, an event of probability at least $e^{-1}(1-2e^{-1})$. Then $L(1)\ge3$. The subsequent drift of $L$ is at least $L-1$, so $\mathbb E[L(1+t)-1\mid\mathcal F_1]\ge e^t(L(1)-1)$ by the integrating factor and stopped compensation. The same Yule comparison from the finite degree at time one gives conditional uniform integrability. It follows that $\mathbb E[eW\mid\mathcal F_1]\ge L(1)-1$ there. In particular $A=\mathbb E_\beta W^\mu\in(0,\infty)$.

For $T_k=\inf\{t:D(t)\ge k\}$ the age mixture gives $\pi(D_o\ge k)=\mathbb E_\beta[e^{-\mu T_k};T_k<\infty]$. On unbounded-degree paths $D(T_k)=k$ and $T_k\to\infty$, so $ke^{-T_k}\to W$; on bounded-degree paths the terms are eventually zero and $W=0$. The maximal bound gives dominated convergence and the cumulative tail. For a point probability we need a separate flux identity, not subtraction of two tail equivalents. For $k\ge2$ the newborn has degree below $k+1$, and the count above $k$ increases only by a jump from degree $k$. Compensation in the cell and the eigenmeasure identity give
\[
 \mu\pi(D_o\ge k+1)
 =\pi\left[\mathbf1_{\{D_o=k\}}\left(k-c_o+\sum_{y\sim o\,\mathrm{in\ core}}D_y^{-1}\right)\right].
\]
The rate on $D_o=k$ is at most $k$, and at least $(1-\epsilon)k$ on $c_o\le\epsilon k$. The complementary probability decays exponentially by the core-degree moment. Multiply the resulting bounds by the suitable power of $k$: the liminf of $k^{1+\mu}p_k$ is at least $\mu A$, and its limsup is at most $\mu A/(1-\epsilon)$. Send $k\to\infty$ and then $\epsilon\downarrow0$. Sources: \proofsrc{round2/proofs/astra/source/round2/standalone/claim-3.md}{fixed-size mean and tail proofs} and \proofsrc{round2/proofs/astra/source/round2/degree-law.md}{birth age and crossing flux}.
\end{proof}

\paragraph{The degree-ratio conjecture.}
Krapivsky and Redner \cite[Sec.~5.1, Eqs.~(32)--(33)]{KR} predict sublinear growth of the mean degree counts and a power-law tail for their normalized counts $N_k/\mathcal N$, where $N_k$ counts degree exactly $k$ and $\mathcal N=\sum_{j\ge2}N_j$ is the nucleus size, not the full tree size. In our notation the ratio of means is $\E N_{=k}(n)/\E C_n$. Addario-Berry et al.\ \cite[Sec.~7, Question~2]{AB} attribute to Krapivsky and Redner the prediction that each fixed-degree ratio among non-leaves has an almost-sure limit and that these limits have order $k^{-(1+\mu)}$. In their notation $X_n^k=N_{=k}(n)$ and $X_n^{\ge2}=C_n$; the proposed limit and tail are
\[
  \frac{X_n^k}{X_n^{\ge2}}\longrightarrow p_k \quad\text{almost surely for each fixed } k\ge2,
  \qquad p_k=\Theta(k^{-(1+\mu)})\quad(k\to\infty).
\]
Theorems~\ref{th:core-amplitude} and~\ref{th:core-means} give this limit and the representation $\lim_{k\to\infty}k^{1+\mu}p_k=\mu A$. This represents the coefficient without evaluating $A$ numerically. First let $n\to\infty$ at each fixed $k$, then let $k\to\infty$ in the resulting law. Theorems~\ref{th:core-amplitude} and~\ref{th:core-means} do not give a joint pointwise exact-degree equivalent with $k=k(n)$; Section~\ref{sec:further} records a supplementary equivalent for the mean only. Since every $p_k$ is positive, consecutive exact-degree counts have the same order for fixed $k$. Since $\sum_{j\ge k}p_j\to0$, the limiting fraction above a large degree tends to zero. These answer both bounded-degree questions in Section~7(2).

\subsection{Renewal, uniform expansion and the spatial law}
\begin{lemma}[Marked-fringe first moments]\label{lem:core-renewal}
In any finite ordinary seed with a core, orient away from a chosen core founder $o$. At the promotion of a leaf by a degree-$D$ core assign the new core the permanent mark $a=1/m(D)$, with either $m(D)=D$ or $m(D)=\min(D,20)$. Initialize old nonfounders with compatible marks, for instance all $1/2$. A child's mark always dominates the parent's current reciprocal degree. The full and residual first-moment semigroups below satisfy \eqref{eq:core-fringe-identity} and \eqref{eq:core-weights}. For either choice of marks, the reproduction kernel has its unique critical discount at $\mu$ on $(51/100,\infty)$.
\end{lemma}
\begin{proof}
A planted type $S$ at a nonfounder contains its descendants and an exterior core slot at its root, counted in both total and core degree. The intrinsic generator $L^0$ suppresses incoming arrows from that slot and gives a null transition on choosing it; all internal arrows keep their physical rates. Write $p_S=1/D_{\rm root}$, $b_S=L_{\rm root}/D_{\rm root}$, $a_S$ for the root mark and $Y$ for immediate core-child types counted with multiplicity. Let $Y^+$ add a leaf to the root of $Y$, and $J_a$ be the planted newborn with one leaf and mark $a$. The auxiliary root mark zero is used only for the uncapped kernel, by continuity, and is never assigned at a physical birth. In the residual family a particle changes intrinsically, dies at rate $a_S$, and emits additional particles $Y,Y^+$ at rates $a_Y-p_S,p_S$ respectively. The full family additionally emits $J_{1/m(D_{\rm root})}$ at rate $b_S$. The generator on residual scores is $Rf=L^0f-a_Sf(S)+\sum_Y[(a_Y-p_S)f(Y)+p_Sf(Y^+)]$. Denote by $R_t,P_t$ the residual and full first-moment semigroups. Write $(B_0g)(S)=b_Sg(1/m(D_{\rm root}))$ and $(Jf)(a)=f(J_a)$ so the full generator is $R+B_0J$.

For a physical tree $T$ the founder supplies the row
\[
 \mathcal I(T)f=b_o f(J_{1/m(D_o)})+\sum_{Y\text{ at }o}[(a_Y-p_o)f(Y)+p_o f(Y^+)],
 \qquad I_{T,u}f=\mathbb E_T\mathcal I(T(u))f.
\]
In the sum over nonfounder fringes, a gain into an old root from its parent is missing in its intrinsic row. For an edge from a parent of type $S$ to its old core child of type $Y$, the two terms in the parent row sum with the death term in the child row to $(a_Y-p_S)f(Y)+p_Sf(Y^+)-a_Yf(Y)=p_S[f(Y^+)-f(Y)]$. This is exactly the missing leaf-gain term for that directed arrow. At a founder edge the same terms come from $\mathcal I(T)$ and the child row. Promotions also create new terms $f(J_a)$, accounted for at the promoter. The identity is therefore
\begin{equation}\label{eq:core-fringe-identity}
 \mathbb E_T\sum_{v\ne o,\,\mathrm{core}}f(S_v(t))
 =\sum_{v\ne o,\,\mathrm{core\ at\ }0}P_tf(S_v(0))+\int_0^t I_{T,u}P_{t-u}f\,du.
\end{equation}
This is a first-moment identity, not a coupling of independent physical fringes.

For a weight take $d=3/20$, $q=9/25$, $\rho=q+d=51/100$, and $h(0),h(1),h(2),h(3),h(c\ge4)=2,1,17/20,39/50,7/10$. Let $w_x$ be the product of $a_y/(a_y+d)$ on the path to $x$ excluding the root; put $W(S)=(a_S+d)^{-1}\sum_x w_xh(c_x)$ over cores. Factoring the first edge gives $\sum_Y a_YW(Y)=(a_S+d)W(S)-h(c_{\rm root})$. Gains of leaves do not change $W$. At a promotion, the rate times the change divided by the old path factor and root normalization is
\[
 (1-c/D)\big[h(c+1)-h(c)+(1+d m(D))^{-1}\big]\le qh(c).
\]
For $D\ge20$ the left side is at most $1/4<q(7/10)$, for either marking convention. Otherwise the uncapped expression suffices. For $c=0$ it is nonpositive. For $c\ge4$ it decreases at fixed $D$ with $c$, so only $c=1,2,3,4$ need checking. With $r=qh(c)$, $v=h(c+1)-h(c)$, multiplication by $D(1+dD)$ gives the quadratic $(r-v)dD^2+(r-v-1+cdv)D+c(1+v)$ for the difference. For $c=2,3,4$ its leading coefficient is positive and the discriminant is respectively $-3591/10^6,-110617/781250,-2831/62500$. For $c=1$ its positive leading coefficient is $153/2000$, its vertex lies between $3$ and $4$, and its values at those integers are $1/1000,3/125$. Thus the inequality holds analytically. In particular $L^0W\le qW$, $RW\le\rho W$, $B_0JW\le1/d$, and $1<W(S)\le |S|/d$. The corresponding stopped exponential bounds on total weight rule out an unbounded count in bounded time. With bounded count the death intensity is bounded, hence only finitely many emitted particles are possible; each intrinsic process has size rate at most its size. Removing stops by Fatou gives
\begin{equation}\label{eq:core-weights}
 R_tW\le e^{\rho t}W,\qquad P_tW\le e^{(\rho+1/d)t}W,
 \qquad I_{T,t}W\le 1/d+2C(T)e^{qt}.
\end{equation}
For the last bound use the same promotion inequality on the physical sum $\sum_x w_xh(c_x)$ without the denominator, initially at most $2C(T)$. The copy terms at the founder have combined weight equal to that sum minus $h(c_o)$, and the direct birth costs at most $1/d$.

For $|f|\lesssim W$, the first-event integral on each finite type gives a continuous $P_v f$ and its pointwise backward derivative: the initial rate and the weight after any one event are $O(|S|+1)$, giving a derivative bound $O_t((1+|S|)^2)$ up to time $t$, using \eqref{eq:core-weights}. Write $A_g(T)=\sum_{v\ne o,\,\mathrm{core}}g(S_v)$ and let $\mathcal L$ be the physical generator. The edge cancellation gives $\mathcal L A_g=A_{(R+B_0J)g}+\mathcal I(T)g$. Now put $g_u=P_{t-u}f$. By the first-event backward equation $\partial_u g_u=-(R+B_0J)g_u$, and hence
\[
 (\partial_u+\mathcal L)A_{g_u}(T)=\mathcal I(T)g_u.
\]
Integrate this identity along the physical chain stopped at a finite size. The polynomial bounds on scores, derivatives and jump increments times intensity are dominated by finite-time Yule moments, so dominated convergence removes the stop. The endpoint scores are $A_f(T(t))$ and $A_{P_tf}(T)$; the expected integral is the immigration term of \eqref{eq:core-fringe-identity}.

For $\Re z>\rho$ put $G_z f=\int_0^\infty e^{-zt}R_tf\,dt$, $U_z=G_z B_0$, and
\begin{equation}\label{eq:core-kernel}
 K_z=JU_z,\qquad (K_zg)(a)=\int_0^\infty e^{-zt}R_t(B_0g)(J_a)\,dt.
\end{equation}
Integrals on scores are componentwise. In weighted supremum norm, $\|G_z\|_W\le(\Re z-\rho)^{-1}$, and the exponential Taylor series under this bound gives norm holomorphy without assuming strong continuity of $R_t$ in that norm. To see compactness in the uncapped case, separate the original particle from the families of its emitted copies. Its own mark $a$ affects only its death clock. The row is the integral with factor $e^{-(z+a)t}$ of the intrinsic expectation of
\[
 b_Sg(1/m(D_{\rm root})) +\sum_Y[(a_Y-p_S)U_zg(Y)+p_SU_zg(Y^+)].
\]
For $\|g\|_\infty\le1$, the absolute expectation is at most $1+e^{qt}/(\Re z-\rho)$ by the child-weight identity. Thus rows are Lipschitz in $a$ uniformly for bounded $g$, including at $0$ (multiply the integrand bound by $te^{-t\Re z}|a-a'|$). Arzel\`a--Ascoli gives compactness. Likewise $JG_s1$ is continuous by the same decomposition with original toll $1$ and continuation $G_s1$.

At each real $s>\rho$, every row charges every nonzero newborn mark: the original root can survive, gain leaves alone to the corresponding degree, and then hold there with positive direct-birth rate. Thus the operator is strongly positive on the cone in $C(\mathcal A)$ (and strictly positive as a finite matrix in the capped case). The compact strongly positive Krein--Rutman theorem, or finite-dimensional Perron--Frobenius in the capped case, gives an algebraically simple spectral radius and a positive eigenfunction with positive minimum. For $s_1<s_2$, $(K_{s_1}-K_{s_2})g_2$ on the positive $s_2$ eigenfunction is strictly positive at every mark, and dominates a strictly positive multiple of $g_2$, proving strict decrease by iteration. Norm continuity and comparison on the fixed eigenfunction at $s$ give continuity of the radius.

Discounted full-family occupation at newborns, grouped by the number of direct births in the lineage, is $\sum_{j\ge0} K_s^j JG_s1$, with positive forcing bounded below by $(s+1/2)^{-1}$ by original-particle survival. It converges in norm if $r(K_s)<1$; the analogous arbitrary-seed integral is then bounded by a multiple of $W$ by decomposing at the first direct birth. Otherwise the forcing dominates a multiple of the positive eigenfunction and the series diverges at every mark. From $T_3$ there are no initial nonfounder fringes, and the founder has positive direct immigration in mean on an interval. By Tonelli in \eqref{eq:core-fringe-identity} and \eqref{eq:core-weights}, the physical Laplace integral of $F_3(t)-1$ is finite exactly when $r(K_s)<1$. Its abscissa is $\mu>\rho$ by the earlier logarithmic exponent. Continuity and strict decrease give the claimed criticality. Sources: \proofsrc{round2/proofs/astra/source/round2/fringe/renewal.md}{renewal identity and weight}, \proofsrc{round2/proofs/astra/source/round2/fringe/finite-kernel.md}{finite permanent marks, Section 1}. No finite truncation of the degrees or of the offspring histories is used.
\end{proof}

\begin{lemma}[Uniform exponential mean remainder]\label{lem:core-uniform}
Use $m(D)=\min(D,20)$ in Lemma~\ref{lem:core-renewal}. There are $\rho<\sigma<\mu$ and a bounded positive linear functional $d(f)$ on scores of finite norm $\|f\|_W=\sup_S|f(S)|/W(S)$, such that for every bounded root score depending only on $(L,c)$ and every finite ordinary seed $T$ with a core,
\begin{equation}\label{eq:core-unif}
 \mathbb E_T\sum_{x\,\mathrm{core}} f(L_x(t),c_x(t))
 =e^{\mu t}\Psi(T)d(f)+O(\|f\|_\infty C(T)e^{\sigma t}),
 \qquad 0<aC(T)\le\Psi(T)\le b C(T),
\end{equation}
with constants independent of the seed and score. Treat the same $f$ on planted types as the score at their root. For $f_k=\mathbf1_{\{L+c\ge k\}}$, the numbers $d_k=d(f_k)$ satisfy $d_k>d_{k+1}>0$ and $d_k\downarrow0$.
\end{lemma}
\begin{proof}
The matrix $K_z$ has nineteen coordinates (marks $1/2,\ldots,1/20$); all actual send rates are unchanged. Put $k(t)=JR_tB_0$. A high-frequency bound follows by differentiating the residual mean of $B_0g$, $\|g\|_\infty\le1$. For $D<20$ the root gain rate times the absolute score change is bounded by $2D\le38$. For $D\ge20$ the label in the toll stays fixed, and the absolute change is at most $c/[D(D+1)]$, with gain rate $\le D$. A root promotion changes the toll by at most $1/D$ and has rate at most one; other intrinsic jumps do not change it. Kill and copy terms contribute at most $1/2+(13/20)W(S)$ by the weight identity. Hence the total absolute score-change intensity is at most $42W$. In a residual family its expected integral on bounded intervals is finite by \eqref{eq:core-weights}; stopped compensation and removal of the stops give $R_t B_0g=B_0g+\int_0^t R_uRB_0g\,du$. Integration by parts and $JW\le 1/d$ yield $\|K_z\|\le (1+280/\epsilon)/|z|$ for $\Re z\ge\rho+\epsilon>\rho$, in the induced row-sum norm. The same estimate before applying $J$ gives $\|U_z\|=O_\epsilon(|z|^{-1})$ in weighted norm.

Let $g,\ell$ be positive right and left Perron vectors at $\mu$, with $\ell g=1$. The contour projection around the isolated simple eigenvalue one is locally analytic by the inverse series, so there is an analytic eigenvalue branch with derivative $\ell K'_\mu g<0$: its integral weights births by minus their positive ages. Thus $(I-K_z)^{-1}$ has a simple pole at $\mu$ of residue $D_*=(g\otimes\ell)/(-\ell K'_\mu g)$. On $\Re z>\mu$ the Neumann series works by domination by the real-discount matrix. On $z=\mu+i\tau$, if $K_zv=v$ then $|v|\le K_\mu|v|$. The positive vector $\ell$ forces equality, hence $|v|$ is a positive multiple of $g$ for a nonzero $v$. Every row contains a mark-$1/2$ birth contribution with strictly positive time density on $(0,\infty)$: keep the original newborn unchanged and alive, with direct-birth rate $1/2$. For $\tau\ne0$ this contribution makes the triangle inequality strict because $e^{-i\tau t}v(1/2)$ does not have constant phase. This rules out any other pole on the line. The high-frequency bound and compactness of the bounded-frequency portion of the line, together with the local simple pole decomposition, give a single $\sigma\in(\rho,\mu)$ such that the closed half-plane $\Re z\ge\sigma$ contains no other pole and the inverse is bounded on its boundary.

Write $v(t)=\sum_{j\ge1} k^{*j}(t)$ with matrix multiplication in time order. First-event continuity and the bound on $k$ on compact time intervals give locally uniform convergence (bound the $j$th term by $B^j t^{j-1}/(j-1)!$). The Laplace transform of $v-k$ is $(I-K_z)^{-1}K_z^2$ on $\Re z>\mu$. The time function is continuous (including zero after extension by zero) and after exponential discount integrable by the positive generation series. The transform is $O(|\Im z|^{-2})$ at large height uniformly between $\sigma$ and any fixed $s>\mu$. With $G(z)=(I-K_z)^{-1}K_z^2$, Fourier inversion of the discounted time function on $\Re z=s$ and the residue theorem give
\[
 v(t)-k(t)=\frac1{2\pi}\int_{\mathbb R}e^{(s+iy)t}G(s+iy)\,dy
 = e^{\mu t}D_*+\frac1{2\pi}\int_{\mathbb R}e^{(\sigma+iy)t}G(\sigma+iy)\,dy.
\]
Indeed the horizontal integrals at height $\pm L$ are $O((s-\sigma)e^{st}/L^2)$ for fixed $t$, and $G$ has residue $D_*K_\mu^2=D_*$ since $\ell K_\mu=\ell$. On $\sigma$ the integrand after removing $e^{\sigma t}$ is bounded in norm by the integrable function $\|G(\sigma+iy)\|$. Consequently the second integral is $O(e^{\sigma t})$. Adding $k(t)=O(e^{\rho t})$ gives $v(t)=e^{\mu t}D_*+O(e^{\sigma t})$. It is the subtraction of $k$ that gives the second power of frequency decay.

The newborn first-generation equation is $JP_tf=JR_tf+(v*JR_\cdot f)(t)$; for any starting type, $P_tf=R_tf+\int_0^t R_uB_0JP_{t-u}f\,du$. In the newborn equation, the leading convolution is $e^{\mu t}D_*\int_0^t e^{-\mu u}JR_u f\,du$. Extending to infinity gives $e^{\mu t}g d(f)$, with $d(f)=\ell JG_\mu f/(-\ell K'_\mu g)$. The discarded tail costs at most a constant times $\|f\|_W e^{\mu t}\int_t^\infty e^{-(\mu-\rho)u}du=O(\|f\|_W e^{\rho t})$. The remainder convolution is bounded by a constant times $\|f\|_W e^{\sigma t}\int_0^t e^{-(\sigma-\rho)u}du$. Inserting this newborn expansion into the second equation gives leading term $e^{\mu t}d(f)\int_0^t e^{-\mu u}R_u B_0g(S)\,du$. The same two integrals bound its tail and error, now with the extra factor $W(S)$ from \eqref{eq:core-weights}. Absorb the residual term $R_t f$ in the error. This yields, with $H=U_\mu g$,
\[
 |P_tf(S)-e^{\mu t}H(S)d(f)|\lesssim e^{\sigma t}W(S)\|f\|_W.
\]
Each exact-degree indicator has strictly positive $d$ since an original newborn can remain alive and gain leaves alone until it has that degree, then hold for an interval. Also $d_k\downarrow0$ by dominated convergence under residual occupation, bounded by the discounted $W$-mean. This proves all coefficient assertions.

Apply \eqref{eq:core-fringe-identity} and set $\Psi(T)=\sum_{v\ne o}H(S_v(0))+\int_0^\infty e^{-\mu u}I_{T,u}H\,du$, with the sum over old nonfounder cores. For each core $x$, the contribution to the weight of its ancestor's fringe at distance $j$ is at most $(2/d)(10/13)^j$: $h(c_x)\le2$, the denominator is at least $d$, and every path factor is at most $10/13$. There is at most one ancestor per distance, so $\sum_{v\ne o}W(S_v(0))\le (2/d)C(T)/(1-10/13)$. Summing the individual remainders therefore costs $O(\|f\|_\infty C(T)e^{\sigma t})$, using $W\ge1$. The immigration bound in \eqref{eq:core-weights} is $I_{T,u}W=O(C(T)e^{qu})$. The immigration error is bounded by a constant times $\|f\|_\infty C(T)e^{\sigma t}\int_0^t e^{-(\sigma-q)u}du$, while the discarded leading tail uses $e^{\mu t}\int_t^\infty e^{-(\mu-q)u}du=O(e^{qt})$. Thus neither bound depends on the seed shape. The founder toll is bounded and absorbed in the error. This proves the mean formula and upper comparison.

For the lower comparison, call a core eligible if $c\le2$ with a leaf, or if it is bare of degree two and has a bare degree-two neighbour. For $C\ge2$, let $N_1$ be the number of pendants of the core tree. The number of $c\ge3$ is at most $N_1-2$. Bare vertices with no bare degree-two neighbour use two distinct edges per such vertex, so there are at most $(C-1)/2$ of them. All pendants are eligible. These bounds imply at least $(C+5)/4$ eligible cores: the ineligible high-core-degree count is at most the eligible count minus two. For a nonfounder of the first kind, the original residual particle surviving one unit and its first root ring choosing a leaf yields a uniform positive lower bound on $H$, by compensation of the direct-birth toll: the choice probability is at least $1/3$ and the toll is at least $\min g$. For a bare nonfounder with a bare degree-two child, silence the root and the child's other neighbour for one unit, let that child's first ring select the root, and require survival. The child's degree is two until its choice. This gives a leaf to the root and the previous bound in the next unit. Any other bare eligible nonfounder with no such child can be charged to its bare parent with such a child, at most twice per parent. Allowing for the founder leaves at least $C/12-1$ good nonfounders. Thus $\Psi$ dominates a multiple of $C$ for large $C$. For bounded $C\ge2$, at least one pendant of the core tree is a nonfounder with a leaf, giving a uniform positive bound. When $C=1$, the founder's promotion rate is one until its first promotion, so direct immigration gives $\Psi\ge(\min g)/(1+\mu)$ (use $H(J_a)=g(a)$). This supplies a positive uniform $a$ in \eqref{eq:core-unif}. Source: \proofsrc{round2/proofs/astra/source/round2/standalone/claim-6.md}{all-seed contour and weight estimates, Sections 2--4}.
\end{proof}

\begin{lemma}[Construction and tightness of the spatial law]\label{lem:core-spatial}
On the marked-core space of Theorem~\ref{th:core-spatial}, the cell operators are well-defined on $C_b$, satisfy $Q_tQ_u=Q_{t+u}$ and $1\le Q_t1\le e^t$, and admit a probability eigenmeasure at $e^{\mu t}$. This measure is a weak subsequential limit of the mean tilts below and has $\pi(e^{\theta c_o})<\infty$ for $0<\theta<\log(1+\mu)$.
\end{lemma}
\begin{proof}
On each oriented old core edge $x\to y$ use a Poisson proposal at rate $1/c_x^0$ and accept with probability $c_x^0/D_x(t-)$, where $c_x^0$ is the initial core degree. These are admissible since total degree does not decrease. All other transitions are within one cell, driven by independent drivers: leaf-source gains at rate $L_x$, internal core sends at $1/D_x$ per edge and promotions at $L_x/D_x$. The core-source choices can be realized by disjoint subintervals of a rate-one clock; leaf-source gains can also use individual leaf clocks. In a finite union of cells there are finitely many incoming old-edge proposals in bounded time, and endogenous births have total rate at most the current ordinary population there. Hence finite old-ball truncations with missing slots retained and inbound proposals across the cut suppressed are nonexplosive. The root cell grows at rate at most its core count, so its first and second moments of count are at most $e^t$ and $2e^{2t}$, respectively. Relabelling at intermediate time partitions future descendants, giving the semigroup by conditional expectation; independence of whole cells is not required.

Couple all truncations by their common drivers. For a score of ball radius $k$, only cells based in the old $B_k(o)$ can matter for any member of the root cell. A discrepancy from beyond radius $R\ge k$ must reach one of these cells along a time-ordered chain of proposals; erase any loops and count only the last $j=R-k+1$ edges. A simple path of length $j$ has expected number of ordered tuples $t^j/j!$ times the product of reciprocal old sender core degrees. Reversing paths identifies the product with the reversed simple random walk probability times $c_w^0/c_{\rm start}^0\le c_w^0$, where $w$ is the endpoint. The sum over reversed simple paths is at most one in probability. Thus disagreement through time $t$ has probability at most $p=\min(1,(t^j/j!)\sum_{w\in B_k(o)}c_w^0)$; for a singleton there are no paths. The event from the old proposals controls all larger truncations simultaneously. Its probabilities are summable in $R$ for fixed $t,k,G,o$, giving almost-sure stabilization on every fixed old ball and time interval, hence a Markov evolution with independent future Poisson increments. The construction also gives the conditional continuation at a promotion: use the stopped Poisson increments and then the same finite-ball localization. By Cauchy--Schwarz the cylinder score error is at most $2\sqrt{2}e^t\|f\|_\infty\sqrt p$. The old ball of radius $R+1$ determines the missing slots of the $R$-ball. For a cylinder this proves continuity at every initial tree. For a bounded continuous score and a converging sequence of initial trees, take the same drivers on matching balls. For each fixed output ball radius the probability of agreement tends to one by choosing increasing matching old radii. Equivalently on any subsequence one can take a further subsequence with summable disagreement probabilities at successively larger radii, giving almost-sure local convergence of every rooted output in the matched root cell. The cell count is uniformly integrable by the second-moment bound, so the expected sums converge as well. This proves $C_b$ invariance.

Let $m_t(f)=\mathbb E_{T_3}\sum_{x\,\mathrm{core}} f(G_t,x)$, so $m_t(1)=F_3(t)$ and $m_{t+u}=m_tQ_u$. For $0<\delta<\mu$ and $T>0$ set $w(t)=\exp(-\mu t-\delta|t-T|)$, $Z=\int_0^\infty w(t)F_3(t)dt$, and $\pi_{\delta,T}=Z^{-1}\int w(t)m_tdt$. The exponent gives finiteness of $Z$ and for fixed $\delta$ gives $w(0)/Z\to0$ as $T\to\infty$ (integrate on $[T,T+1]$ and use monotonicity). Choose $\delta\downarrow0$ and then $T$ large for each $\delta$ so this ratio goes to zero. For fixed $u$, changing variables in the tilt, $w(t-u)/w(t)$ is between $e^{(\mu\pm\delta)u}$ for $t\ge u$, and the missing mass on $[0,u]$ is $O_u(w(0)/Z)$. This gives $\|\pi_{\delta,T}Q_u-e^{\mu u}\pi_{\delta,T}\|_{\rm TV}\to0$. Every bounded time interval has vanishing tilt mass by the same estimate.

Conditional on the full history through a core's promotion (or on the initial tree for the centre), its mean degree $a$ time units later is at most $2e^a$ by degree-rate domination. Its core degree is dominated by $1+\operatorname{Pois}(a)$ using the fresh increments of its own clock. Both estimates are conditional on the past, not on any subsequent environment. Thus each of $D^\alpha$ and $e^{\theta c}$ has conditional mean at most $B e^{\kappa a}$ for $(B,\kappa)=(2^\alpha,\alpha)$ and $(e^\theta,e^\theta-1)$, with $0<\alpha<\min(1,\mu)$ and $0<\theta<\log(1+\mu)$. The total promotion intensity is at most the core count. Compensate births and use $w(t+a)\le w(t)e^{-(\mu-\delta)a}$; the tilt of each moment is bounded by $B(1+w(0)/Z)/(\mu-\delta-\kappa)$ for $\delta<\mu-\kappa$. The term with $w(0)$ is the initial centre. Besides tight root degree this gives uniform smallness of $\pi_{\delta,T}(c_o\mathbf1_{D_o>J})$ for $J$ large by Cauchy--Schwarz. Exchange the two endpoints in the finite root sums: the mean number of simple paths of length $i$ from a root with successive bounds $J_0,\ldots,J_{i-1}$ to a vertex with $D>J_i$ is at most $(J_1\cdots J_{i-1})\pi_{\delta,T}(c_o\mathbf1_{D_o>J_i})$, by counting from the bad endpoint. Choose cutoffs successively to have summable, arbitrarily small errors. At each radius there are finitely many marked balls with those bounds. Consistent ball records give a Polish local space (boundary degree constraints are checked at the next radius), and the simultaneous bounds form a compact set. Prokhorov gives a subsequential probability limit $\pi$. Bounded cylinders are determining; continuity of their images and the tilt relation give $\pi Q_u=e^{\mu u}\pi$ for each fixed $u$ on the same weak subsequence. Truncating the exponential root moment gives the asserted finite moment. Source: \proofsrc{round2/proofs/astra/source/round2/operator.md}{spatial cutoff and mean tilts, Sections 1--2}.
\end{proof}

\sectag{DEN}
\section{Absolute continuity of the amplitude}\label{sec:core-den}
\tagnext{DEN}
\begin{theorem}[Absolute continuity of the core amplitude]\label{th:core-density}
For the edge start, the law of $X$ in Theorem~\ref{th:core-amplitude} is absolutely continuous with respect to Lebesgue measure. The proof uses the founder hub theorem and the hub-children property specified below. No density regularity or lower bound is asserted here.
\end{theorem}
\begin{proof}
Two facts supply a pair of adjacent hubs. First, almost surely $Z_1+Z_2>0$, where $Z_v=\lim_n D_v(n)/n$ (Addario-Berry et al.~\cite[Theorem~3.2 at the founder edge, arXiv v1; see Theorem~3.1 for the degree limits]{AB}). Second, on the event $Z_v>0$, each fixed vertex $v$ promotes infinitely many leaves that themselves have positive limiting shares. For the second fact we follow Sections~1 and~3 of \proofsrc{round2/proofs/astra/source/round2/fringe/hub-children.md}{hub children}. A \emph{silent planted newborn} is a core with one leaf and one exterior slot that counts in its degree, sends nothing and whose incoming arrows are null; all other rates are physical. Its leaf stock gains at rate at least the stock and loses at total rate at most one, and each tracked leaf is lost at rate at most $1/2$. So the stock dominates a linear birth--death process with rates $1,1/2$, which survives with probability $1/2$. On survival the stock hits every level $k$, and at such a hit the newborn's mass $W^0=\lim e^{-t}D(t)$ is positive with conditional probability at least $(k-1)^2/[k(k+1)]$, by the comparison in the proof of Theorem~\ref{th:hub-order}. Hence $q=\Prob(W^0>0)\ge1/2$. Now fix $v$ and start at the first time $\tau$ at which the tree has at least three vertices and $v$ is a core vertex. This is a stopping time, finite on $\{Z_v>0\}$. For the founder that becomes the centre at the first birth, this means starting at $T_3$; for the other founder, wait for its promotion. On $\{\tau<\infty\}$, conditional on the history through $\tau$, the future ring times $\tau+t_j$ of $v$ form a rate-one Poisson process. To every ring, before knowing what it chooses, attach an independent silent planted newborn with mass $W^0_j$. When that ring attaches a new leaf to an old leaf of $v$, start the copy at the promoted old leaf, with the new leaf as its one leaf and $v$ as its exterior core neighbour. Arrows from the child into $v$ are generated on a separate channel of rate $1/D_{\rm child}$ and leave the child unchanged. The child leaves its copy only if $v$ later sends it a leaf, and then continues with fresh drivers by the strong Markov property. This construction gives the physical law. It is made before conditioning on hubhood. On $\{Z_v>0\}$ we have $c_v(t)\le c_v(\tau)+\#\{j:t_j\le t-\tau\}=O(1+t)$ and $e^{-t}D_v(t)\to YZ_v>0$, where $Y=\lim e^{-t}|T(t)|$. So the rate $c_v/D_v$ of sends into existing core neighbours has finite integral, and localization at bounded compensator leaves only finitely many such sends. Hence all but finitely many rings of $v$ promote a child that follows its copy forever, and that child has $Z_w=e^{-(\tau+t_j)}W^0_j/Y$. The marks $W^0_j$ are i.i.d., independent of the clock, and positive with probability $q>0$. So almost surely infinitely many children of a hub are hubs. By the Yule clock, along a hub--hub edge the two integrated directional rates $\int dt/D_v(t)$ are finite: past any fixed time they decay exponentially eventually, and before then the rates are bounded. Localization at bounded compensator implies finitely many crossings almost surely. Thus the events of a given finite labelled seed at size $n$ and a given core--core edge never crossed after that size form a countable full-probability cover. No random choice of a last crossing is used as a stopping time.

Fix one such seed and edge with positive event probability. Restart elapsed time at the size-$n$ epoch, conditioning first only on the finite labelled seed there. By the strong Markov property the future has the law of a fresh process from this seed. The edge is fixed, not selected using future crossings. Run the two cut components independently with the exterior slots included in the degrees. An attempted arrow out of a slot kills that side; until any kill the superposition has exactly the physical rates. Each side starts at ordinary size $m_i\ge2$. At current size $m$, internal rate $m-1/D_{\rm root}$ plus kill rate $1/D_{\rm root}$ equals $m$. Preassign independent $\operatorname{Exp}(m)$ waits and independent arrow-choice uniforms on each side, selecting each arrow with probability its rate divided by $m$. Survival forever depends only on the uniforms: the statement for finitely many steps extends by decreasing intersection, and independence then holds for the full sequences. Conditional on fixed-cut survival, the sides remain independent, their size clocks have their original laws, and the clocks remain independent of the surviving marked chains. This is not the unconditioned law of the process with the slot choice simply deleted.

For clarity we derive the side limits under survival without an inert-tree limit input. Counts add exactly with slots restored. By Theorem~\ref{th:core-amplitude} and the Yule clock the total normalized count $e^{-\mu t}(C_A(t)+C_B(t))$ has a positive finite limit. Freeze all randomness except the first wait on side $A$. Fubini gives convergence of the total for a full set of first waits. After a first wait $s$ its count is a common function $F_A(t-s)$. Take two admissible waits $s_1<s_2$, put $d=s_2-s_1$, and subtract the totals. For $H(t)=e^{-\mu t}F_A(t)$ this gives $H(t)-e^{-\mu d}H(t-d)\to\ell$ for some finite $\ell$. Nonnegativity bounds $H$ eventually by either convergent total after a time shift. Subtract $\ell/(1-e^{-\mu d})$ and take the limsup of the absolute value; it is bounded by itself times $e^{-\mu d}<1$, hence is zero. The same argument on $B$ gives both time limits. Dividing by the Yule masses to power $\mu$ gives the nonnegative finite limits $X_A,X_B$ along local size epochs, each a function only of its surviving marked chain.

In fact each side limit is positive under survival. Glue two copies of this one-slot seed into an ordinary finite tree, which is reachable from an edge by a finite sequence of positive-probability leaf additions. Each copy has positive survival probability, as already follows from the fixed event under consideration. Thus the doubled tree has a positive-probability no-cross event with iid surviving side limits. A positive probability of a zero side limit would make both vanish with positive probability, contradicting positivity for the reachable free seed. Positivity from any reachable seed follows by conditioning the original almost-sure limit on that seed.

The independent masses are $W_A\sim\operatorname{Gamma}(m_A,1)$ and $W_B\sim\operatorname{Gamma}(m_B,1)$, independent of $(X_A,X_B)$. The change of variables to sum and proportion yields $U=W_A/(W_A+W_B)\sim\operatorname{Beta}(m_A,m_B)$ independent of the side coefficients. Thus in the law conditional on this specified seed and edge and on no subsequent crossing of the edge,
\[
 X=U^\mu X_A+(1-U)^\mu X_B.
\]
For fixed positive coefficients and $0<\mu<1$ the map of $U$ is strictly concave with one interior critical point, and has a locally Lipschitz inverse on each open monotone branch. Exhaust the branches by compact intervals, ignoring the critical point of zero Beta mass. Every Lebesgue-null set has a null preimage in Beta probability. Integrate over the coefficients and sum the null masses over the countable cover (overlap is harmless). Since $X>0$ a.s., this proves absolute continuity on the whole line. Sources: \proofsrc{round2/proofs/neon2/limit-density/README.md}{one-edge Beta split} and \proofsrc{round2/proofs/neon2/limit-density/REPORT.md}{clock factorization and side-limit argument}.
\end{proof}

\sectag{MASS}
\section{Unit hub mass}\label{sec:conservation}
Write $\mathcal F_n$ for the history through size $n$, $W_n$ for the attachment recipient at the next step, $D_v(n)$ for degree, and $C_n$ for the number of cores (vertices of degree at least two). Put $D_v(n)=0$ before birth. Set
\[
 P_n(v)=\Prob(W_n=v\mid\mathcal F_n)=\frac1n\sum_{u\sim v}\frac1{D_u(n)}\le \frac{D_v(n)}n
\]
for $v\le n$, and set $P_n(v)=0$ otherwise. This is also the conditional law given $T_n$. After birth $D_v(n)/n$ is a nonnegative supermartingale: its next conditional mean is $(D_v(n)+P_n(v))/(n+1)$. Consequently the limits $Z_v$ exist simultaneously almost surely. Call $v$ a hub if $Z_v>0$; no degree-divergence characterization is needed here.
\tagnext{MASS}
\begin{theorem}[Unit hub mass and attachment limits]\label{th:hub-mass}
On the event $\sum_{n\ge3} C_n/n^2<\infty$, almost surely
\[
 \sum_{v\ge1} Z_v=1,\qquad
 \|P_n-Z\|_{\mathrm{TV}}\longrightarrow0,\qquad
 \mathcal L(D_{W_n}(n)/n\mid T_n)\Longrightarrow \sum_{v\ge1}Z_v\delta_{Z_v}.
\]
In the edge-seeded friend tree that event has probability one. In particular the conditional recipient's mean normalized degree tends to $\sum_v Z_v^2$, and the unconditional mean tends to $\E\sum_v Z_v^2$.
\end{theorem}
\begin{proof}
This combines the \proofsrc{round2/proofs/astra/source/round2/standalone/claim-7.md}{hub-mass argument} with the \proofsrc{round2/proofs/neon2/claim14-proofs/README.md}{pathwise summability criterion}. For $n\ge3$ the core induces a tree: every intermediate vertex on a path between cores is a core. Put $q_n(v)=L_v(n)/n$ on cores, where $L_v(n)$ counts leaf neighbours, and zero otherwise. Every leaf has a unique core neighbour, so $\sum_v q_n(v)=1-C_n/n$. The vector $P_n-q_n$ is nonnegative with mass $C_n/n$ (the contribution of core initiators). In contrast, $\sum_v D_v(n)/n=2(n-1)/n$: degrees count both ends of each edge. The law of a uniformly sampled edge end would weight $v$ by $D_v(n)/[2(n-1)]$, which is not the recipient law $P_n$. The theorem sums the fixed-label limits, not the finite-size degree totals. On the stated event, $C_n/n\to0$ because $\sum_{r=n}^{2n} C_r/[r(r+1)]\ge C_n(1/n-1/(2n+1))$. Thus $q_n(v)\to Z_v$ for each fixed $v$: a core loses at most $C_n/n$ by counting only its leaf neighbours, and a leaf has normalized degree $1/n$. In particular $\sum_v Z_v\le1$ there. We must still rule out escape of mass to late labels.

For integer $K\ge1$ stop at $\tau=\inf\{r\ge3:\sum_{j=3}^r C_j/[j(j+1)]>K\}$, with $\inf\varnothing=\infty$, and write $B_r=C_{r\wedge\tau}$. Before the crossing the sum for $B$ is at most $K$; its frozen tail from a finite crossing has sum $C_\tau/\tau\le1$. Hence $\sum_{r\ge3} B_r/[r(r+1)]\le K+1$ on every path. For $m\ge3$ start a sum $Y_m=0$ and at step $r>m$, if $r\le\tau$, update the normalized degrees of all already counted vertices and add $2/r$ for the new promotion if there is one. Freeze the entire sum, including the denominator, after the crossing step. The active indicator $\mathbf1_{\{r\le\tau\}}$ is $\mathcal F_{r-1}$-measurable, so the previously counted contribution has nonpositive conditional drift. The new contribution is exactly $2(B_r-B_{r-1})/r$. Summation by parts, followed by $B_n/n\le\sum_{r\ge n}B_r/[r(r+1)]$, gives
\[
 \E Y_n\le 2\E\sum_{r=m+1}^n\frac{B_r-B_{r-1}}r
 =2\E\left[\frac{B_n}n-\frac{B_m}{m+1}+\sum_{r=m+1}^{n-1}\frac{B_r}{r(r+1)}\right]
 \le 2\E\sum_{r>m}\frac{B_r}{r(r+1)}.
\]
On $\{\tau=\infty\}$, $Y_n\ge\sum_{v>m}q_n(v)$, since every such core was promoted after $m$. The latter expression tends to $1-\sum_{v\le m}Z_v$. Fatou's lemma bounds its expectation restricted to $\{\tau=\infty\}$ by the last display. That bound tends to zero by dominated convergence as $m\to\infty$, giving unit mass on $\{\tau=\infty\}$. The union over $K$ contains the summability event. For the edge seed, the sublinear mean bound $\E C_n=O(n^\beta)$ for some $\beta<1$ follows from Theorem~\ref{th:exist}, so Tonelli gives summability almost surely. Alternatively Addario-Berry et al.\ \cite[Prop.~6.22]{AB} prove $n^{-\beta}C_n\to0$ almost surely for $\sqrt{2}-1/2<\beta<1$. Any such $\beta$ gives summability by comparison with $\sum_n n^{\beta-2}$. Thus their earlier bound suffices for this theorem without the new exponent or either enclosure.

Coordinatewise convergence and convergence of the total masses imply $\|q_n-Z\|_1\to0$: split at a finite label $m$, then let the limiting tail beyond $m$ decrease to zero. Since $\|P_n-q_n\|_1=C_n/n$, the label TV statement follows. For bounded continuous $f$ on $[0,1]$, the discrepancy of the two integrals is bounded by
\[
 \|f\|_\infty\|P_n-Z\|_1+\sum_v Z_v|f(D_v(n)/n)-f(Z_v)|\longrightarrow0
\]
by dominated convergence against the fixed probability vector $Z$ on the path. Apply this with $f(x)=x$ and then use bounded convergence for the unconditional assertion. Also, the fraction of all vertices that are leaves adjacent to hubs equals $\sum_{v:Z_v>0}q_n(v)$ for $n\ge3$, and tends to one by the same $\ell^1$ convergence. These are conditional laws, not a pathwise convergence assertion for the actual recipients' degrees. No necessity of the summability hypothesis is asserted.
\end{proof}

\sectag{SPA}
\section{A conditional spectral characterization}\label{sec:spa}
This section proves a spatial eigenmeasure statement, which is unconditional, and a spectral-radius identity, which takes two finite-exterior estimates as hypotheses. Neither the hypotheses nor the identity is used anywhere else; Theorem~\ref{th:core-means} uses only the unconditional Lemma~\ref{lem:core-spatial}.
\tagnext{SPA}
\begin{theorem}[Spatial spectral characterization, conditional]\label{th:core-spatial}
Let $\mathcal G$ be the space of finite or countably infinite connected locally finite rooted core trees marked by $L_x\in\mathbb Z_{\ge0}$, with $D_x=L_x+c_x\ge2$, in the topology of agreement on marked balls. A singleton is allowed. The positive bounded operators $Q_t$ on $C_b(\mathcal G)$ count in mean the future cell of the old root, assigning each promotion to its promoter's cell. There is a Borel probability $\pi$ with $\pi Q_t=e^{\mu t}\pi$ for all $t\ge0$; this part is unconditional. Under the finite-exterior planting and one-edge hypotheses of Lemma~\ref{lem:core-ports}, also $r_{C_b(\mathcal G)}(Q_1)=e^\mu$, and any $r>0$ with $\nu Q_1=r\nu$ for a Borel probability $\nu$ equals $e^\mu$. Uniqueness of $\pi$ is not asserted.
\end{theorem}
\begin{proof}
Construction and eigenmeasure existence are in Lemma~\ref{lem:core-spatial}; the conditional spectral-radius upper bound and the matching lower bound on any probability eigenvalue are in Lemma~\ref{lem:core-ports}. The cutoff estimate used for construction is stated there with its pointwise hypotheses; it is not a separate uniform operator-norm truncation claim. Source: \proofsrc{round2/proofs/astra/source/round2/operator.md}{cell operator and one-edge comparisons}.
\end{proof}
\begin{lemma}[Finite-exterior hypotheses and spectral comparison]\label{lem:core-ports}
Use the following two finite-tree inputs, always with the physical degree including the edge across the cut.
\begin{enumerate}
\item For every integer $M\ge2$, $0<\epsilon<1$ and fixed integer $k\ge6M/\epsilon$, for each sufficiently large elapsed $s$ there is a rooted $k$-vertex branch $B_s$ such that when joined by one edge to any finite exterior parent of degree at least $M$, its conditional mean core count after $s$ is at least
\[
 k^{-1}e^{-\eta s}\big(F(s)-k-\theta_k(s)\big),\quad
 \eta=7\log 2/\log(1/\epsilon)+9/M,\quad
 \theta_k(s)=\sqrt{k(k+1)(2k+1)}e^{-s}.
\]
The rooted shape is chosen before planting, independently of exterior, from the finite shapes of that size. This is the weighted one-edge comparison (Lemma 4.1(c), Corollary 5.3 and Lemma 6.1 of \proofsrc{proofs/existence.md}{the finite-seed proof}).
\item For every integer $j\ge4$ there are nondecreasing $A_j(t)>0$ with $\log A_j(t)=O_j(\log(2+t))$ and $0\le\Lambda_j\le6/j+2/j^2$, such that the mean core count on the newborn side of a core--core edge, started just after promotion with one newborn core and one leaf on that side, is at most $A_j(t)e^{\Lambda_jt}F(t+1)$. Uniformity is over the entire finite exterior at the birth, not over arbitrary newborn-side shapes. This is the input of \proofsrc{round2/proofs/astra/source/round2/finite/one-port-upper.md}{one-edge upper estimate, equations (1)--(2)}.
\end{enumerate}
Under these inputs, $r_{C_b(\mathcal G)}(Q_1)=e^\mu$, and no other positive eigenvalue of $Q_1$ admits a Borel probability eigenmeasure.
\end{lemma}
\begin{proof}
Both estimates extend to a locally finite exterior for fixed elapsed time. Complete an increasing finite old ball by appending a core with an ordinary leaf to each missing core-neighbour slot; this preserves the degrees on the ball, includes the prescribed branch and yields genuine finite marked trees. The cutoff coupling gives convergence of the expected branch count: all future branch cores lie in cells of its finitely many old branch cores, whose counts are uniformly integrable by Yule domination. The original parent is already core so the cut never transmits a cell label. For a finite seed branch with old leaves, any such leaf promoted later still joins an old branch core cell. The same argument after conditioning at the end of a finite planting script transfers the lower bound.

For any initial $(G,o)$, the root's cell can plant $B_s$ with a fixed positive probability in a fixed elapsed time, uniformly over the finitely many possible shapes at the chosen size. Here are the planting details (compare Lemma 3.1 of the finite-seed proof). If $o$ has no leaf, choose a neighbour, silence all its neighbours including $o$ for one unit, and require its first ring in that unit to select $o$. The neighbour's degree is unchanged until then, so the event has positive probability. If $o$ has a leaf just silence $o$ for the unit. In a second unit silence $o$ and require at least $M$ rings of a tagged leaf, yielding at least $M+1$ leaves. In a third unit require the first ring of $o$ to promote one of them, with probability of a leaf choice at least $(M+1)/(M+1+c_o)$, and silence $o$, that core and its newborn after the promotion until unit end. The new branch is an ordinary edge attached to $o$. In the next $k-2$ units, script the remaining additions in connected order, one per unit. For each addition choose an inside neighbour of the attachment target to send to it at its only ring; silence all other inside vertices, the newborn after birth and $o$. Before its ring the sending degree is fixed by these silences. All named clock conditions have strictly positive probability, bounded below at fixed $G,o,M,k$ independently of the chosen branch; exterior activity cannot reach the branch while $o$ is silent. All new cores belong to the root cell. Consequently $Q_{s+h}1$ is bounded below by a positive constant times the bound in input 1. Taking $s\to\infty$ with parameters fixed gives $\liminf t^{-1}\log Q_t1(G,o)\ge\mu-\eta$. Then send $M\to\infty$ and $\epsilon\downarrow0$ to get the lower exponent $\mu$. If $\nu Q_1=r\nu$, Fatou applied to $e^{-(\mu-a)n}Q_n1$ for $a>0$ gives $r\ge e^\mu$.

Conversely all descendants of the old root other than itself lie in branches of its direct promotions, whose mean counts are bounded by input 2 at birth. Compensate those promotions at rate at most one. By monotonicity of $F,A_j$ and the nonnegative $\Lambda_j$, uniformly in rooted initial states,
\[
 Q_t1\le1+t A_j(t)e^{\Lambda_jt}F(t+1).
\]
One can also first prove this for finite completions and pass to the limit by the cutoff estimate. Since $\|Q_1^n\|=\|Q_n1\|_\infty$, take $n\to\infty$ then $j\to\infty$ in the spectral-radius formula. This gives $r(Q_1)\le e^\mu$, and the eigenmeasure of Lemma~\ref{lem:core-spatial} gives the reverse. Any probability eigenvalue is bounded above by the same norm estimate and below by the preceding Fatou argument. The finite-exterior estimates remain hypotheses of this spectral comparison.
\end{proof}

\sectag{DRF}
\section{Finite local drift tables}\label{sec:drf}
This section shows that finite rational drift tests of a specific phase-averaged kind can certify exponents arbitrarily close to $\mu$ on both sides. It is not used elsewhere, and it gives no runtime bound.
\tagnext{DRF}
\begin{theorem}[Finite local drift completeness]\label{th:core-drift}
There is a family of finite rational linear drift tests on phase-averaged cut tables whose certified upper exponents have infimum $\mu$ and certified lower exponents have supremum $\mu$. Specifically for $h\to\infty$ and $M=R=\lceil e^{3h}\rceil$, the tests certify rational exponents on either side of $\mu$ at distance
\[
 O\big(e^{-(\mu-\sigma)h}+(M^{-1}+R^{-1})e^{(2-\mu)h}\big),
\]
where $\sigma<\mu$ is from Lemma~\ref{lem:core-uniform}. More generally the same error bound applies for $M,R$ making its right side sufficiently small. This concerns the tables constructed here, not completeness of a previously used descriptor family.
\end{theorem}

Long previously proved a terminating two-sided rational approximation theorem for $\mu$ \cite{LongKernel}. The result here constructs a distinct family of phase-averaged local drift tables and gives a rate in their cutoff parameters, not a runtime bound for that earlier search.
\begin{proof}
Orient the physical core from its founder and assign depth to each core at promotion. For each phase $r$ modulo $R$, cut a low--low edge if the child's depth equals $r$ modulo $R$; low means degree less than $M$. Remove edges incident to high vertices from the low forest but retain their slots. Low components have height at most $R-1$ and degree bounded by $M-1$. Record the full degrees, leaves, slots, parent direction and relative depth residues. Each high vertex is a separate particle recording its depth residue. In the auxiliary process, lows gain leaves from their own leaves and active neighbours only; a gain reaching $M$ splits the component and replaces that vertex by a high particle. Low promotions have their true rate and the new edge is cut by the same phase rule. A high particle produces a low singleton, with one leaf, one inactive parent slot and the next depth, at rate one. After a split the resulting particles evolve independently. This finite rational generator on first-moment scores is denoted $G$ (it acts on column scores). Write $c_i$ for the core count of a particle and $\Phi_w$ for the phase-averaged sum of particle scores $w_i$ on the physical tree.

Let $A(w)$ be the maximum absolute score change from adding one leaf anywhere in a low particle, including splitting on conversion to high, and $B(w)$ the maximum absolute newborn score at a high birth. At high sources the missing promotion rate and the total outgoing core-arrow rate are each bounded in sum by $\sum_{\rm high}c_x/D_x\le2C/M$. Each low--low edge is blocked in one phase, with sum of the two directional rates at most one. All other events project exactly. Consequently, for the physical generator $\mathcal L$,
\[
 |\mathcal L\Phi_w-\Phi_{Gw}|\le D(w) C,\qquad D(w)=2(A(w)+B(w))/M+A(w)/R.
\]
Sufficient tests are $\Phi_{Gw-\lambda w}+D(w)C\le0$ for an upper exponent and $\Phi_{Gw-\lambda w}-D(w)C\ge0$ for a lower one, on every finite rooted core tree, with $w>0$. Positivity and finiteness give $\Phi_w\asymp C$. Since a discrete step has drift $\mathcal L\Phi_w/n$, multiplying the corresponding one-step factors for sufficiently large $n$ bounds the mean exponent on the stated side.

To make the tests finite, label each vertex by its high flag or low $(D,c)$, founder flag, parent high flag or absence and downward low neighbourhood through $R$ levels, stopping each branch at a high vertex. A high vertex need not list its children. Parent--child compatibility and all low arities are finite. Each vertex contributes its phase-averaged component-root score: for a low with a low parent only its cut phase counts; a low with high parent or no parent counts in all phases. These contributions use only relative depths. Convert the upper test (or the negative of the lower test) to an additive reward, including the term $D(w)$ per vertex. The grammar of legal trees has finitely many symbols and rules: bounded lists for lows; for a high vertex, an arbitrary finite list of children reporting a high parent, encoded by empty-list and head-plus-tail rules of zero reward. Every such high arity is realizable by adding sufficiently many leaves. Include the start rule and remove unproductive or unreachable symbols. For each rule $i\to(j_1,\ldots,j_b)$ with reward $a$, impose $v_i\ge a+\sum_j v_{j}$ with multiplicity, and $v_{\rm start}\le0$. Induction proves sufficiency. Conversely if every start reward is nonpositive, each productive reachable symbol has finite supremum of its finite derivation rewards (insert it in a fixed start context). The suprema satisfy the inequalities by independent choice of child derivations. Absolute score-change bounds are pairs of linear inequalities in slack variables. Homogeneous rescaling permits $w_i\ge1$; at rational $\lambda$, feasible rational linear systems admit rational witnesses.

To prove completeness set $w=e^{hG}c$, temporarily real. The auxiliary per-core promotion rate is at most one, so $c_i\le w_i\le c_i e^h$ and $B(w)\le e^h$. Couple two auxiliary copies started with a single additional leaf at a low vertex, retaining anatomical neighbour slots as annotations even across splits. Start with the target pair dirty. Matched clean sources have identical choices; a send to a dirty recipient may be ignored in one copy without affecting any clean pair. The phases depend on permanent depths. The only possible new dirt from a dirty ring is a child or a recipient pair. Thus $A(w)\le2e^{2h}$ by the same compensation as above, even when the initial addition converts its target to high.

Next couple the physical tree and each auxiliary phase from the same projected seed, starting with no dirt. At a clean low source only a blocked low--low send into a clean low recipient can introduce dirt, at most two marks. Clean high recipients absorb actual leaf gains (their auxiliary degree is not tracked). At a clean high source synchronize physical promotions with auxiliary births; on a choice of a core neighbour charge at most five marks (source pair, unmatched auxiliary child and possibly recipient pair). In each comparison exceptional high rings have intensity at most $2C_{\mathrm{true}}/M$. Summed over phases, the expected clean seam intensity is at most $\mathbb E C_{\mathrm{true}}$, since the true marginal is the same and each physical low edge has one cut phase. Dirty sources add at most two marks per ring. With $C_0$ initial cores and $\mathbb E C_{\mathrm{true}}(v)\le C_0e^v$, integration gives average expected dirt bounded by $d_h C_0$, where $d_h=(10/M+2/R)(e^{2h}-e^h)$. Thus
\[
 |\Phi_w-\mathbb E_T C(T(h))|\le d_h C_0,\qquad
 |\Phi_{Gw}-\mathbb E_T\textstyle\sum_x b_x(T(h))|\le(d_h+2e^h/M)C_0.
\]
The second inequality uses $Gw=e^{hG}Gc$: clean low intensities agree, dirty intensities are bounded by their count, and clean high discrepancies sum to at most $2C_{\mathrm{true}}/M$. By Lemma~\ref{lem:core-uniform} the two physical means equal $e^{\mu h}\Psi(T)d_2$ and $e^{\mu h}\Psi(T)d(b)$, respectively, with errors $O(C_0 e^{\sigma h})$. Integrating $d\mathbb E C(T(h))/dh=\mathbb E\sum b_x(T(h))$ from the three-vertex seed, where $\Psi>0$, proves $d(b)=\mu d_2$. If the stated error is small, $\Phi_w\gtrsim C_0 e^{\mu h}$ and $|\Phi_{Gw-\mu w}|+D(w)C_0\lesssim C_0(e^{\sigma h}+(1/M+1/R)e^{2h})$. Choose rational trial exponents a sufficiently large multiple of this relative error above and below $\mu$ to get strict margins proportional to $C_0$. For fixed $M,R$ all coefficients and score-change bounds vary continuously with the finite table, and there are at most $C_0$ particles in a projection. Rational perturbations preserve the margins uniformly in the physical tree. This proves the claimed rate and completeness. Source: \proofsrc{round2/proofs/astra/source/round2/fringe/local-certificates.md}{phase-averaged cut tables}.
\end{proof}

\part{The birth-order module}

\sectag{BIRTH}
\section{Hub probabilities and positive moments}\label{sec:hubs}\label{sec:hub-proofs}
This part uses from Part~I only the existence of the limits $Z_v$ (Section~\ref{sec:conservation}).
In \cite[Sec.~7]{AB}, Question~3 asks first about hub probabilities and then about stochastic order, Question~4 asks about conservation of hub mass, and Question~5 asks about the forest left after deleting vertices of large degree. We answer Question~4 and both parts of Question~3. Hub probabilities tie at the two founders and decrease strictly from label two onward (Theorem~\ref{th:hub-order}), and so do the tails $\Prob(Z_v>x)$ at every threshold $0\le x<1$ (Theorem~\ref{th:stoch-order}). Addario-Berry et al.\ \cite[Sec.~3]{AB} already observed that the conjectured conditional limit law for the next recipient's normalized degree is equivalent to unit hub mass. Theorem~\ref{th:hub-mass} also gives total-variation convergence of the recipient's label law. The moment comparison of Theorem~\ref{th:hub-moment-order} is stronger than a mean comparison but does not by itself give the stochastic order; Theorem~\ref{th:stoch-order} is proved by a separate threshold-by-threshold argument. For Question~5 we give a conditional obstruction to total-variation approximation, not a description of all components or a sharp cutoff transition.
Let $H_v=\{Z_v>0\}$ and $h_v=\Prob(H_v)$. The founder symmetry below is equality in law, not equality on the same realization.

\tagnext{HUB}
\begin{theorem}[Birth order of hub probabilities]\label{th:hub-order}
On every history through size $n\ge2$, before either of the next two labels is placed,
\[
 \Prob(H_{n+1}\mid\mathcal F_n)-\Prob(H_{n+2}\mid\mathcal F_n)
 \ge \frac{371}{12500 n(n+1)}.
\]
The limiting shares $Z_1$ and $Z_2$ have the same marginal law, $h_2-h_3\ge127/2440$, and $h_1=h_2>h_3>h_4>\cdots$.
\end{theorem}

\tagnext{MOM}
\begin{theorem}[Birth order of every positive moment]\label{th:hub-moment-order}
For every real $p>0$ and every history through $n\ge2$,
\[
 \E[Z_{n+1}^p-Z_{n+2}^p\mid\mathcal F_n]
 \ge \frac9{5000}\frac{\Gamma(p+1)\Gamma(n)}{\Gamma(n+p+2)}.
\]
Also $\E[Z_2^p-Z_3^p]\ge 9/[5000(p+1)(p+2)]$. Thus the moments tie at the founders and strictly decrease thereafter.
\end{theorem}

\paragraph{The common two-placement calculation.}
In the old tree let $d_u$ be the degree, $r_u=\sum_{v\sim u}1/d_v$, $w_u=r_u/(d_u+1)$ and $s_u=w_u/d_u$. Use any equivariant continuation score on a marked vertex in a completed tree. Let $I(u)$ denote its middle-minus-tip difference when the second birth attaches to the first, whose old parent is $u$. Let $J(u,v)$ be the difference of the scores of two new leaves at distinct old vertices $u,v$ in that order. Then the conditional difference, times $n(n+1)$, is
\begin{equation}\label{eq:hub-signed}
 \mathcal N=\sum_u w_u I(u)+\sum_{\{u,v\}\in E(T_n)}(s_v-s_u)J(u,v).
\end{equation}
Indeed the second target being the first newborn has rate $1/(d_u+1)$. Repeated old targets give exchangeable leaves. For distinct old targets pair the two reverse orders on the same unlabelled completion. Nonadjacent pairs cancel. For an edge, the second target's rate loses $1/[d_u(d_u+1)]$ after a birth at $u$ (and the corresponding quantity at $v$ in the reverse order); subtracting gives exactly $s_v-s_u$. In particular it is not legitimate to compare arbitrary individual completions as if every summand had the same sign.

In continuous time every vertex has a rate-one clock with a uniform neighbour mark. The embedded chain is the prescribed tree; at each fixed size its holding time is exponential with rate equal to the size and independent of the jump choices. From a fixed completion of size $m$, the normalized population converges to $S\sim\mathrm{Gamma}(m,1)$ independently of the embedded chain. Indeed the population is Yule: from one it is geometric with parameter $e^{-t}$ on $\{1,2,\ldots\}$, and from $m$ it is a sum of $m$ independent copies. The normalized population is a martingale bounded in every fixed $L^q$, $q>1$, by these geometric laws, and the scaled geometric converges in law to an exponential. In these clocks $W=\lim_{t\to\infty} e^{-t}D(t)=SZ$. For the hub score use the probability of $W>0$; for powers use $m_p=\E W^p/\Gamma(1+p)$. With $(a)_p=\Gamma(a+p)/\Gamma(a)$, the latter numerator in \eqref{eq:hub-signed} is converted back to a conditional share difference by $\Gamma(1+p)/[n(n+1)(n+2)_p]$. The common completion size is essential. At a random return or promotion use the physical remaining law and multiply a moment by the clock discount $e^{-p t}$, not a size-factor substitution.

\begin{proof}[Proof of Theorem~\ref{th:hub-order}]
The exact sign certificate for \eqref{eq:hub-signed} at hubhood is detailed in Section~\ref{app:hub-certificates}: it gives $\mathcal N\ge b+2a=371/12500$ on every multi-core tree and at least the same floor on the edge and all stars. Averaging this inequality starts the marginal strict comparisons at $(3,4)$.

For the separate founder seam, at size three label 2 is the centre with probability $1/2$; otherwise it is an exchangeable leaf with label 3. The centre's two tracked leaf families dominate a linear birth--death process with per-particle rates $1,1/2$ (track only living leaves and thin deaths), so its hub probability is at least $3/4$. For use here and in Section~\ref{app:hub-certificates}, a linear birth--death process from one with rates $1,\lambda<1$ has extinction probability $\lambda$: the embedded walk hits zero with that probability, by solving the finite gambler's ruin recursion and sending its upper barrier to infinity. Survival of the tracked population implies hubhood even though this comparison has a slower growth exponent. In fact from $k$ physical leaves compare below by a chain with births $j$, total death rate one for $j>0$, absorption at zero. Its scaled stock is a nonnegative supermartingale, bounded in $L^2$ by the Yule comparison from $k$; its limit has mean at least $k-1$ by integrating the death drift, and second moment at most $k(k+1)$. Thus the conditional hub probability at a hit of $k$ leaves is at least $(k-1)^2/[k(k+1)]$ by Cauchy--Schwarz. On survival the slower linear process hits every $k$ (it exits each finite positive band almost surely), and the probability of a nonhub whose stock hits $k$ tends to zero by this bound and the strong Markov property.

For the exchangeable leaf, at its centre track its other sibling and subsequent gains from the tracked leaves and the marked leaf. Before the first send to the marked leaf, at stock $j$ the gain rate is at least $j+1$, losses at most $j/\max(2,j+1)$ and the send rate at most $1/\max(2,j+1)$. The test $f(j)=\min(1,2/(j+1))$ with payoff one at the send bounds the chance of ever sending. Indeed it has nonpositive drift at $j=0,1$ and at $j\ge2$ the upper drift is $-(j-2)/[(j+1)(j+2)]$. To sharpen the no-send bound from $j=1$, stop at stock four with payoff $1-f(4)=3/5$, or on sending with payoff zero. Set $g_i=i+1$, $b_i=i/\max(2,i+1)$, $k_i=1/\max(2,i+1)$. For the comparison rates, the probability of advancing from $i$ to $i+1$ before sending is $q_i=g_i/(g_i+b_i+k_i-b_iq_{i-1})$ (omit the last term at zero), by the first-step equation and the strong Markov property after a downward jump. The resulting test $(3/5)\prod_{i=j}^3q_i$ is increasing in $j$, so larger gains and smaller losses or killing in the actual process give a subsolution. The stopped process exits almost surely: inside the bounded stock range there is uniformly positive probability to reach four by successive gains in a unit interval, since the gain is bounded below and the competing rates are bounded above. Bounded stopping gives no-send probability at least $(3/5)q_1q_2q_3=108/305$. A leaf cannot be a hub without a send. Taking half the centre--leaf gap gives $127/2440$; exchanging the initial edge gives the founder tie in law.
\end{proof}

\begin{proof}[Proof of Theorem~\ref{th:hub-moment-order}]
Section~\ref{app:hub-certificates} gives the exact bounds for normalized moments in \eqref{eq:hub-signed}, uniformly in $p>0$, including the seed and stars. Multiplication by the common factor gives the stated gamma ratio since $n(n+1)\Gamma(n)=\Gamma(n+2)$.

For the founder seam, from size three the centre starts with two tracked leaves. The stock tests in Section~\ref{app:hub-certificates} give $\E W_{\rm centre}^p/\Gamma(1+p)\ge(2+p)/e$ for $0<p\le1$, and at least one for $p\ge1$. For $0<p\le1$ use the hub cap $a_0=197/305$ just proved and the exact mean cap $b_0=387/1000$ from the \proofsrc{round2/proofs/astra-late/source/round2/open/age-order/hub-ranking/tip-moments/check.py}{tip-moment checker} (the physical recurrence is given in Section~\ref{app:hub-certificates}). H\"older gives $\E W_{\rm leaf}^p\le a_0^{1-p}b_0^p$. If $E$ is mean-one exponential then $\E\log E\ge-\int_0^1(-\log x)\,dx=-1$, hence $\Gamma(1+p)\ge e^{-p}$ by Jensen. The log ratio of the raw centre lower $(2+p)e^{-1-p}$ to the leaf upper is concave in $p$. Both endpoint ratios exceed $51/50$, even with $e<11/4$. The raw gap is therefore at least $b_0/50$ since $a_0>b_0$. Using $\Gamma(1+p)\le1$ by log convexity on $[0,1]$, division by $2(3)_p$ yields a gap of at least $(b_0/50)/[(p+1)(p+2)]$.

For $p\ge1$ the leaf gain rate at degree $D$ is at most $D-1+1/2$ because its original parent keeps degree at least two. The affine pure-birth comparison from degree one gives normalized moment at most $(1/2)_p/\Gamma(1+p)$, the $p$-th moment of $\mathrm{Beta}(1/2,1/2)$, at most $1/2$. Thus the half-difference after common size-three normalization is at least $1/[2(p+1)(p+2)]$. Both ranges imply the claimed founder bound. Symmetry and averaging finish the proof.
\end{proof}

\paragraph{Telescoping, mixtures and scope.}
For $2\le m\le i-1$, $i<j$, condition each successive comparison on the common $\mathcal F_m$ and sum. The hub bound is $(371/12500)(1/(i-1)-1/(j-1))$. The moment bound is
\[
 \frac{9}{5000(p+1)}\{\mathrm B(i-1,p+1)-\mathrm B(j-1,p+1)\},
\]
by the beta recurrence. These bounds apply unconditionally also at $i=2$ using the stronger founder seam; the conditional versions apply at an a.s. finite stopping size with fixed positive label offsets by partitioning by that size. Integration gives strict expectation order for every nonzero finite positive mixture of powers, for example $z\mapsto1/(a-\log z)$ with $a>0$ and value zero at zero. Taking $p\downarrow0$ gives hubhood with a weaker constant than Theorem~\ref{th:hub-order}. Neither positive-moment order nor the zero-threshold comparison supplies all positive-threshold inequalities; see the \proofsrc{round2/proofs/astra-late/source/round2/open/age-order/hub-ranking/beta-tests/README.md}{threshold and moment comparison}. Theorem~\ref{th:stoch-order} below supplies them by a separate argument. The conditional comparisons proved here average over the two placements; they are not conditioned on the tree after both births.

\sectag{CERT}
\section{Exact certificates for hub probabilities and moments}\label{app:hub-certificates}
The finite comparisons below are in the units of the signed numerator \eqref{eq:hub-signed}. A cutoff in a comparison process is not a cutoff on the physical tree or on the birth labels. All decimal constants in the inequalities are terminating rationals. The linked programs use exact integer or rational arithmetic, including directed rounding where indicated; a displayed decimal diagnostic is not the quantity tested.

\subsection{Physical continuation bounds}

\paragraph{Hubhood.}
At a marked core of degree $d$ track $j$ actual leaves; the tracked birth and loss intensities are $j$ and $j/d$. On an extra gain track the fresh leaf. When the last tracked leaf is actually promoted, it has degree two and can send back to the mark. Until the mark next gains, that neighbour, at degree $k$, sends back at rate $1/k$ and grows at rate at most $k-1+1/d$. A return comparison with terminal zero at $k=32$ uses
\[
 R_d(k)=\frac{1+k(k-1+1/d)R_d(k+1)}{1+k(k-1+1/d)}.
\]
The backward values decrease in $k$ and lie in $[0,1]$, so faster competing growth gives a lower bound, and any other gain of the mark is at least as good as the specified return. At cutoff degree 20 the payoff $B(20,j)=1-20^{-j}$ follows by the linear birth--death comparison and the leaf-hit argument in the proof of Theorem~\ref{th:hub-order}. Use $B(d,0)=R_d(2)B(d+1,1)$ only after the last actual loss, and for $j>0$ use $B(d,j)=[d B(d+1,j+1)+B(d,j-1)]/(d+1)$. The checked inequality $B(d+1,j+1)\ge B(d,j)$ makes extra gains favourable. Each race exits almost surely: in the finite return range there is a positive lower bound on its send intensity, and at positive stock a tracked jump or gain has rate at least one. There are finitely many stages below the cutoffs since a successful gain increases degree and otherwise tracked stock falls.

An adverse marked leaf at a parent of degree $k$ and with $j$ other tracked leaves has first-send upper $U(k,j)$. The parent has a permanent nonleaf neighbour and $k\ge j+2$. Before the send, the parent gains at rate at least $j+1$, loses tracked siblings at rate $j/k$, and sends to the mark at rate $1/k$. The test $1/(j+1)$ with send payoff one is a supersolution: for $j\ge1$ the drift is at most $-1/(j+2)+1/k$, and the case $j=0$ has payoff already one. For the finite improvement use that barrier at $k=112$ and unwind
\[
 U(k,j)=\frac{1+k(j+1)U(k+1,j+1)+j U(k,j-1)}{1+k(j+1)+j}.
\]
The exact rows decrease along a gain and in $j$; the gain inequality and the terminal supersolution prove the physical bound by stopping, since the marked leaf's own clock forces an exit even without an exterior gain. Use the unrounded barrier in analytic tails (a rounded table can be one tick higher). A favourable tip at parent degree $k$ has lower $L(k)=R_1(k)B(2,1)$ (zero for $k\ge32$): the parent's gain rate is at most $k$ and promotion earns the physical stock-one lower. The serial gap lower $G(k)$ is $B(2,1)-U(2,0)$ plus the credit
\[
 \frac{k^{-1}\{B(3,2)-B(2,1)\}}{k+3/2+k^{-1}}.
\]
Indeed, the middle's tracked gain and loss rates $1,1/2$ balance at the baseline, growth of the old neighbour leaves the baseline unchanged, and the extra rate $1/k$ into the middle earns the displayed improvement on the first monitored event. These are the recurrences of the \proofsrc{round2/proofs/astra-late/source/round2/open/age-order/global-transport/hostile/README.md}{physical hub bounds} and the \proofsrc{round2/proofs/astra-late/source/round2/open/age-order/preplacement-transfer/README.md}{first-send proof}.

\paragraph{Positive powers.}
Put $A_p(a)=(a)_p/\Gamma(1+p)$. For stock $j$ the two useful normalized terminal lowers are
\[
 \phi_p(j)=\begin{cases}
 e^{-1}\sum_{i=0}^{j-1}A_p(j-i)/i!,&0<p\le1,\\
 A_p(j-1),&p\ge1,\ j\ge2,
 \end{cases}
\]
with value zero at $j=0$ and, for the second test, also at $j=1$. At $p=1$ either test is allowed. The first is increasing in stock and satisfies $j\phi_p(j+1)+\phi_p(j-1)=(j+1+p)\phi_p(j)$ for $j\ge1$: its generating series is $e^{z-1}z(1-z)^{-p-1}$, which satisfies the corresponding first-order differential equation by differentiation. For the second, the rate-$j$ birth and unit death generator minus $p$ applied to the test is $p\phi_p(j)(1/(j-1)-1/(j+p-2))\ge0$ for $j\ge2$; the rows at zero and one have nonnegative drift too. Thus in the physical tree, where gain is at least $j$ and loss at most one, both discounted tests are submartingales. Both grow as $j^p/\Gamma(1+p)$ by the gamma ratio (for the first, divide the sum by $A_p(j)$ and dominate the terms by $1/i!$). Domination by the scaled total Yule population in higher moments justifies localization and passage to infinite time. Since the stock is at most degree, these tests bound $\E W^p/\Gamma(1+p)$ from below.

For the middle, truncate its degree at 40 in the fractional harmonic comparison and at 30 with the shifted test for $p\ge1$. Use the return recursion with terminal zero at 96 or 40, respectively,
\[
 R_{d,p}(k)=\frac{1/k+(k-1+1/d)R_{d,p}(k+1)}{p+1/k+k-1+1/d},\quad
 B_p(d,0)=R_{d,p}(2) B_p(d+1,1),
\]
\[
 B_p(d,j)=\frac{j B_p(d+1,j+1)+(j/d)B_p(d,j-1)}{p+j+j/d}.
\]
The return rows decrease in $k$ because each next value is at most $1/(1+p(k+1))$ by backward induction. The stock row at a fixed degree is increasing and bounded above by the gain successor. To see this, write $v_j=B_p(d,j)$, $g_j=B_p(d+1,j+1)$; $(1+p)v_0\le g_1$, and if $v_{j-2}\le v_{j-1}$ the preceding row implies $g_{j-1}\ge(1+p/(j-1))v_{j-1}$. Hence $g_j\ge(1+p/j)v_{j-1}$ and $v_{j-1}\le v_j\le g_j$ by the displayed weighted average. These inequalities justify optional stopping with extra physical gains, as in hubhood. Favourable parallel tips use the same discounted return with $d=1$, cutoff 32 and postpromotion payoff the uniform middle lower. For $p<1/4$ in the stock-sensitive comparison, the serial host of degree $k$ can be retained in a preliminary race: gain intensity $1+1/k$, loss $1/2$, and host growth with unknown intensity in $[0,k-1/2]$. Test both ends of this interval (drift is affine in that intensity) with gain payoff $B_p(3,2)$, loss payoff $B_p(2,0)$ in the actual return state, and continuation at $k+1$ on growth. Cut at 64 and take the maximum with the separately valid uniform middle bound.

For a moment upper at a leaf with parent degree $k$, discount up to first promotion. After promotion its degree is two, and its future gain rate at degree $D$ is at most $D-1+1/k$. An affine birth process with rate equal to current effective mass and unit increments has scaled Gamma limit of shape its initial mass $a$ (the increment count at time $t$ has generating function $[e^{-t}/(1-(1-e^{-t})z)]^a$). Thus from promotion use $A_p(a)$, $a=1+1/k$, or the following sharper cap:
\[
 Q_p(k)=\frac{a A_p(a+1)+(A_p(a-1/2)+A_p(1+1/(k+1)))/2}{p+a+1}.
\]
For this bound race a rate-$a$ candidate incoming gain, accepting real incoming gains by thinning, against the marked vertex's own rate-one ring. At a candidate its degree is at most three. At its own ring without a preceding gain, either its fresh leaf becomes a core or the old parent's degree grows, each with probability one half. The effective masses in those two cases are at most $a-1/2$ and $1+1/(k+1)$. The independent candidate intensity is used only before the first event; it is an incoming channel and does not replace the tag's outgoing clock. This proves the cap by integration of the exponential discount. By the gamma recurrence it is a mixture of $A_p(a)$ and the average of the two smaller values, with weight $1-1/(p+a+1)$ on the larger value. All three values and this weight decrease with $k$, since Gamma moments increase with shape. Thus before promotion use
\[
 U_p(k)=\frac{Q_p(k)/k+U_p(k+1)}{p+1+1/k},
 \qquad U_p(K)=Q_p(K)/(1+pK).
\]
This is decreasing by the backwards bound $U_p(k+1)\le Q_p(k+1)/(1+p(k+1))$ and the first-step equation; extra parent gains have nonpositive drift. The boundary follows from the promotion intensity bound $1/K$ and monotonicity of $Q_p$. The identical argument applies with the simpler $A_p(a)$ cap. At $p=1$ with $K=64$ these rational tests certify $U_1(2)<387/1000$; see the \proofsrc{round2/proofs/astra-late/source/round2/open/age-order/hub-ranking/tip-moments/README.md}{tip proof} and \proofsrc{round2/proofs/astra-late/source/round2/open/age-order/hub-ranking/tip-moments/check.py}{tip checker}.

In the stock-sensitive fractional range use the finer state $(k,j)$ before promotion: replace the numerator of the hub first-send recurrence by $Q_p(k)+k(j+1)U_p(k+1,j+1)+jU_p(k,j-1)$ and add $pk$ to its denominator. At cutoff use $(1+1/k)\min(1/(j+1),1/(1+pk))$, from the stock barrier and promotion hazard respectively; $A_p(1+1/k)\le1+1/k$ for $p\le1$. The interval implementation enlarges the upper to at least its gain successor and checks the gain and stock inequalities. Thus it is a physical supersolution even with extra immigration.

\subsection{From local inequalities to uniform signed bounds}

Let $l=d-c$ at an old core of degree $d$ and core degree $c$, and put $s_l=1/(2d)$, the coefficient at each old leaf. Use $G(k)$ for the serial lower (middle lower minus tip upper), $L(k)$ for a favourable leaf and $U(k,j)$ for an adverse one; for the stock-independent moment ranges omit $j$ from $U$. The old-core credit, after paying its old-leaf edges and including those leaves' serial credits, is
\begin{align*}
 C={}& wG(d+1)+l s_l G(2)\\
 &-l(s_l-s)_+ (U(2,0)-L(d+1))_+
    -l(s-s_l)_+(U(d+1,l-1)-L(2))_+ .
\end{align*}
Omit the debit terms at $l=0$. In a parallel completion on an old leaf edge, the old leaf is promoted, so its new tip has a degree-two parent and no other leaf siblings; the other parent has only $l-1$ old leaf siblings left. On an old core edge with opposite profile $(e,j,t)$, the paired debit is bounded by
\[
 \max\{0,(s-t)(U(d+1,l)-L(e+1))_+,\
             (t-s)(U(e+1,j)-L(d+1))_+\} .
\]
For the outward orientation from the first endpoint one may replace $t$ by its lower $1/[d e(e+1)]$; the coarser bound $sU(d+1,l)$ also suffices where used. Call the resulting directed upper $J$. With baseline $b$, the direct edge costs at most $b$ plus the two directed excesses $(J-b)_+$. No edge pays two baselines.

Contract maximal chains of internal bare degree-two cores. Retain the original degrees and receives at the endpoints, and pay a separate fee at each incident chain port. A chain of length $k=1$ has bare coefficient $(x+y)/6$ for endpoint reciprocals $x,y$; if $k\ge2$ its end coefficients are $(x+1/2)/6,(y+1/2)/6$, and each further interior coefficient is $1/6$. Each bare site earns $2sG(3)$ and each bare--bare edge costs at most $(U(3,0)-L(3))_+|s-t|$. Add the paired endpoint debits, subtract the bare-site credits and both endpoint fees, and check the result is at most $b$ for $k=1,2,3$. Every extra site then increases credit without increasing debit. These tests use convexity of each paired debit in the ranged coefficients, keeping exact neighbour degrees up to the stated transfer cutoffs and interval endpoints beyond them. At each retained vertex check
\[
 C+a(c-2)-b-\sum_{\rm direct}(J-b)_+-\sum_{\rm chain}\mathrm{fee}
     +\mathrm{receipts}-\mathrm{donations}\ \ge0 .                 \tag{*}
\]
Only actual direct slots can fund receipts, each funded at most once. At fixed degree, stock and port list, the variable reciprocal sum enters the credit concavely and the costs convexly, so the endpoint tests suffice. For transfer counts one may use the vertices of their count simplex when the same concavity applies. The contracted tree still has at least two vertices (original core terminals cannot be suppressed); its $K-1$ edges and $\sum(c-2)=-2$ give $\mathcal N\ge b+2a$ by summing (*). This proves the implication from the certified inequalities, even on an arbitrarily long core path.

\paragraph{Exact hub inequalities.}
Use $a=.00634$, $b=.017$, per-chain fees $.021$ at degree two, $.0037$ at degrees 3--9 and $.00625$ at degree at least 10. The latter endpoints can instead pay $.00356$ on direct slots to fund the flagged low profiles $(d,c)=(3,2),(3,3),(4,4)$. The \proofsrc{round2/proofs/astra-late/source/round2/open/age-order/global-transport/hostile/edge-dual/check.py}{direct and high-port checker} and \proofsrc{round2/proofs/astra-late/source/round2/open/age-order/global-transport/hostile/interfaces/low-end.py}{low-end and chain checker} test (*): at low degree enumerate multisets of slots $2,\ldots,9$ and a high symbol with reciprocal in $[0,1/10]$; pay the maximum of the direct excess and chain fee on each degree-two slot. High donors through degree 112 use the receive endpoints and the five fee/receive corners (direct unpaid receive $0,1/2$, direct donor $1/4,1/2$, chain $1/2$). The high endpoint debit in a chain is bounded in its outward direction by $.012$ (finite stock tests through degree 41, then $sU\le1/[2(d+1)]$); in the other direction keep $L(e+1)$ through degree 30, and beyond use $L=0$ and reciprocal endpoints $0,1/31$. Low endpoints keep actual stocks and paired slope and range their other reciprocal sum from 0 to $(c-1)/2$. The high--high chain cost before fees is at most $70627717243/2400000000000<.017+2(.00625)$; the low--high and low--low costs after fees are below $.016905$ and $.010694$, respectively.

For clarity, the high-degree reduction uses the unrounded bounds $sU\le1/[2(d+1)]$ and $U(d+1,l-1)\le1/l$ when $l>0$. The first follows from $s\le(d+l)/[2d(d+1)]$ and $U(d+1,l)\le1/(l+1)$. For $d\ge113$ direct excess is zero. If $l/d\le1/2$, the old-leaf debit is at most $(l/d) U(2,0)/2$ (in the other orientation it vanishes for $l\ge4$ by $L(2)>1/4$, and for $l\le3$ use $s-s_l\le(l-1)/[2d(d+1)]$). Thus in (*) the loss from stock is bounded by half the negative bracket $g\,113/114+G(2)/2-U(2,0)/2$, with $g=.16694235$ a uniform serial lower. Direct-donor ports leave $a-.00356$ per port and there are at least $113/2$ ports in this case. If $l/d>1/2$, leaf serial credit pays the possible old-leaf debit, since $s_l-s\le1/[2d(d+1)]$ when positive, while the other debit vanishes; centre serial credit alone is at least $113g/228$. Chain ports have actual reciprocal $1/2$; to mix bills apply concavity in the port-count simplex, using the all-direct and all-chain endpoints at each stock. At the all-chain endpoint the possible old-leaf debit is paid by leaf credit for \emph{every} stock, using $s_l-s\le1/[2d(d+1)]$ and the preceding $l\le3$ bound for the other orientation. Then for $l/d\le1/2$ use $(a-.00625)113/2$ plus the stock credit rather than the direct-donor estimate: equivalently, minimize the affine lower $g\,113(1+l/d)/228+(a-.00625)d(1-l/d)$, which exceeds $2a+b$ at both stock endpoints. For high terminals, $l=d-1$ gives $C\ge(9/11)g+(9/20)G(2)$, and the checker subtracts the sole port charge. The full \proofsrc{round2/proofs/astra-late/source/round2/open/age-order/global-transport/hostile/edge-dual/README.md}{edge proof} and \proofsrc{round2/proofs/astra-late/source/round2/open/age-order/global-transport/hostile/interfaces/low-end.md}{chain proof} give the row definitions. Import the scalar and stock tests through the edge checker, and run the low-end checker for simultaneous endpoint bills.

For the seed the serial weights sum to one and edge coefficients tie; for stars of degrees 2 and 3 the \proofsrc{round2/proofs/astra-late/source/round2/open/age-order/hub-ranking/check.py}{seed and star checker} verifies numerator lowers $.18940948,.20181908,.22577249$ in the three respective cases. On a star of degree $d\ge4$, the adverse tip at its centre has upper at most $1/d<L(2)$, so the old-edge debit vanishes and old-leaf serial credit alone gives $G(2)/2>371/12500$. This completes the hub certificate.

\paragraph{Exact moment inequalities.}
The certified bounds for the same signed score, including exceptional trees, are as follows. The proofs and executable tests for each row are linked immediately below.
\[
\begin{array}{c|c|l}
 p\text{ range}&\mathcal N\text{ lower}&\text{checks}\\ \hline
 (0,473/512]&33/1000&473\text{ stock-sensitive cells (multi-core)}\\
 {}[473/512,1]&1/40&39\text{ harmonic-stock cells (multi-core)}\\
 (0,1]&13/100&\text{edge seed and all stars}\\
 {}[1,3]&1/40&256\text{ cells}\\
 {}[3,7]&1892360505093/10^{15}&2048\text{ cells}\\
 {}[7,\infty)&22890923/1058158080&\text{analytic comparison}
\end{array}
\]
In all cells lower boundary products $A_p(j)$ at integer $j$ are evaluated at the lower exponent and discounts in lower recursions at the upper exponent. Use the alternating lower through term 17 for $e^{-1}$ and round down. The first cell uses the zero endpoint in polynomials only, not the convention $Z^0=1$. For upper Gamma ratios at fractional $p$ use the Euler product $\prod_{i\ge0}(1+p(a-1)/[(i+1)(i+a+p)])$. It follows by taking the ratio of Euler's finite gamma products; each factor increases in $p$ if $a>1$ and decreases if $a<1$. Beyond 48 terms the upward tail is at most $(1-p(a-1)/48)^{-1}$ if $a>1$ (sum of inverse squares at most $1/48$ and $e^x\le(1-x)^{-1}$); if $a\le1$ discard the tail for an upper bound. For $p\ge1$ and $0<a\le1$ use the chord between integer powers of the Beta moment $A_p(a)$, decreasing and convex in $p$, at the lower cell endpoint. For a shape $1+a$ multiply by $1+p/a$ at the upper endpoint. Use positive mixtures for the improved cap with the appropriate interval ends of each weight, and round upper rows up. Thus the tests bound entire real intervals, not just exponent samples.

\emph{Through $473/512$.} Use stock at the adverse parent, cutoff 48 in the first 36 cells and 32 thereafter. The residual (*) uses $a=b=.015$ in the first 36 cells and $.011$ thereafter. The chain fees are $.007$ and (at a terminal of degree two) $.022$ or $.020$, respectively. In the first 36 cells degree at least 15 budgets $.005$ per direct slot; a direct opposite degree at most 3 receives it. First use $J\le sU$ at receive endpoints $k/2,c/2$ for $k$ chain slots. If insufficient, retain exact direct degrees through 14 for $d\le4$ and through six otherwise; range the reciprocal sum of higher slots and drop their favourable lower and opposing coefficient in $J$. For large $d\ge D$ ($D=96,64$ respectively), write $x=l/d$, $T=(1+1/(D+1))/(D+1)$, $g$ for the uniform serial gap and $z=Dg/(D+1)+(g-U(2,0))/2$. The stock and hazard caps show direct excess vanishes by $sU\le T/2$ (use $l+c/2$ in $s$ and the minimum of stock and degree bounds); the core-side old-leaf debit is at most $T$ in total, using $U(d+1,l-1)\le(1+1/(d+1))/l$ and $s\le1/(d+1)$. For $x\le1/2$ the credit before port allowances is at least $\min(0,z)/2-T$, and at least $D/2$ ports still contribute $a-.007$ less a possible donor fee. For $x>1/2$ the opposite old-leaf debit is at most $U(2,0)/[4(D+1)]$, since $w\ge l/(d+1)$; the serial credit is at least $Dg/[2(D+1)]+g/4$. Subtract both debit bounds. These tail inequalities pay $2a+b$ on every cell. Chain tests retain stock at low endpoints, exact reciprocals and maximal stock exposure through the tip cutoff, and reciprocal endpoints from zero to the inverse cutoff beyond, with zero favourable lower there. See the \proofsrc{round2/proofs/astra-late/source/round2/open/age-order/hub-ranking/mean-ledger/fractional/README.md}{fractional stock proof} and \proofsrc{round2/proofs/astra-late/source/round2/open/age-order/hub-ranking/mean-ledger/fractional/check.py}{473-cell checker}.

\emph{From $473/512$ to 3.} The upper fractional cells use $a=.0085,b=.008$, fees $.004,.0131$ (the latter for degree-two terminals). Use the degree-only upper, count the chain ports, and at fixed count take both reciprocal endpoints as above. At $d\ge32$ direct excess is zero and $U(d+1)<L(2)$ by the validated cap $u=1/(33p_{\rm lo})$, where $p_{\rm lo}$ is the lower endpoint of the box for the exponent $p$. With $z=(32/33)g+(g-U(2))/2$ the credit with available port allowances is bounded below by $xz+32(1-x)(a-.004)$; both ends pay $2a+b$. For $p\ge1$ the degree cap is $U_p(k)\le1/k$: $A_p(1/k)\le1/k$ by the Beta moment and the recurrence gives $A_p(1+1/k)/(1+pk)\le1/k$. This gives the same suppression test and tail with $u=1/33$. Alternatively reserve $[w+l/(2d)]/40$ per core and charge $csU(d+1)$; the reserves sum to at least $1/40$ since $\sum r_u=n$ and $d_u+1\le n$. The finite sitewise test uses both receive endpoints below 32; above, the outgoing centre-side debit totals at most $w/33$ and opposite old-leaf debit at most $lU(2)/(2d)$. With $h=g-1/40$ the checked sufficient tests are $h>1/33$ and $(32/33)(h-1/33)+(h-U(2))/2\ge0$. Chain exposures at degrees 10--30 can use $sU\le U(d+1)/(d+1)$; beyond, use the reciprocal endpoint $1/31$ and the analytic degree cap. The \proofsrc{round2/proofs/astra-late/source/round2/open/age-order/hub-ranking/fractional-middle/README.md}{harmonic proof} and \proofsrc{round2/proofs/astra-late/source/round2/open/age-order/hub-ranking/fractional-middle/check.py}{39-cell signed checker} give the upper fractional range. The \proofsrc{round2/proofs/astra-late/source/round2/open/age-order/hub-ranking/mean-ledger/real/README.md}{real-power proof} and \proofsrc{round2/proofs/astra-late/source/round2/open/age-order/hub-ranking/mean-ledger/real/check.py}{256-cell checker} discharge the next range by the sitewise test on 87 cells and the chain test on 169. The seed and finite star profiles and the unbounded star line for fractional powers are checked separately by the \proofsrc{round2/proofs/astra-late/source/round2/open/age-order/hub-ranking/mean-ledger/fractional/star_bound.py}{512-cell star checker}: their exact respective minima are $406542520567/2500000000000$, $2085176012479/12000000000000$, and $265041348327659/2012097656250000$, each greater than $13/100$. The star tail uses centre serial credit minus outgoing old-leaf debit (the centre coefficient exceeds each leaf's); the real-power checker uses the same star sign with $1/40$ reserve.

\emph{From 3 to infinity.} Let $A=\sum w_u\ge1$. The sum of absolute paired coefficients is at most $A=\sum d_us_u$. The part $X$ adverse at an old leaf is at most $A/3$: the old leaves at a core of degree $d$ have total weight $A_l=l/(2d)$ and adverse absolute coefficient at most $A_l(1-2l/(d+1))_+\le A_l(d-1)/(d+1)$ for $l>0$. Three times the latter is at most $A_l+w_u$ since $w_u\ge l/(d+1)$. On the edge seed $X=0$. If $M_p$ bounds the middle below and $U_{2,p}\ge U_{3,p}$ bound tips above for parent degrees at least 2 and at least 3, respectively, this gives $\mathcal N\ge A[M_p-(4U_{2,p}+2U_{3,p})/3]$. On $[3,7]$ compute the uniform caps by the degree-only first-promotion recurrence with unrefined $A_p(1+1/k)$ and parent cutoff 48, using the closed cell rules already described. The \proofsrc{round2/proofs/astra-late/source/round2/open/age-order/hub-ranking/mellin/check.py}{signed box checker} regenerates the $[1,3]$ lower inputs and checks the exact minimum on $[3,7]$ (compare the printed fraction to the lower endpoint in the table).

For $p\ge7$ use the absorbed stock chain with births $j$ and unit death. Put $P_q(j)=A_q(j)$ for $j\ge1$, $P_q(0)=0$, with $P_0(j)=\mathbf1_{\{j>0\}}$. The gamma recurrence gives $(\mathcal L-q)P_q=-P_{q-1}$ for $q\ge1$, also at $j=1$. With $u=1-e^{-t}$ the function $r_q(u)=e^{-qt}\E P_q(j(t))$ has derivative $-r_{q-1}(u)$ and starts at one; $0\le r_q\le1$ by the pure Yule comparison. Seven integrations therefore give the alternating lower $\sum_{i=0}^7(-1)^i/i!$ at $u=1$, with uniform integrability as for the stock tests. A leaf at parent degree at least $k$ has the locked affine cap $A_p(1/k)$ from initial degree one. This decreases in $p$, so substitute $p=7$ in the two adverse terms. The exact alternating lower minus the resulting Beta cost is the last rational in the table. See the \proofsrc{round2/proofs/astra-late/source/round2/open/age-order/hub-ranking/mellin/README.md}{large-power proof}. The table entries are all at least $9/5000$.

For reproduction use Python 3 with assertions enabled (no \texttt{-O}), invoking the linked checkers with \texttt{-IB}. If an internal historical digest authenticates a Markdown input before checking, use the \proofsrc{round2/verify/astra-late/with-originals.py}{authenticated-input launcher} from the repository root, passing the checker path relative to the original source root. The \proofsrc{round2/proofs/astra-late/source/round2/open/age-order/hub-ranking/all-powers/positive/check.py}{power-seam checker} tests the joins and normalizations and authenticates inputs; it is not a substitute for the individual box and star replays. The \proofsrc{round2/proofs/astra-late/independent/audit_claim26.py}{independent hub audit} and \proofsrc{round2/proofs/astra-late/independent/audit_claim27.py}{independent moment audit} give additional exact reconstruction and physical controls, not a finite-tree sampling argument for all sizes.

\sectag{ORD}
\section{Stochastic order at every threshold}\label{sec:ord}\label{app:stoch-order}
\tagnext{ORD}
\begin{theorem}[Strict stochastic order by birth label]\label{th:stoch-order}
For every $n\ge2$, every history through size $n$, before either of the next two labels is placed, and every $0\le x<1$,
\[
 \Prob(Z_{n+1}>x\mid\mathcal F_n)>\Prob(Z_{n+2}>x\mid\mathcal F_n).
\]
Also $\Prob(Z_2>x)>\Prob(Z_3>x)$ for every $0\le x<1$. Consequently $Z_1$ and $Z_2$ have the same law and
\[
 \Prob(Z_u>x)>\Prob(Z_v>x)\qquad\text{for all }2\le u<v\text{ and }0\le x<1,
\]
so $Z_2,Z_3,Z_4,\ldots$ decrease strictly in first-order stochastic order.
\end{theorem}
Its proof has been checked by one model family other than its author's and by no human (Section~\ref{sec:provenance}).
For $x\ge1$ both tails vanish and for $x<0$ both equal one, so the strict inequalities are for the whole nontrivial range. At $x=0$ the theorem recovers the strict order of Theorem~\ref{th:hub-order} without its quantitative gap. Integrating the tails against $p x^{p-1}\,dx$ also recovers the strict order of positive moments, but not the quantitative bounds of Theorems~\ref{th:hub-order} and~\ref{th:hub-moment-order}; conversely, moment order does not imply this theorem. As in Theorem~\ref{th:hub-order} the comparison conditions on the tree before both placements; conditioning on the completed tree after both births can reverse it (Section~\ref{sec:further-order}).

\subsection{Structure of the proof}\label{sec:ord-guide}
The complete proof is reproduced in Section~\ref{sec:ord-proof}. This subsection summarizes it, separates the analytic reduction from the finite checks, and records the review. The proof uses only the existence of the limits $Z_v$ from Section~\ref{sec:conservation} and none of the other numbered theorems of the paper or this supplement.
\begin{proof}[Proof outline]
Taking expectations in the conditional inequality gives $\Prob(Z_b>x)>\Prob(Z_{b+1}>x)$ for every $b\ge3$; the founder inequality covers $b=2$, transitivity covers every pair $2\le u<v$, and exchanging the two initial labels gives the founder tie. The conditional inequality starts from \eqref{eq:hub-signed} with the tail $\Prob(Z>x\mid\cdot)$ of the marked vertex as the continuation score. Its positive serial terms are bounded below, and its adverse parallel and tip terms above, by three physical comparison laws: a Poisson mixture of Beta tails for a vertex with tracked leaves, a Beta tail locked at the current attachment weight, and a first-promotion bound for a pending leaf. Piecewise-linear price tables, run on two logarithmic clocks in the tree size, extend these laws to every future size, and a tree ledger of the same kind as in Section~\ref{app:hub-certificates} reduces every finite tree to local inequalities. Two exact integer certificates check those inequalities on a finite cover of the thresholds, an analytic bound handles the remaining high thresholds, and a separate size-four decomposition proves the founder inequality. The rest of this section gives the guide (\S\ref{sec:ord-guide}), the review (\S\ref{sec:ord-review}), the complete proof (\S\ref{sec:ord-proof}) and the sources (\S\ref{sec:ord-sources}).
\end{proof}
\paragraph{Signed identity.}
Fix $n\ge2$ and a tree $T_n$, write $b=n+1$ and $m=n+2$, and use the old-tree quantities $d_u$, $r_u$, $w_u=r_u/(d_u+1)=d_us_u$ of \eqref{eq:hub-signed}. With the tail $\Prob(Z>x\mid\cdot)$ of the marked vertex as the continuation score, \eqref{eq:hub-signed} takes the form (Lemma~2.1 of the proof)
\begin{align*}
 &n(n+1)\bigl(\Prob(Z_b>x\mid T_n)-\Prob(Z_m>x\mid T_n)\bigr)\\
 &\qquad=\sum_u w_u\Delta_{\rm ser}(u)-\sum_{\{u,v\}}\bigl[(s_u-s_v)_+\Delta_{\rm par}(v,u)+(s_v-s_u)_+\Delta_{\rm par}(u,v)\bigr].
\end{align*}
Here $\Delta_{\rm ser}(u)$ is the middle-minus-tip tail difference when $b$ attaches to $u$ and $m$ to $b$, and $\Delta_{\rm par}(v,u)$ is the later-minus-earlier difference when $b$ attaches to $v$ and $m$ to $u$. Twins and nonadjacent pairs cancel exactly. In every adverse parallel term the host of the later label has a permanent non-leaf neighbour, so its degree $k$ and its number $j$ of other leaf neighbours satisfy $k\ge j+2$ for ever.

\paragraph{Three comparison laws.}
Write $B_{N,r}(x)=\Prob(\mathrm{Beta}(r,N-r)>x)$. (i) A vertex with $j$ tracked leaf neighbours at size $N$ has $\Prob(Z>x\mid T_N)\ge e^{-1}\sum_{i=0}^{j-1}B_{N,j-i}(x)/i!$, by a submartingale built from the P\'olya-urn identity. (ii) A vertex with current attachment weight $r<N$, the sum of its neighbours' reciprocal degrees, has $\Prob(Z>x\mid T_N)\le B_{N,r}(x)$, since that weight can grow only through its own gains. (iii) A leaf whose host has degree $k$ and $j$ other leaves has $\Prob(Z>x\mid T_N)\le\Prob(Z>0\mid T_N)\le U_{\rm hub}(k,j)$, a first-promotion bound computed backwards from the caps $1/(j+1)$ on the branch $k\ge j+2$ and $\min(1,2/(j+1))$ otherwise.

\paragraph{From the laws to every future size.}
Put $\tau=(m-1)\log(1/(1-x))$, or $y=\log(1/(1-x))$ for small $m$. Price tables are piecewise linear in a logarithmic time variable and nonincreasing in it. Two clocks, $\log(N/m)$ and $\log((N-1)/(m-1))$, advance by less and by more than $1/N$ at each step, which is what a lower and an upper comparison need. Nodal generator inequalities at both ends of each time slab then make the tables a submartingale and a supermartingale along the physical chain at every size $N\ge m$. A step that crosses a knot is affine in its overshoot, so both ends of the overshoot are checked. The lower tables track leaf stock, the return of a promoted leaf and a serial middle; the upper tables track a pending adverse leaf and the race after its promotion. At its cutoffs each table stops at one of the three laws or at another valid price, such as zero for a return or a hand-off to a simpler row.

\paragraph{Tree ledger and coverage.}
The serial credits pay the adverse terms through site credits, edge debits, fees on contracted chains of bare degree-two cores, and the identity $\sum(c-2)=-2$ over the sites of the contracted core tree, as in Section~\ref{app:hub-certificates}. Concavity or convexity in the continuous parameters reduces each site and chain to finitely many endpoint tests; chains longer than three, and sites of degree above a fixed cutoff, are handled analytically. The edge seed and the stars are checked separately. For $m\ge200$ the program \ordersrc{round3/q3-birth-order/cert_all_m.cpp}{\texttt{cert\_all\_m.cpp}} checks 52 contiguous $\tau$-boxes covering $[0,6]$, uniformly in $m$. For $4\le m\le199$ the program \ordersrc{round3/q3-birth-order/cert_small_m.cpp}{\texttt{cert\_small\_m.cpp}} runs an exact backward induction over integer sizes on 84 contiguous $y$-boxes covering $[0,2]$. When $\tau\ge6$ an analytic bound applies to every $b\ge3$ and every tree: the total adverse parallel weight is at most the total serial weight, at most a third of it comes from old leaves, and the stock lower then beats the Beta upper bounds of shapes one half and one third by an explicit positive margin. The first program evaluates that margin exactly, and the two coverages meet without a gap.

\paragraph{Founder seam.}
After the third and fourth births, $\Prob(Z_2>x)-\Prob(Z_3>x)$ is one third of the centre-minus-leaf difference on the four-vertex star plus one twelfth of the middle-minus-leaf difference on the four-vertex path; the remaining branches are symmetric. Both differences are positive by the small-size certificate for $y\le2$ and by the analytic bound at size four for $y\ge2$.

\paragraph{What is analytic and what is computed.}
The signed identity, the three comparison laws, the passage from nodal inequalities to all future sizes, the ledger summation over arbitrary finite trees, the high-threshold bound and the founder decomposition are proofs. The computer checks finitely many inequalities: the price tables on each box, the local site and chain tests, and the scalar margins. Both programs use signed 128-bit integers with directed rounding and enclose the exponential, logarithm, Gamma and incomplete-Gamma values they need; no floating-point value is trusted. A passing run therefore establishes the finite inequalities, and the analytic reduction turns them into Theorem~\ref{th:stoch-order} for every tree, label and threshold. A run does not check the reduction.

\subsection{Review and recorded repairs}\label{sec:ord-review}
\paragraph{Review and recorded repairs.}
GPT audited the \ordersrc{round3/q3-birth-order/REVIEW-Q3.md}{full argument} and the \ordersrc{round3/q3-birth-order/REVIEW-IMPLEMENTATION.md}{implementation} on these exact bytes, and returned \texttt{HOLDS} for the conditional and the marginal statements. It recorded four local points, which the outline above already respects. The strict parameter monotonicities of $B_{N,r}$ stated before Lemma~3.1 of the proof fail at $x=0$; non-strict monotonicity on $[0,1)$ and strict monotonicity on $(0,1)$ suffice, and strictness at $x=0$ comes from the ledger. In Lemma~3.3 the first-jump mixture is fractional-linear, not affine, in the host growth rate; clearing its positive denominator gives the same two endpoint tests. For hosts beyond the finite table, the caps $1/(j+1)$ on the branch $j\le k-2$ and $\min(1,2/(j+1))$ otherwise are the intended extension. The 12{,}105 printed checks include the 798 attempts on subdivided boxes. The implementation audit also showed that the positive-overshoot test at a frozen lower row is needed, by an actual crossing at which the zero-overshoot test alone passes with equality. No model family other than the author's and the reviewer's, and no human, has checked this proof.

\subsection{The complete proof}\label{sec:ord-proof}
This subsection reproduces the complete proof as one text. The repository keeps it as a single frozen file, split for layout only into two Markdown files. Sections~1--9 are reproduced. Section~10, the replay commands and observed output, stays in the repository; the replay is described in Section~\ref{sec:ord-sources}. The four points of Section~\ref{sec:ord-review} apply to this text. The first three are also printed as framed review notes where they apply: the remark on monotonicity of $B_{N,r}$ before Lemma~3.1, the first-jump mixture in the proof of Lemma~3.3, and the caps for hosts beyond the finite table. The fourth, the count of 12{,}105 printed checks, concerns the program output in Section~\ref{sec:ord-sources}. A further framed note marks an informal parenthetical in the proof of Lemma~4.1. The notes are the only additions to the source text.
\begin{center}\footnotesize\begin{tabular}{@{}ll@{}}\toprule
Frozen proof (SHA-256) & \texttt{e85e9cebfbde5439fad9246da2fec2c6eb7f014b09248e960c20d016411ce191}\\
Sections 1--3 & \ordersrc{round3/q3-birth-order/PROOF.md}{stochastic-order proof, first file}, lines 12--166\\
\quad SHA-256 & \texttt{5e9a5e388b1570902424d68f653aaeca6d93a304c113c604507c1ae04223ed33}\\
Sections 4--9 & \ordersrc{round3/q3-birth-order/PROOF-PART-2.md}{stochastic-order proof, second file}, lines 9--236\\
\quad SHA-256 & \texttt{9a3a43c2534ea087fb23a735c1cde4975147bf87387b8bda2d16fbc00d8a2653}\\
Commit & \texttt{62081cad01112c42c2011ac3eed11aa3aa665e09}\\
\bottomrule\end{tabular}\end{center}
% BEGIN TRANSCRIPTION round3/q3-birth-order/PROOF.md lines 12-166 sha256 5e9a5e388b1570902424d68f653aaeca6d93a304c113c604507c1ae04223ed33 commit 62081cad01112c42c2011ac3eed11aa3aa665e09
\begin{transcription}
\subsubsection*{1. Exact Model and Main Theorems}

\trhead{1.1 The Random Friend Tree Process}
Start from the two-vertex tree $T_2$ consisting of a single edge $\{1, 2\}$. For each integer $n \ge 2$, conditionally on $T_n$ on vertex set $\{1, \dots, n\}$ with degree function $d_n(u) := \deg_{T_n}(u)$, form $T_{n+1}$ by adding vertex $n+1$ attached to a random vertex $W_n \in V(T_n)$ sampled by picking $U_n \in V(T_n)$ uniformly at random and then picking $W_n$ uniformly among the neighbors of $U_n$ in $T_n$. Equivalently,

\fitdisplay{%
\mathbb{P}(n+1 \sim u \mid T_n) = \frac{w_n(u)}{n}, \qquad w_n(u) := \sum_{v \sim u} \frac{1}{d_n(v)}.
}

Since $\sum_{u \in V(T_n)} w_n(u) = \sum_{v \in V(T_n)} d_n(v)\cdot \frac{1}{d_n(v)} = n$, this defines a probability distribution on $V(T_n)$.

For each vertex label $u \ge 1$, the sequence $(d_N(u)/N)_{N \ge u}$ is a nonnegative supermartingale (its conditional expectation at step $N+1$ is $(d_N(u) + w_N(u)/N)/(N+1) \le d_N(u)/N$ because $w_N(u) \le d_N(u)$), so the limiting degree share

\fitdisplay{%
Z_u := \lim_{N \to \infty} \frac{d_N(u)}{N} \in [0, 1]
}

exists simultaneously for all $u \ge 1$ almost surely (Addario-Berry et al., \emph{Random Structures \& Algorithms} 69(1), e70082 (2026); arXiv:2403.20185v1 \S{}3).

\trhead{1.2 Main Results}

\textbf{Theorem 1 (Universal Pre-Placement Conditional Stochastic Dominance).}\\
\emph{For every integer $n \ge 2$, every finite tree $T_n$ on $n$ vertices, and every threshold $x \in [0, 1)$,}

\fitdisplay{%
\mathbb{P}(Z_{n+1} > x \mid T_n) > \mathbb{P}(Z_{n+2} > x \mid T_n).
}

\textbf{Theorem 2 (Complete Resolution of Addario-Berry et al. \S{}7 Question 3).}\\
\emph{The limiting degree shares $(Z_u)_{u \ge 1}$ satisfy exact founder symmetry and strict first-order stochastic dominance from birth label $2$ onward:}

\fitdisplay{%
Z_1 \stackrel{d}{=} Z_2 \succ_{\mathrm{st}} Z_3 \succ_{\mathrm{st}} Z_4 \succ_{\mathrm{st}} \cdots,
}

\emph{that is, $\mathbb{P}(Z_1 > x) = \mathbb{P}(Z_2 > x)$ and $\mathbb{P}(Z_u > x) > \mathbb{P}(Z_v > x)$ for every $2 \le u < v$ and every $x \in [0, 1)$ (with trivial equality at $x \ge 1$, where both sides are $0$, and at $x < 0$, where both sides are $1$).}

\emph{Proof of Theorem 2 from Theorem 1 and the Founder Seam (Section 8).}\\
The graph automorphism swapping $1 \leftrightarrow 2$ in $T_2 = \{1, 2\}$ preserves the transition law, so $Z_1 \stackrel{d}{=} Z_2$. Section 8 proves $\mathbb{P}(Z_2 > x) > \mathbb{P}(Z_3 > x)$ for every $x \in [0, 1)$. For every $b = n + 1 \ge 3$ ($n \ge 2$) and every $x \in [0, 1)$, taking expectation over $T_n$ in Theorem 1 gives $\mathbb{P}(Z_b > x) > \mathbb{P}(Z_{b+1} > x)$. Transitivity gives $\mathbb{P}(Z_u > x) > \mathbb{P}(Z_v > x)$ for all $2 \le u < v$ and $x \in [0, 1)$. $\square$


\subsubsection*{2. Exact Pre-Placement Signed Identity}

Fix $n \ge 2$ and an arbitrary finite tree $T_n$. Write $b := n+1$, $m := n+2 \ge 4$, $d(u) := d_n(u)$, $w(u) := w_n(u) = \sum_{v \sim u} 1/d(v)$, and define the \textbf{edge-drop weight}

\fitdisplay{%
s(u) := \frac{w(u)}{d(u)\big(d(u)+1\big)} \qquad (u \in V(T_n)).
}

Let $T_m(b \sim u, m \sim v)$ be the tree formed from $T_n$ by attaching leaf $b$ to $u$ and leaf $m$ to $v$.

\textbf{Lemma 2.1 (Exact Pre-Placement Decomposition).}\\
\emph{For every finite tree $T_n$ ($n \ge 2$) and every $x \in [0, 1)$,}

\fitdisplay{\begin{multlined}[\linewidth]%
n(n+1)\Big(\mathbb{P}(Z_b > x \mid T_n) - \mathbb{P}(Z_m > x \mid T_n)\Big) = \ftbreak \sum_{u \in V(T_n)} d(u)s(u)\, \Delta_{\mathrm{ser}}(u) \ftbreak - \sum_{\{u, v\} \in E(T_n)} \Big[(s(u) - s(v))_+ \Delta_{\mathrm{par}}(v, u) + (s(v) - s(u))_+ \Delta_{\mathrm{par}}(u, v)\Big],
\end{multlined}}

\emph{where:}
\begin{enumerate}
\item \emph{$\Delta_{\mathrm{ser}}(u) := \mathbb{P}(Z_b > x \mid T_m(b \sim u, m \sim b)) - \mathbb{P}(Z_m > x \mid T_m(b \sim u, m \sim b))$ is the serial gap at $u$: $b$ has degree $2$, $1$ leaf neighbor $m$, and $1$ host neighbor $u$ of degree $d(u)+1$, while $m$ is a leaf on host $b$ of degree $2$ with $0$ other leaf neighbors.}
\item \emph{$\Delta_{\mathrm{par}}(v, u) := \mathbb{P}(Z_m > x \mid T_m(b \sim v, m \sim u)) - \mathbb{P}(Z_b > x \mid T_m(b \sim v, m \sim u))$ is the adverse parallel gap on $\{b \sim v, m \sim u\}$: $m$ is a leaf on host $u$ of degree $d(u)+1$, and $b$ is a leaf on host $v$ of degree $d(v)+1$.}
\end{enumerate}

\emph{Proof.} Condition on $b \sim u$ (probability $w(u)/n$). Vertex $m$ then attaches to:
\begin{itemize}
\item \textbf{$u$ (twin):} $b$ and $m$ are symmetric leaves on $u$ in $T_m$, so their marginal laws match and contribute $0$.
\item \textbf{$b$ (serial):} $U_{n+1}$ must be $u$ (the unique neighbor of $b$) and $W_{n+1} = b$ (probability $1/((n+1)(d(u)+1))$). Multiplying by $w(u)/n$ yields $\frac{w(u)}{n(n+1)(d(u)+1)} = \frac{d(u)s(u)}{n(n+1)}$.
\item \textbf{$v \neq u$ not adjacent to $u$:} attaching $b$ to $u$ does not change any neighbor degree of $v$, so $\mathbb{P}(b \sim u, m \sim v \mid T_n) = \frac{w(u)w(v)}{n(n+1)} = \mathbb{P}(b \sim v, m \sim u \mid T_n)$. Swapping $b \leftrightarrow m$ is a tree isomorphism between $T_m(b \sim u, m \sim v)$ and $T_m(b \sim v, m \sim u)$, so every non-adjacent pair cancels identically.
\item \textbf{$v \sim u$ (an edge $\{u, v\} \in E(T_n)$):} attaching $b$ to $u$ increases $d(u)$ to $d(u)+1$, reducing $w(v)$ by $\frac{1}{d(u)} - \frac{1}{d(u)+1} = \frac{1}{d(u)(d(u)+1)}$. Hence

\fitdisplay{%
\mathbb{P}(b \sim u, m \sim v \mid T_n) = \frac{w(u)w(v) - s(u)}{n(n+1)}, \qquad \mathbb{P}(b \sim v, m \sim u \mid T_n) = \frac{w(u)w(v) - s(v)}{n(n+1)}.
}
\end{itemize}

Pairing the two orientations cancels their common mass by the $b \leftrightarrow m$ isomorphism, leaving net probability $\frac{(s(u) - s(v))_+}{n(n+1)}$ on $\{b \sim v, m \sim u\}$ and $\frac{(s(v) - s(u))_+}{n(n+1)}$ on $\{b \sim u, m \sim v\}$. $\square$

\textbf{Leaf-stock invariant at completion $T_m$.} Let $l(u)$ be the number of leaf neighbors of $u$ in $T_n$. In the adverse configuration $\{b \sim v, m \sim u\}$ along edge $\{u, v\} \in E(T_n)$:
\begin{itemize}
\item If $d(v) \ge 2$, then $v$ remains a non-leaf neighbor of $u$ in $T_m$, so host $u$ has degree $k = d(u)+1$, at least one non-leaf neighbor ($v$), and $j = l(u) \le d(u) - 1 = k - 2$ other leaf neighbors besides $m$.
\item If $d(v) = 1$, then $b \sim v$ turns $v$ into a degree-$2$ non-leaf neighbor of $u$ in $T_m$, so host $u$ has degree $k = d(u)+1$, at least one non-leaf neighbor ($v$), and $j = l(u) - 1 \le d(u) - 1 = k - 2$ other leaf neighbors besides $m$.
\item In the serial tip $\{b \sim u, m \sim b\}$, host $b$ has degree $k = 2$, one non-leaf neighbor ($u$), and $j = 0 = k - 2$ other leaf neighbors besides $m$.
\end{itemize}

Thus \textbf{every adverse pending leaf in Lemma 2.1 satisfies the permanent invariant $k \ge j + 2$} (the host has a permanent non-leaf neighbor, since vertex degrees never decrease).


\subsubsection*{3. Exact Endpoint Comparison Laws}

For real $N > r > 0$ and $x \in [0, 1)$, let

\fitdisplay{%
B_{N, r}(x) := \mathbb{P}\big(\mathrm{Beta}(r, N-r) > x\big) = \frac{\Gamma(N)}{\Gamma(r)\Gamma(N-r)} \int_x^1 u^{r-1}(1-u)^{N-r-1}\,du,
}

with $B_{N, r}(x) := 0$ for $r \le 0$ and $B_{N, N}(x) := 1$. For integer $r \in \{1, \dots, N-1\}$, $B_{N, r}(x) = \mathbb{P}(\mathrm{Bin}(N-1, x) \le r-1) = (1-x)^{N-1}\sum_{i=0}^{r-1}\binom{N-1}{i}\left(\frac{x}{1-x}\right)^i$. Note that $B_{N, r}(x)$ is strictly decreasing in $x \in [0, 1)$, strictly decreasing in real $N > r$, and strictly increasing in $r \in (0, N)$.
\reviewnote{At $x=0$ the strict monotonicities in $N$ and $r$ stated here fail. Non-strict monotonicity on $[0,1)$ and strict monotonicity on $(0,1)$ suffice, and strictness at $x=0$ comes from the ledger (Section~\ref{sec:ord-review}).}

\textbf{Lemma 3.1 (Poisson-Beta Submartingale Lower Bound).}\\
\emph{At any size $N \ge 2$, for any vertex $v \in V(T_N)$ with at least $j \in \{1, \dots, N-1\}$ leaf neighbors and any $x \in [0, 1)$,}

\fitdisplay{%
\mathbb{P}(Z_v > x \mid T_N) \ge H_{N, j}(x) := e^{-1}\sum_{k=0}^{j-1}\frac{1}{k!} B_{N, j-k}(x).
}

\emph{Proof.} Set $H_{N, 0} \equiv 0$. From the exact P\'olya-urn identity $N B_{N, r} = r B_{N+1, r+1} + (N-r) B_{N+1, r}$ and $k(e^{-1}/k!) = e^{-1}/(k-1)!$, telescoping $\Delta B(r) = B_{N+1, r} - B_{N+1, r-1}$ gives the discrete harmonic identity for all $1 \le j < N$:

\fitdisplay{%
N H_{N, j}(x) = j H_{N+1, j+1}(x) + (N - j - 1) H_{N+1, j}(x) + H_{N+1, j-1}(x).
}

Designate $j$ current leaf neighbors of $v$ as tracked leaves, add each subsequent leaf attached to $v$ to the tracked set, and remove a tracked leaf when a new vertex attaches to it. At tree size $s$ with tracked stock $i \ge 1$, the probability of attaching to $v$ is at least $i/s$ (each tracked leaf has degree $1$), the disjoint probability of attaching to a tracked leaf is $i/(s d_s(v)) \le 1/s$, and $H_{s+1, i}$ increases in $i$. Thus $(H_{s, I_s}(u))_{s \ge N}$ is a bounded submartingale for every $u \in (x, 1)$. On $\{Z_v < u\}$, $I_s/s \le d_s(v)/s \to Z_v < u$, so $H_{s, I_s}(u) \le B_{s, I_s}(u) = \mathbb{P}(\mathrm{Bin}(s-1, u) < I_s) \to 0$. Reverse Fatou gives $H_{N, j}(u) \le \mathbb{P}(Z_v \ge u \mid T_N)$, and letting $u \downarrow x$ by continuity of $H_{N, j}$ completes the proof. $\square$

\textbf{Lemma 3.2 (Locked Beta Upper Bound).}\\
\emph{At any size $N \ge 2$, for any vertex $v \in V(T_N)$ with current attachment weight $r_0 = w_N(v) < N$ and any $x \in [0, 1)$,}

\fitdisplay{%
\mathbb{P}(Z_v > x \mid T_N) \le B_{N, r_0}(x).
}

\emph{Proof.} Any attachment to a neighbor $u \sim v$ can only increase $d_s(u)$, hence can only decrease $1/d_s(u)$, while each attachment to $v$ adds a new leaf neighbor of weight $1$. Thus $w_s(v) \le r_0 + (d_s(v) - d_N(v))$ for all $s \ge N$, so $(d_s(v) - d_N(v))_{s \ge N}$ is stochastically dominated by a P\'olya urn starting with $(r_0, N-r_0)$ balls, whose limiting fraction is $\mathrm{Beta}(r_0, N-r_0)$. $\square$

\textbf{Lemma 3.3 (Static Time-Independent Hubhood Upper Bound).}\\
\emph{Fix $K_{\max} \ge 2$. Define $U_{\mathrm{hub}}(K_{\max}+1, j) := \frac{1}{j+1}$ for $0 \le j \le K_{\max}-1$, $U_{\mathrm{hub}}(K_{\max}+1, K_{\max}) := \min(1, \frac{2}{K_{\max}+1})$, and recursively for $k = K_{\max}$ down to $2$ and $0 \le j < k$:}

\fitdisplay{%
U_{\mathrm{hub}}(k, j) := \min\left(c(k, j),\; \max_{g \in \{j+1,\, (k+j+1)/2\}} \frac{\frac{1}{k}\cdot 1 + \frac{j}{k} U_{\mathrm{hub}}(k, j-1) + g U_{\mathrm{hub}}(k+1, j+1)}{\frac{1+j}{k} + g}\right),
}

\emph{where $U_{\mathrm{hub}}(k, -1) := 0$ and $c(k, j) := \frac{1}{j+1}$ if $j \le k - 2$, while $c(k, k-1) := \min\left(1, \frac{2}{k}\right)$. Then for every $N \ge 4$, every $x \in [0, 1)$, and every leaf $v \in V(T_N)$ attached to a host $w$ of degree $k \ge 2$ with $j \in \{0, \dots, k-1\}$ other leaf neighbors,}

\fitdisplay{%
\mathbb{P}(Z_v > x \mid T_N) \le \mathbb{P}(Z_v > 0 \mid T_N) \le \mathbb{P}(v \text{ ever receives a child} \mid T_N) \le U_{\mathrm{hub}}(k, j).
}

\emph{Proof.} A leaf $v$ that never receives a child has $d_N(v) \equiv 1$, hence $Z_v = 0$. Track the state $(k, j)$ of host $w$ (degree $k$, $j$ other leaf neighbors besides $v$) until $v$ receives its first child (payoff $1$). Before $v$ receives a child, the only events changing $(k, j)$ or stopping are:
\begin{itemize}
\item attachment to $v$ (rate $1/k$, stops with payoff $1$);
\item attachment to one of the $j$ other leaves of $w$ (rate $j/k$, moves to $(k, j-1)$);
\item attachment to $w$ (rate $g = (j+1) + \sum_{\text{non-leaf } z \sim w} 1/d(z) \in [j+1, (k+j+1)/2]$, moves to $(k+1, j+1)$).
\end{itemize}

First, verify the analytic supermartingale caps $c(k, j)$:
\begin{enumerate}
\item \textbf{Permanent non-leaf invariant ($j \le k - 2$):} Since degrees never decrease, $k \ge j + 2$ persists on every transition. Test $f_1(j) := \frac{1}{j+1}$ with payoff $1$ at promotion. For $j = 0$, $f_1(0) = 1$ bounds all probabilities. For $j \ge 1$, the generator drift is

\fitdisplay{%
g\left(\frac{1}{j+2} - \frac{1}{j+1}\right) + \frac{j}{k}\left(\frac{1}{j} - \frac{1}{j+1}\right) + \frac{1}{k}\left(1 - \frac{1}{j+1}\right) \le -\frac{1}{j+2} + \frac{1}{k(j+1)} + \frac{j}{k(j+1)} = -\frac{1}{j+2} + \frac{1}{k} \le 0
}
\end{enumerate}

because $g \ge j + 1$ and $k \ge j + 2$. By optional stopping (the marked leaf $v$ itself triggers an event at rate $1/k > 0$, or host $w$ grows at rate $\ge 1$, so the chain exits every bounded region a.s.), $\mathbb{P}(v \text{ ever promoted}) \le \frac{1}{j+1}$.
\begin{enumerate}\setcounter{enumi}{1}
\item \textbf{General case including $j = k - 1$ (no non-leaf neighbor):} Test $f_2(j) := \min\left(1, \frac{2}{j+1}\right)$. For $j \in \{0, 1\}$, $f_2(j) = 1$. For $j \ge 2$, using $g \ge j + 1$ and $k \ge j + 1$, the drift is

\fitdisplay{\begin{multlined}[\linewidth]%
g\left(\frac{2}{j+2} - \frac{2}{j+1}\right) + \frac{j}{k}\left(\frac{2}{j} - \frac{2}{j+1}\right) + \frac{1}{k}\left(1 - \frac{2}{j+1}\right) \ftbreak \le -\frac{2}{j+2} + \frac{j + 1}{k(j+1)} = -\frac{2}{j+2} + \frac{1}{k} \ftbreak \le -\frac{2}{j+2} + \frac{1}{j+1} = -\frac{j}{(j+1)(j+2)} < 0.
\end{multlined}}
\end{enumerate}

Hence $\mathbb{P}(v \text{ ever promoted}) \le \min(1, \frac{2}{j+1})$ unconditionally for all $(k, j)$.

Now backward induction on $(K_{\max}+1 - k, j)$ proves $\mathbb{P}(v \text{ ever promoted} \mid k, j) \le U_{\mathrm{hub}}(k, j)$: the base $k = K_{\max}+1$ is $c(K_{\max}+1, j)$, and for $k \le K_{\max}$, conditioning on the first jump of $\{v, \text{other } j \text{ leaves}, w\}$ gives an affine-in-$g$ mixture whose maximum over $g \in [j+1, (k+j+1)/2]$ is achieved at the two endpoints, and taking the minimum with the independently valid bound $c(k, j)$ preserves validity. Note that states with $j \le k - 2$ never transition to $j' = k' - 1$. $\square$
\reviewnote{The first-jump mixture is fractional-linear, not affine, in $g$; clearing its positive denominator gives the same two endpoint tests. For hosts beyond the finite table the intended caps are $1/(j+1)$ on the branch $j\le k-2$ and $\min(1,2/(j+1))$ otherwise (Section~\ref{sec:ord-review}).}
\end{transcription}
% END TRANSCRIPTION round3/q3-birth-order/PROOF.md
% BEGIN TRANSCRIPTION round3/q3-birth-order/PROOF-PART-2.md lines 9-236 sha256 9a3a43c2534ea087fb23a735c1cde4975147bf87387b8bda2d16fbc00d8a2653 commit 62081cad01112c42c2011ac3eed11aa3aa665e09
\begin{transcription}
\subsubsection*{4. Exact Discrete-Time Continuation Bounds for All $m \ge m_0 = 200$ (\texttt{cert\_all\_m.cpp})}

To close all $m = n+2 \ge m_0 := 200$ simultaneously without any continuous-time Gamma/$S_m$ Chernoff union bound, we work \textbf{directly in discrete tree size $N \ge m$} on \textbf{raw survival probabilities $\mathbb{P}(Z_v > x \mid T_N)$}, parameterized by the dimensionless birth-scale threshold

\fitdisplay{%
\tau := (m - 1)\log\frac{1}{1 - x} \in [0, \infty) \qquad \Longleftrightarrow \qquad q := 1 - x = e^{-\tau/(m-1)}.
}

Fix a closed rational box $[\tau_{\mathrm{lo}}, \tau_{\mathrm{hi}}] \subset [0, 6]$, a uniform time step $h := 1/128$, and a finite horizon $T = M_{\mathrm{steps}} h$ with grid knots $t_i = i h$ ($i = 0, \dots, M_{\mathrm{steps}}$). Note that $m_0 - 1 = 199 > 128 = 1/h$.

\trhead{4.1 Exact Raw Endpoint Enclosures Uniform in $m \ge m_0$}

\textbf{Lemma 4.1 (Uniform Lower Enclosure of $H_{N, j}(x)$).}\\
\emph{Let $t_{\mathrm{lo}}(N) := \log(N/m)$. For any $m \ge m_0$, any $\tau \le \tau_{\mathrm{hi}}$, and any integer $N \ge m$ with $t_{\mathrm{lo}}(N) \le t_{\mathrm{hi}}$, set $a_+ := \frac{m_0 e^{t_{\mathrm{hi}}} - 1}{m_0 - 1}$. Then}

\fitdisplay{%
H_{N, j}(x) \ge e^{-1} e^{-\tau_{\mathrm{hi}}(a_+ + 1/(m_0 - 1))} \sum_{k=0}^{j-1} \frac{1}{k!} \sum_{i=0}^{j-k-1} \frac{\tau_{\mathrm{hi}}^i}{i!} \prod_{s=0}^{i-1} \left(a_+ - \frac{s}{m_0 - 1}\right)_+.
}

\emph{Proof.} Since $H_{N, j}(x)$ is decreasing in both $x$ and $N$, it suffices to bound $H_{N, j}(x)$ at $\tau = \tau_{\mathrm{hi}}$ and $N \le m e^{t_{\mathrm{hi}}}$. For $N \le m e^{t_{\mathrm{hi}}}$ and $m \ge m_0$,

\fitdisplay{%
\frac{N - 1}{m - 1} \le \frac{m e^{t_{\mathrm{hi}}} - 1}{m - 1} = e^{t_{\mathrm{hi}}} + \frac{e^{t_{\mathrm{hi}}} - 1}{m - 1} \le a_+.
}

Evaluating the exact binomial tail $B_{N, r}(x) = q^{N-1}\sum_{i=0}^{r-1}\binom{N-1}{i}(x/q)^i$ at the single upper envelope $N_* := 1 + (m-1)a_+ \ge N$ (or noting $B_{N, r}(x)$ is decreasing in $N$, and for $N = \lfloor m e^{t_{\mathrm{hi}}} \rfloor \le 1 + (m-1)a_+$, we have $N \le 1 + (m-1)a_+$ while $N + 1 > m e^{t_{\mathrm{hi}}}$? More directly: $B_{N, r}(x)$ is decreasing in real $N$, so evaluate at real $N_1 := \lceil 1 + (m-1)a_+ \rceil \ge N$, which satisfies $(N_1 - 1)/(m-1) \in [a_+, a_+ + \frac{1}{m-1}] \subset [a_+, a_+ + \frac{1}{m_0-1}]$!). At this integer $N_1 \ge N$:
\reviewnote{The parenthetical remark in this paragraph is reproduced as written. Its operative step: evaluate at the integer $N_1=\lceil 1+(m-1)a_+\rceil\ge N$ and use monotonicity in $N$.}
\begin{enumerate}
\item $q^{N_1 - 1} = \exp\big(-\tau_{\mathrm{hi}}\frac{N_1 - 1}{m - 1}\big) \ge \exp\big(-\tau_{\mathrm{hi}}(a_+ + \frac{1}{m_0 - 1})\big)$;
\item $\frac{x}{q} = e^{\tau_{\mathrm{hi}}/(m-1)} - 1 \ge \frac{\tau_{\mathrm{hi}}}{m - 1}$;
\item For each $0 \le s \le i-1$, $\frac{N_1 - 1 - s}{m - 1} \ge a_+ - \frac{s}{m - 1} \ge \left(a_+ - \frac{s}{m_0 - 1}\right)_+$.
Multiplying these inequalities yields the stated bound (\texttt{H\_raw\_lower\_int}). $\square$
\end{enumerate}

\textbf{Lemma 4.2 (Uniform Upper Enclosure of $B_{N, r}(x)$).}\\
\emph{Let $t_{\mathrm{up}}(N) := \log\frac{N-1}{m-1}$. For any $m \ge m_0$, any $\tau \ge \tau_{\mathrm{lo}} > 0$ with $u_0 := \frac{\tau_{\mathrm{lo}}}{m_0 - 1} \le \frac{1}{2}$, any $r \in (0, 3)$, and any integer $N \ge m$ with $t_{\mathrm{up}}(N) \ge t_{\mathrm{lo}}$, set $L_{\mathrm{lo}} := \tau_{\mathrm{lo}} e^{t_{\mathrm{lo}}}$ and $B_0 := (m_0 - 1)e^{t_{\mathrm{lo}}} + 1 - r$. Then substituting $1 - u = q e^{-v/B}$ ($B := N_{\mathrm{lo}} - r$ at real $N_{\mathrm{lo}} := 1 + (m-1)e^{t_{\mathrm{lo}}} \le N$) gives the exact identity}

\fitdisplay{\begin{multlined}[\linewidth]%
B_{N, r}(x) \le B_{N_{\mathrm{lo}}, r}(x_{\mathrm{lo}}) \ftbreak = e^{-L_{\mathrm{lo}}} \cdot \frac{\Gamma(B+r)}{\Gamma(B)B^r} \cdot \frac{1}{\Gamma(r)} \int_0^\infty \big(w + B(1 - e^{-v/B})\big)^{r-1} e^{-v}\,dv, \ftbreak \qquad w := B\frac{x_{\mathrm{lo}}}{1 - x_{\mathrm{lo}}},
\end{multlined}}

\emph{which is bounded above uniformly in $m \ge m_0$ by:}
\begin{enumerate}
\item \textbf{\emph{For $0 < r < 1$:}} \emph{$\frac{\Gamma(B+r)}{\Gamma(B)B^r} \le 1$ (Gautschi's inequality), $B \ge B_0$, and $w > L_{\mathrm{lo}}$. Splitting $[0, \infty)$ at $V_{\mathrm{cut}} := \min(24, \lfloor B_0/2 \rfloor)$ with $\alpha := 1 - \frac{V_{\mathrm{cut}}}{2 B_0}$ (using $B(1 - e^{-v/B}) \ge \alpha v$ on $[0, V_{\mathrm{cut}}]$ and $\ge \alpha V_{\mathrm{cut}}$ on $[V_{\mathrm{cut}}, \infty)$) gives}

\fitdisplay{%
B_{N, r}(x) \le e^{-L_{\mathrm{lo}}} \left[ \alpha^{r-1} e^{L_{\mathrm{lo}}/\alpha} Q\!\left(r, \frac{L_{\mathrm{lo}}}{\alpha}\right) + \frac{(L_{\mathrm{lo}} + \alpha V_{\mathrm{cut}})^{r-1} e^{-V_{\mathrm{cut}}}}{\Gamma(r)} \right].
}

\item \textbf{\emph{For $r = k + a \in [1, 3)$ ($k \in \{1, 2\}, a \in [0, 1)$):}} \emph{$B \le (m-1)e^{t_{\mathrm{lo}}}$, so $w \le w_{\mathrm{hi}} := L_{\mathrm{lo}}\frac{e^{u_0} - 1}{u_0}$, $B(1 - e^{-v/B}) \le v$, and $\frac{\Gamma(B+r)}{\Gamma(B)B^r} \le C_r(B_0)$ where $C_{1+a}(B_0) = 1 + \frac{a}{B_0}$ and $C_{2+a}(B_0) = (1 + \frac{1+a}{B_0})(1 + \frac{a}{B_0})$, giving}

\fitdisplay{%
B_{N, r}(x) \le e^{-L_{\mathrm{lo}}} C_r(B_0) \cdot e^{w_{\mathrm{hi}}} Q(r, w_{\mathrm{hi}}).
}
\end{enumerate}

\emph{Here $Q(r, v) := \frac{1}{\Gamma(r)}\int_v^\infty t^{r-1}e^{-t}\,dt$ is evaluated outward-rounded via odd-truncated alternating Taylor series for $v \le 2$ and the $(2K+1)$-st Stieltjes continued-fraction convergent ($K = 45$, guaranteed strict upper bound for $0 < a < 1$) for $v > 2$, extended to $r \in [1, 3)$ by the exact recurrence $e^v Q(1+a, v) = e^v Q(a, v) + v^a/\Gamma(1+a)$ (\texttt{B\_raw\_upper\_int}).}

\trhead{4.2 Exact Discrete-Time Affine Knot-Crossing Theorem}

The key observation that eliminates all time-discretization error for \textbf{all $m \ge m_0 = 200$} is that the discrete time increments satisfy the universal one-sided inequalities

\fitdisplay{%
N\Delta t_{\mathrm{lo}}(N) = N\log\!\left(1 + \frac{1}{N}\right) < 1, \qquad N\Delta t_{\mathrm{up}}(N) = N\log\!\left(1 + \frac{1}{N-1}\right) > 1 \qquad (\forall N \ge 2),
}

while both step sizes satisfy $\Delta t(N) < \frac{1}{m_0 - 1} = \frac{1}{199} < \frac{1}{128} = h$ so a single step $N \to N+1$ crosses \textbf{at most one grid knot $t_{i+1}$}.

\textbf{Lemma 4.3 (Affine Knot-Crossing Lemma).}\\
\emph{Let $f(t, s)$ be continuous and piecewise-linear on $[0, T + h]$ with nonpositive slope $-\sigma_i(s) \le 0$ on each slab $[t_i, t_{i+1}]$ (so $f(t, s) = f_{i+1}(s) + \sigma_i(s)(t_{i+1} - t)$ for $t \in [t_i, t_{i+1}]$), and consider a Markov transition at tree size $N$ from $(t, s)$ ($t \in [t_i, t_{i+1})$) to $(t', S_{N+1})$ with $t' = t + \Delta t$, $0 < \Delta t < h$, and transition probabilities $p_{s \to s'} = r_{s \to s'}/N$ ($s' \neq s$, $\sum_{s' \neq s} r_{s \to s'} = r_{\mathrm{tot}} \le N$).}
\begin{enumerate}
\item \textbf{\emph{(Lower submartingale, $N\Delta t < 1$):}} \emph{Suppose slab $i$ satisfies the two endpoint conditions}

\fitdisplay{%
\text{(L-R)}\quad \sigma_i(s) \le \sum_{s' \neq s} r_{s \to s'}\big(f_{i+1}(s') - f_{i+1}(s)\big), \qquad \text{(L-L)}\quad \sigma_i(s) \le \frac{\sum_{s' \neq s} r_{s \to s'} f_i(s') - r_{\mathrm{tot}} f_{i+1}(s)}{1 + h r_{\mathrm{tot}}}.
}
\end{enumerate}

\emph{Then $\mathbb{E}[f(t', S_{N+1})] - f(t, s) \ge 0$ for EVERY $t' \in [t_i, t_{i+1}]$, AND for every knot-crossing step $t' \in (t_{i+1}, t_{i+2}]$ provided either (a) $(i+1, s)$ is also a continuing state satisfying (L-R) and (L-L) on slab $i+1$, or (b) $(i+1, s)$ is a stopping state satisfying the explicit knot condition}

\fitdisplay{%
\text{(L-KNOT)}\quad -\sigma_{i+1}(s) + \sum_{s' \neq s} r_{s \to s'}\Big[f_{i+1}(s') - f_{i+1}(s) - \big(\sigma_{i+1}(s') - \sigma_{i+1}(s)\big)\delta_2\Big] \ge 0 \qquad (\delta_2 \in \{0, \Delta t_{\max}^{\mathrm{cor}}\}).
}

\begin{enumerate}\setcounter{enumi}{1}
\item \textbf{\emph{(Upper supermartingale, $N\Delta t > 1$):}} \emph{Suppose slab $i$ satisfies}

\fitdisplay{%
\text{(U-R)}\quad \sigma_i(s) \ge \sum_{s' \neq s} r_{s \to s'}\big(f_{i+1}(s') - f_{i+1}(s)\big), \qquad \text{(U-L)}\quad \sigma_i(s) \ge \frac{\sum_{s' \neq s} r_{s \to s'} f_i(s') - r_{\mathrm{tot}} f_{i+1}(s)}{1 + h r_{\mathrm{tot}}}.
}
\end{enumerate}

\emph{Then $\mathbb{E}[f(t', S_{N+1})] - f(t, s) \le 0$ for EVERY $t' \in [t_i, t_{i+1}]$, AND for every knot-crossing step $t' \in (t_{i+1}, t_{i+2}]$ whenever $(i+1, s)$ is a continuing state satisfying (U-R) and (U-L) on slab $i+1$. At the terminal horizon $t_M = T$, lifting $f_M^-(s) := f_{\mathrm{stop}}(s) + \Delta t_{\max}(M-1)\big(\sum_{s' \neq s} r_{s \to s'}(f_{\mathrm{stop}}(s') - f_{\mathrm{stop}}(s))\big)_+$ guarantees $\mathbb{E}[f_{\mathrm{stop}}(S_{N+1})] \le f(t, s)$ for every step crossing $t_M$.}

\emph{Proof.} Write the exact one-step drift $N\big(\mathbb{E}[f(t', S_{N+1})] - f(t, s)\big) = N\big(f(t', s) - f(t, s)\big) + \sum_{s' \neq s} r_{s \to s'}\big(f(t', s') - f(t', s)\big)$.
\begin{itemize}
\item \textbf{Interior ($t' \in [t_i, t_{i+1}]$):} $N(f(t', s) - f(t, s)) = -(N\Delta t)\sigma_i(s)$ (which is $\ge -\sigma_i(s)$ if $N\Delta t < 1$, and $\le -\sigma_i(s)$ if $N\Delta t > 1$), while $\sum_{s' \neq s} r_{s \to s'}(f(t', s') - f(t', s))$ is affine in $t' \in [t_i, t_{i+1}]$, so its extrema occur at $t' = t_{i+1}$ (giving (L-R)/(U-R)) and $t' = t_i$ (using $f_i(s) = f_{i+1}(s) + h\sigma_i(s)$, giving (L-L)/(U-L)).
\item \textbf{Knot crossing ($t' \in (t_{i+1}, t_{i+2}]$):} Write $\delta_1 = t_{i+1} - t \ge 0$, $\delta_2 = t' - t_{i+1} > 0$, $\delta_1 + \delta_2 = \Delta t$. Then

\fitdisplay{%
N\big(\mathbb{E}[f(t', S_{N+1})] - f(t, s)\big) = -N\sigma_i(s)(\Delta t - \delta_2) - N\sigma_{i+1}(s)\delta_2 + \sum_{s' \neq s} r_{s \to s'}\big(f(t_{i+1} + \delta_2, s') - f(t_{i+1} + \delta_2, s)\big),
}
\end{itemize}

which is \textbf{affine in $\delta_2 \in [0, \Delta t]$}! At $\delta_2 = 0$, it reduces to the $t' = t_{i+1}$ condition on slab $i$ ((L-R)/(U-R)). At $\delta_2 = \Delta t$, it equals the interior drift on slab $i+1$ at $t' = t_{i+1} + \Delta t \in [t_{i+1}, t_{i+2}]$ (already certified if $(i+1, s)$ is continuing, or directly checked by (L-KNOT) if $(i+1, s)$ is a lower stopping state).
\begin{itemize}
\item \textbf{Terminal lift ($t < t_M \le t'$ in the upper bound):} For $t \in [t_{M-1}, t_M)$, $f(t, s) \ge f_M^-(s)$ and $1/N < 1/(N-1) \le \Delta t_{\max}(M-1)$, so $\mathbb{E}[f_{\mathrm{stop}}(S_{N+1})] = f_{\mathrm{stop}}(s) + \frac{1}{N}\sum_{s' \neq s} r_{s \to s'}(f_{\mathrm{stop}}(s') - f_{\mathrm{stop}}(s)) \le f_M^-(s) \le f(t, s)$. $\square$
\end{itemize}

\trhead{4.3 Tracked Markov States (Exact DAG Ordering)}
Both lower and upper state graphs are acyclic within each slab $i$, enabling exact computation of $f_i(s')$ before $f_i(s)$:
\begin{enumerate}
\item \textbf{Lower submartingale} (at $\tau_{\mathrm{hi}}$, returning $M_{\mathrm{lo}} := \max(V_0(2,1), H_{\mathrm{raw\_lower}}(1, 0))$, $H_{\mathrm{host}}(p) := \max(H_0(p), M_{\mathrm{lo}})$, $L_{\mathrm{lo}}(k) := R_0(1, k)$):
\begin{itemize}
\item \textbf{Tracked positive stock $V(d, j)$} ($2 \le d \le D_{\max}, 1 \le j < d$): degree $d$, $j$ tracked leaves; gain rate $j \to V(d+1, j+1)$, loss rate $j/d \to V(d, j-1)$ (or $R(d, 2)$ if $j = 1$); immediate-stop option $H_{\mathrm{raw\_lower}}(j, t_{i+1})$ (mandatory at $d = D_{\max}$ and whenever (L-KNOT) fails).
\item \textbf{Zero-stock return $R(d, k)$} ($1 \le d < D_{\max}, 2 \le k \le R_{\max}$): degree $d$, $0$ leaves, one tracked neighbor of degree $k$; return rate $1/k \to V(d+1, 1)$, neighbor growth rate $\le k - 1 + 1/d \to R(d, k+1)$ (zero at $k = R_{\max}+1$).
\item \textbf{Serial host bonus $H(p)$} ($2 \le p \le H_{\mathrm{cut}}$, handing off to $V(2, 1)$ at $p = H_{\mathrm{cut}}+1$): degree-$2$ serial middle with $1$ leaf and $1$ host neighbor of degree $p$; gain rate $1 + 1/p \to V(3, 2)$, loss rate $1/2 \to R(2, 2)$, host growth rate $r_{\mathrm{host}} \in [0, p - 1/2] \to H(p+1)$ (both endpoints checked).
\item \textbf{Pathwise monotonicity:} At every knot $i$, \texttt{cert\_all\_m.cpp} explicitly asserts $V_i(d+1, j+1) \ge V_i(d, j)$, $R_i(d, k) \ge R_i(d, k+1)$, $V_i(d+1, 1) \ge R_i(d, 2)$, and $H_i(p) \ge H_i(p+1)$. By piecewise linearity, these hold for all real $t \in [0, T+h]$, justifying dropping extra gains and inflating competitor growth rates.
\end{itemize}
\item \textbf{Upper supermartingale} (at $\tau_{\mathrm{lo}} > 0$; at $\tau_{\mathrm{lo}} = 0$ uses Lemma 3.3 \texttt{U\_hub}):
\begin{itemize}
\item \textbf{Post-promotion 1-race $P(k)$} ($2 \le k \le K_{\max}$; stopped at $k = K_{\max}+1$ using $t_{\max(0, i-1)}$ so linear interpolation on $[t_i, t_{i+1}]$ stays $\ge B_{\mathrm{raw\_upper}}(1 + 1/(K_{\max}+1), t_i)$): degree-$2$ promoted vertex with $1$ leaf and $1$ host of degree $k$; gain rate $1 + 1/k \to$ stop with $B_{\mathrm{raw\_upper}}(2 + 1/k, t_i)$, leaf-loss rate $1/2 \to$ stop with $B_{\mathrm{raw\_upper}}(1/2 + 1/k, t_i)$, host-growth rate $g_w \in \{1/2, k - 1/2\} \to P(k+1)$ (both endpoints checked; monotone lift $P_i(k) \ge P_i(k+1)$ enforced).
\item \textbf{Pending adverse leaf $U(k, j)$} ($2 \le k \le K_{\max}, 0 \le j < k$; stopped at $k = K_{\max}+1$ with $\min(B_{\mathrm{raw\_upper}}(1/(K_{\max}+1), t_{\max(0,i-1)}), U_{\mathrm{hub}})$): promotion rate $1/k \to P(k)$, sibling loss rate $j/k \to U(k, j-1)$, host gain rate $g \in \{j+1, (k+j+1)/2\} \to U(k+1, j+1)$ (both endpoints checked). At $t = 0$, capped by $\min(U_0(k, j), B_{\mathrm{raw\_upper}}(1/k, 0), U_{\mathrm{hub}}(k, j))$.
\end{itemize}
\end{enumerate}


\subsubsection*{5. Exact Integer-$N$ Backward Induction for $4 \le m \le 199$ (\texttt{cert\_small\_m.cpp})}

For every finite $m = b + 1 \in \{4, \dots, 199\}$, parameterize $x \in [0, 1)$ globally by $y := \log\frac{1}{1 - x} \in [0, \infty)$. By Section 7 (High-$\tau$ Tail), every pair $(m, y)$ with $(m - 1)y \ge 6$ is already certified, so only $y \in [0, 2]$ ($m_{\max}(y) := \min(199, \lfloor 6/y \rfloor + 1)$) remains.

On each exact rational box $[y_{\mathrm{lo}}, y_{\mathrm{hi}}] \subset [0, 2]$:
\begin{itemize}
\item If $199 y_{\mathrm{hi}} \le 6$: initialize all state arrays $(V, R, H, U, P)$ at $N_{\mathrm{start}} = 200$ using Section 4's exact \texttt{init\_at\_200} solver on $[199 y_{\mathrm{lo}}, 199 y_{\mathrm{hi}}]$.
\item If $199 y_{\mathrm{lo}} \ge 6$: initialize directly at $N_{\mathrm{start}} := \max(m_{\max} + 4, \min(900, \lceil 24/y_{\mathrm{lo}} \rceil + 2))$ using exact fixed-$N$ \texttt{H\_exact\_N\_lower} and \texttt{B\_exact\_N\_upper}.
\item Then perform \textbf{one literal exact integer-$N$ backward sweep} $N = N_{\mathrm{start}} - 1, \dots, 4$:
\begin{itemize}
\item Every 1-step transition probability at tree size $N$ is an exact rational with denominator $N d$, $N k d$, or $2 N k$, rounded \textbf{down} for $(V, R, H)$ (taking $\max$ with exact $H_{N, j}(x_{\mathrm{hi}})$) and \textbf{up} for $(P, U)$ (taking $\min$ with exact $B_{N, r}(x_{\mathrm{lo}})$ and $U_{\mathrm{hub}}$).
\item For $4 \le N < 60$, \texttt{B\_exact\_N\_upper} computes $\frac{\Gamma(N)}{\Gamma(r)\Gamma(N-r)B^r}$ in outward-rounded log-space via \texttt{iv\_gamma\_rat} and bounds $\int_0^\infty (w + B(1-e^{-v/B}))^{r-1} e^{-v}\,dv$ by exact outward-rounded monotone quadrature (decreasing integrand for $r < 1$ with exact analytic antiderivative on $[0, 1/64]$; increasing integrand bounded by $(w+B)^{r-1}$ for $r > 1$).
\item At every size $N \in \{4, \dots, m_{\max}\}$, the exact snapshot $(V, R, H, U, P)$ at size $N$ is checked against the exact finite-tree signed ledger for $m = N$ (maximum vertex degree $d \le m - 3$).
\end{itemize}
\end{itemize}


\subsubsection*{6. Dimensionless Normalized All-Tree Signed Ledger}

Fix any box and normalize all lower bounds $(H_{\mathrm{host}}(p), L_{\mathrm{lo}}(k))$ downward and all upper bounds $U_{\mathrm{hi}}(k, j)$ upward by dividing by $M_{\mathrm{lo}} > 0$ (in integer scale $L_{\mathrm{SC}} = 10^9$), so $\widetilde{M} = 1$, $\widetilde{H}(p) \ge 1$, $g_{\mathrm{base}} := 1 - \widetilde{U}(2, 0) > 0$, and for $k > K_{\max}+1$, $\widetilde{U}(k, j) := \min\big(\frac{1}{(j+1)M_{\mathrm{lo}}}, \widetilde{u}_{\mathrm{tail0}}\big)$.

\trhead{6.1 Tree-Potential Charging Identity}
Let $T_n$ be any finite tree of size $n \ge 2$ ($m = n+2$).
\begin{itemize}
\item \textbf{Seed ($n = 2$, $T_2$ = edge):} $s(1) = s(2) = 1/2$, no edge debit; Lemma 2.1 reduces to $\widetilde{H}(2) > \widetilde{U}(2, 0)$.
\item \textbf{Stars ($n = d + 1$, $S_{d+1}$):} center has degree $d \ge 2$, $c = 0$, $l = d$, $r = 0$, $s = \frac{1}{d+1}$, leaf weight $t = \frac{1}{2d}$; Lemma 2.1 reduces to $C_{\mathrm{lo}}(d, 0, 0) > 0$ (checked for all $d \le D_{\mathrm{tail}} = 300$, and for $d > 300$ by $\frac{14}{10} g_{\mathrm{base}} > \frac{1}{2}\widetilde{U}(302, 300)$, since $\frac{d}{d+1} + \frac{1}{2} > 1.49 > 1.4$ and $\frac{d-1}{2(d+1)} < \frac{1}{2}$).
\item \textbf{Multi-core trees (not an edge or star):} Deleting all leaves yields the core tree $C(T_n)$ with $|V(C(T_n))| \ge 2$. Every vertex $u \in C(T_n)$ has $d(u) \ge 2$, core degree $c(u) \ge 1$, leaf count $l(u) = d(u) - c(u)$, and receive sum $r(u) := \sum_{v \sim u,\, v \in C(T_n)} 1/d(v) \in (0, c(u)/2]$. Contracting maximal paths of bare degree-$2$ vertices ($d = c = 2, l = 0$) into chain-edges yields a site tree $\mathcal{T}_S$ on $|S| \ge 2$ sites with $|E(\mathcal{T}_S)| = |S| - 1$ edges and $\sum_{u \in S}(c(u) - 2) = -2$.
\end{itemize}

Choose nonnegative tariffs $(a, b_{\mathrm{base}}, \mathrm{fee}_{\mathrm{lo}}, \mathrm{fee}_{\mathrm{hi}}, \mathrm{fee}_t, \mathrm{xfer}, D_{\mathrm{hi}})$. Define:
\begin{enumerate}
\item \textbf{Site credit (after paying all incident leaf edges):}

\fitdisplay{\begin{multlined}[\linewidth]%
C_{\mathrm{lo}}(d, c, r) := d s(\widetilde{H}(d+1) - \widetilde{U}(2,0)) + \frac{l}{2d}(\widetilde{H}(2) - \widetilde{U}(2,0)) \ftbreak - l(t - s)_+(\widetilde{U}(2,0) - \widetilde{L}(d+1))_+ - l(s - t)_+(\widetilde{U}(d+1, l-1) - \widetilde{L}(2))_+,
\end{multlined}}
\end{enumerate}

where $s = \frac{l+r}{d(d+1)}$ and $t = \frac{1}{2d}$.
\begin{enumerate}\setcounter{enumi}{1}
\item \textbf{Direct site-site edges $\{u, v\}$:} At most one of $s(u) - s(v)$ and $s(v) - s(u)$ is positive, so the edge debit is bounded by $b_{\mathrm{base}} + (J_{u \to v} - b_{\mathrm{base}})_+ + (J_{v \to u} - b_{\mathrm{base}})_+$. Each direct edge between a high site ($d \ge D_{\mathrm{hi}}$) and a low site ($e \le 4$) transfers $\mathrm{xfer}$ from the high site to the low site (net sum $0$).
\item \textbf{Contracted chains:} Each chain port at site $u$ pays fee $F_{\mathrm{ch}}(u)$ ($\mathrm{fee}_t$ if $d=2$, $\mathrm{fee}_{\mathrm{lo}}$ if $3 \le d < D_{\mathrm{hi}}$, $\mathrm{fee}_{\mathrm{hi}}$ if $d \ge D_{\mathrm{hi}}$) from $u$ to the chain (net sum $0$). For every bare degree-$2$ chain $v_1 - \dots - v_L$ ($L \ge 1$) between sites $x, y$, its net debt (sum of internal edge debits + port debits minus internal serial credits minus $F_{\mathrm{ch}}(x) + F_{\mathrm{ch}}(y)$) is verified to be strictly $< b_{\mathrm{base}}$ for $L \in \{1, 2, 3\}$. For $L \ge 4$, every extra interior vertex $v_i$ ($3 \le i \le L-2$) has $s(v_i) = 1/6$ equal to both neighbors (zero extra edge debit) and positive serial credit $\frac{1}{3}(\widetilde{H}(3) - \widetilde{U}(2, 0)) > 0$, so $\mathrm{debt}(L) \le \mathrm{debt}(3) < b_{\mathrm{base}}$.
\item \textbf{Site balance:} Every site $u \in S$ satisfies $R(u) := C_{\mathrm{lo}}(u) + a(c(u) - 2) - b_{\mathrm{base}} - \sum_{\text{chains}} F_{\mathrm{ch}}(u) \pm \mathrm{xfer} \ge \mathrm{exc}(u) := \sum_{\text{direct } v} (J_{u \to v} - b_{\mathrm{base}})_+$.
\end{enumerate}

Summing over all sites $u \in S$ and all edges of $\mathcal{T}_S$ (using $\sum_{u \in S} (a(c(u) - 2) - b_{\mathrm{base}}) = -2a - (|S|-1)b_{\mathrm{base}} - b_{\mathrm{base}}$) yields:

\fitdisplay{\begin{multlined}[\linewidth]%
\frac{n(n+1)}{M_{\mathrm{lo}}}\Big(\mathbb{P}(Z_b > x \mid T_n) - \mathbb{P}(Z_m > x \mid T_n)\Big) \ftbreak \ge \sum_{u \in S}\big(R(u) - \mathrm{exc}(u)\big) + \sum_{\text{chains } P}\big(b_{\mathrm{base}} - \mathrm{debt}(P)\big) \ftbreak + (2a + b_{\mathrm{base}}) \ge 2a + b_{\mathrm{base}} > 0.
\end{multlined}}

\trhead{6.2 Completeness of Finite Checks \& Convexity/Concavity Reductions}
\begin{itemize}
\item \textbf{Low sites ($2 \le d < D_{\mathrm{hi}} = 10$):} For each $(d, c)$ and chain count $k \in \{0, \dots, c\}$, \texttt{R(r) - exc(r)} is concave in $r$ (since $-(t-s)_+$, $-(s-t)_+$, and $-(J_{u \to v} - b_{\mathrm{base}})_+$ are all concave in $s$), so checking the two extreme values of $r$ suffices. If the coarse check fails, exact multisets of direct neighbor degrees in $\{2, \dots, D_{\mathrm{hi}}-1\} \cup \{\ge D_{\mathrm{hi}}\}$ are enumerated with exact receive denominator $\mathrm{lcm}(2, \dots, 10) = 2520$, using opposite coefficient lower bound $s(v) \ge \frac{1}{d e(e+1)}$ and favorable payout $\widetilde{L}(e+1)$.
\item \textbf{High finite sites ($D_{\mathrm{hi}} \le d \le D_{\mathrm{tail}} = 300$):} If $s_{\max}\widetilde{U}(d+1, l) \le b_{\mathrm{base}}$, then $\mathrm{exc} \equiv 0$, so $R$ is jointly concave in $(k, p_{\mathrm{pay}}, r)$ and the 3 simplex vertices $\{(0,0), (c,0), (0,c)\} \times \{r_{\min}, r_{\max}\}$ are exact; otherwise we iterate over all integer $k \in \{0, \dots, c\}$, for each of which $R - \mathrm{exc}$ is jointly concave in $(p_{\mathrm{pay}}, r)$, checking $\{0, c-k\} \times \{r_{\min}, r_{\max}\}$.
\item \textbf{Analytic high-degree site tail ($d > D_{\mathrm{tail}} = 300$ in \texttt{cert\_all\_m.cpp}):}
\begin{enumerate}
\item $s\widetilde{U}(d+1, l) \le \frac{\widetilde{U}(302, 0)}{302} \le b_{\mathrm{base}}$, so $\mathrm{exc} \equiv 0$.
\item Since $s - t \le \frac{l-1}{2d(d+1)}$, core-side leaf debit is at most $T_{\mathrm{cl}} := \min\big(\frac{1}{2(D_{\mathrm{tail}}+2)M_{\mathrm{lo}}}, \frac{\widetilde{U}(302,0)}{2}\big)$.
\item Require $a \ge \max(\mathrm{fee}_{\mathrm{hi}}, \mathrm{xfer})$. If $l \le d/2$ ($c \ge D_{\mathrm{tail}}/2$), serial credit minus leaf-side debit is $\ge \frac{l}{d} z_{\mathrm{lo}} \ge \min(0, z_{\mathrm{lo}}/2)$ with $z_{\mathrm{lo}} = \frac{D_{\mathrm{tail}}}{D_{\mathrm{tail}}+1}g_{\mathrm{base}} + \frac{g_{\mathrm{base}} - \widetilde{U}(2,0)}{2}$, and net tariff adds $\ge (a - \max(\mathrm{fee}_{\mathrm{hi}}, \mathrm{xfer}))\lfloor D_{\mathrm{tail}}/2 \rfloor - 2a - b_{\mathrm{base}}$. If $l > d/2$, then $2l \ge d+1$ so $(t-s)_+ = 0$, while serial credit $\ge (\frac{D_{\mathrm{tail}}}{2(D_{\mathrm{tail}}+1)} + \frac{1}{4})g_{\mathrm{base}}$ pays $T_{\mathrm{cl}} + 2a + b_{\mathrm{base}}$ (we even subtract an extra conservative $\frac{\widetilde{U}(2,0)}{2(D_{\mathrm{tail}}+2)}$).
\end{enumerate}
\item \textbf{Bare degree-$2$ chains ($L \in \{1, 2, 3\}$):} Port debt is convex in $s_x$ (so extreme receives $2r \in \{1, c\}$ suffice) and convex in $(q_x, q_y)$ on $q \in [0, \frac{1}{6(K_{\max}+1)}]$ for the high tail $d > K_{\max}$ (so the two convex corners $q \in \{0, \frac{1}{6(K_{\max}+1)}\}$ with exact stock-coupled \texttt{own\_cap} $\max\big(\frac{\widetilde{u}_{\mathrm{tail0}}}{2(K_{\max}+2)}, \min(\frac{1}{4(K_{\max}+1)M_{\mathrm{lo}}}, \frac{\widetilde{u}_{\mathrm{tail0}}}{K_{\max}+2})\big)$ suffice).
\item \textbf{One-line universal shortcut:} Letting $A := \sum_{u \in V(T_n)} d(u) s(u) > 0$, we have $\sum_{\{u,v\}} |s(u) - s(v)| \le A$ and, for every non-leaf $v$ with $l \ge 1$ leaves, $l(t - s(v))_+ \le \frac{l(d-1)}{2d(d+1)} < \frac{1}{3}\left(\frac{l}{d+1} + \frac{l}{2d}\right) \le \frac{1}{3}(d(v)s(v) + \frac{l}{2d})$, so leaf-originating parallel weight is at most $A/3$. Hence if $\widetilde{U}(2,0) \ge U_3^{\max} := \max_{k \ge 3, 0 \le j < k} \widetilde{U}(k, j)$ and $3 > 4\widetilde{U}(2,0) + 2 U_3^{\max}$, then $G(T_n) \ge A M_{\mathrm{lo}}(1 - \frac{4}{3}\widetilde{U}(2,0) - \frac{2}{3}U_3^{\max}) > 0$ on every finite tree $T_n$.
\end{itemize}


\subsubsection*{7. Self-Contained Analytic High-$\tau$ Tail ($\tau \ge 6$, All $b \ge 3$)}

\textbf{Theorem 7.1.} \emph{For every $b = n + 1 \ge 3$ ($m = b+1 \ge 4$), every tree $T_n$, and every $x \in (0, 1)$ with $\tau := b\log\frac{1}{1-x} \ge 6$,}

\fitdisplay{%
\mathbb{P}(Z_b > x \mid T_n) - \mathbb{P}(Z_{b+1} > x \mid T_n) > \frac{19}{10000}\frac{(1-x)^b}{n(n+1)} > 0.
}

\emph{Proof.} Set $q := 1 - x$. Since $b x/q = b(e^{\tau/b} - 1) \ge \tau \ge 6$, every serial middle has $\mathbb{P}(Z_b > x) \ge H_{m, 1}(x) = e^{-1} q^b$ (Lemma 3.1), while adverse tips on hosts of degree $\ge 2$ and $\ge 3$ satisfy $\mathbb{P}(Z_m > x) \le B_{m, 1/2}(x)$ and $\le B_{m, 1/3}(x)$ respectively (Lemma 3.2, monotonic in shape). By the one-line inequality of Section 6.2 with $A = \sum_u w(u)/(d(u)+1) \ge \frac{1}{n}\sum_u w(u) = 1$,

\fitdisplay{%
n(n+1)\big(\mathbb{P}(Z_b > x \mid T_n) - \mathbb{P}(Z_{b+1} > x \mid T_n)\big) \ge q^b\left(e^{-1} - \frac{4}{3}\frac{B_{m, 1/2}(x)}{q^b} - \frac{2}{3}\frac{B_{m, 1/3}(x)}{q^b}\right).
}

For $r \in \{1/2, 1/3\}$, put $s := 1 - r$, $B := m - r = b + s \ge 7/2$, and $v := q/(B x) \le q/(b x) \le 1/6$. Substituting $u = x + q t$ in $B_{m, r}(x)$ and using Gautschi's bound $\frac{\Gamma(B+r)}{\Gamma(B+1)} \le B^{-s}$ gives

\fitdisplay{%
\frac{B_{m, r}(x)}{q^b} \le \frac{v^s}{\Gamma(r)} \mathbb{E}\!\left[\left(1 + \frac{q}{x} U\right)^{-s}\right], \qquad U \sim \mathrm{Beta}(1, B).
}

Since $(1+z)^{-s} \le 1 - s z + \frac{s(s+1)}{2} z^2$ for $z \ge 0$, $\mathbb{E}[U] = \frac{1}{B+1}$, $\mathbb{E}[U^2] = \frac{2}{(B+1)(B+2)}$, and $\frac{B}{B+1} \ge \frac{7}{9}$ (as $B \ge 7/2$),

\fitdisplay{%
\mathbb{E}\!\left[\left(1 + \frac{q}{x} U\right)^{-s}\right] \le 1 - s v \frac{B}{B+1}\left(1 - (1+s)v\frac{B}{B+2}\right) \le 1 - d_s v, \qquad d_s := \frac{7}{9} s\left(1 - \frac{1+s}{6}\right).
}

On $v \in (0, 1/6]$, $v^s(1 - d_s v)$ is strictly increasing (its derivative is $v^{s-1}(s - (s+1)d_s v) > 0$), so it is bounded above by its value at $v = 1/6$. Exact outward-rounded \texttt{\_\_int128} verification in \texttt{cert\_all\_m.cpp} (\texttt{verify\_high\_tau\_tail\_exact()}) proves

\fitdisplay{%
e^{-1} - \frac{4}{3}\frac{6^{-1/2}}{\Gamma(1/2)}\left(1 - \frac{d_{1/2}}{6}\right) - \frac{2}{3}\frac{6^{-2/3}}{\Gamma(1/3)}\left(1 - \frac{d_{2/3}}{6}\right) > 0.0050397 > \frac{19}{10000}. \quad \square
}


\subsubsection*{8. Self-Contained Founder Seam ($b = 2$: $Z_2 \succ_{\mathrm{st}} Z_3$)}

\textbf{Theorem 8.1.} \emph{For every $x \in [0, 1)$, $\mathbb{P}(Z_2 > x) > \mathbb{P}(Z_3 > x)$.}

\emph{Proof.} At $N = 3$, vertex $3$ attaches uniformly to $1$ or $2$. Conditioned on $3 \sim 1$ (probability $1/2$), $T_3$ is the path $2 - 1 - 3$, in which $2$ and $3$ are symmetric leaves, contributing $0$. Conditioned on $3 \sim 2$ (probability $1/2$), $T_3$ is the path $1 - 2 - 3$. Expanding one more step $T_3 \to T_4$:
\begin{itemize}
\item $U_3 = 1$ or $U_3 = 3$ (total probability $2/3$): vertex $4$ attaches to center $2$, giving the star $S_4$ centered at $2$ with leaves $\{1, 3, 4\}$. Center $2$ has degree $3$ and $3$ leaves (tracking any $2$ of them gives $\mathbb{P}(Z_2 > x \mid S_4) \ge V_4(3, 2)$ and also $\ge H_{4, 3}(x) > H_{4, 1}(x)$), while leaf $3$ is in state $(k=3, j=2)$ (so $\mathbb{P}(Z_3 > x \mid S_4) \le U_4(3, 2)$ and $\le B_{4, 1/3}(x)$).
\item $U_3 = 2, W_3 = 1$ (probability $1/6$): $T_4 = (4 - 1 - 2 - 3)$, where $2$ is a serial middle of state $H_4(2)$ ($\ge H_{4, 1}(x)$) and $3$ is a serial tip of state $U_4(2, 0)$ ($\le B_{4, 1/2}(x)$).
\item $U_3 = 2, W_3 = 3$ (probability $1/6$): $T_4 = (1 - 2 - 3 - 4)$, where $2$ and $3$ are symmetric centers, contributing $0$.
\end{itemize}

Hence for every $x \in [0, 1)$,

\fitdisplay{%
\mathbb{P}(Z_2 > x) - \mathbb{P}(Z_3 > x) = \frac{1}{3}\Big(\mathbb{P}(Z_{\mathrm{center}(S_4)} > x) - \mathbb{P}(Z_{\mathrm{leaf}(S_4)} > x)\Big) + \frac{1}{12}\Big(\mathbb{P}(Z_{\mathrm{mid}(P_4)} > x) - \mathbb{P}(Z_{\mathrm{leaf}(P_4)} > x)\Big).
}

\begin{itemize}
\item For $y = \log\frac{1}{1-x} \in [0, 2]$: \texttt{cert\_small\_m.cpp} certifies BOTH $V_4(3, 2) > U_4(3, 2)$ (using the rigorous $2/(j+1)$ cap at $j = k - 1$ from Lemma 3.3) AND $H_4(2) > U_4(2, 0)$ across all 84 exact contiguous boxes covering $[0, 2]$.
\item For $y \ge 2$: $\tau = 3 y \ge 6$, so Section 7 already proves $H_{4, 1}(x) = e^{-1} q^3 > \frac{4}{3} B_{4, 1/2}(x) > B_{4, 1/2}(x) > B_{4, 1/3}(x)$, making both terms strictly positive. $\square$
\end{itemize}


\subsubsection*{9. Exact Outward-Rounded \texttt{\_\_int128} Arithmetic Guarantees}

Every inequality in \texttt{cert\_all\_m.cpp} and \texttt{cert\_small\_m.cpp} is evaluated in exact 128-bit signed integer arithmetic (\texttt{\_\_int128}) with explicit directed rounding (\texttt{div\_fl} / \texttt{div\_cl}) and \textbf{zero floating-point trust}:
\begin{itemize}
\item Internal transcendental scale $S_2 = 10^{15}$, grid probability scale $S = 10^{12}$, normalized ledger scale $L_{\mathrm{SC}} = 10^9$, chain-coordinate scale $X = 6 \times 10^8$.
\item \texttt{iv\_exp\_pos(x)} ($0 \le x \le 16$, asserted): halving to $u \le 1/8$ + 15-term Taylor series + exact geometric remainder $2 u^{16}/16!$.
\item \texttt{iv\_exp\_neg(x)} (any $x \ge 0$): chunked by $4$ via $e^{-4}$ (overflow-free for all $x \ge 0$), bounded by $[0, 1/S_2]$ for $x \ge 42$ (or $45$).
\item \texttt{iv\_log\_ge1\_raw(y)} ($y \ge 1$): power-of-2 reduction to $[1, 2)$ + 20-term arctanh series in $w = (z-1)/(z+1) \in [0, 1/3]$ + exact remainder $\frac{9}{8} w^{41}/41$.
\item \texttt{iv\_gamma\_rat(p, q)} ($0 < p < q$): shift to $x = 60 + p/q$ using Binet's elementary Bernoulli bounds $\frac{1}{12x} - \frac{1}{360x^3} < \log\Gamma(x) - (x-\frac{1}{2})\log x + x - \frac{1}{2}\log(2\pi) < \frac{1}{12x}$ and outward-rounded log-space reduction.
\item Contiguous interval coverage is explicitly asserted inside both binaries (\texttt{covered\_up\_to == 6 * Q} for $\tau \in [0, 6]$ and \texttt{covered\_up\_to == 2 * YQ} for $y \in [0, 2]$).
\end{itemize}
\end{transcription}
% END TRANSCRIPTION round3/q3-birth-order/PROOF-PART-2.md

\subsection{Sources and replay}\label{sec:ord-sources}
\paragraph{Sources and replay.}
The reviews apply to the following exact bytes; both certificates are exported byte-identical.
\begin{center}\footnotesize
\begin{tabular}{@{}ll@{}}
\toprule
Item & Identifier\\
\midrule
Source commit & \texttt{62081cad01112c42c2011ac3eed11aa3aa665e09}\\
Frozen proof (SHA-256) & \texttt{e85e9cebfbde5439fad9246da2fec2c6eb7f014b09248e960c20d016411ce191}\\
\texttt{cert\_all\_m.cpp} (SHA-256) & \texttt{90689e44d08fdd94bb9edc76cae6d14669a5ac39e84d4da35d19a82d8a517c80}\\
\texttt{cert\_small\_m.cpp} (SHA-256) & \texttt{a10dce484242acb727973baf48c1bdac448dce9e761c506f55326a66b58c3210}\\
\bottomrule
\end{tabular}
\end{center}
From the repository root, \texttt{bash verify/compact\_replay.sh} checks both digests, compiles and runs both programs with \texttt{g++ -std=c++17}, and runs the \ordersrc{round3/q3-birth-order/audit_primitives.py}{primitive audit} and the \ordersrc{round3/q3-birth-order/planted_defects.py}{planted-defect controls} (Section~\ref{app:replay}). Passing runs print \texttt{Boxes PASSED: 52, FAILED: 0} and \texttt{Boxes PASSED: 84, FAILED: 0}. The small-size program counts 12{,}105 attempted box--size ledger checks; 11{,}307 of them lie on accepted boxes, and the other 798 precede subdivision. The \ordersrc{round3/PROVENANCE.md}{provenance manifest} separates byte-identical files from edited exports.

\sectag{EARLY}
\section{An independent finite certificate through label eight}\label{sec:early}
\tagnext{EARLY}
\begin{theorem}[Independent early-label certificate]\label{th:hub-early-order}
The inequalities $h_1=h_2>h_3>\cdots>h_8$ also follow from a finite size-nine calculation with rigorous continuation bounds, independently of the all-tree certificates.
\end{theorem}
\begin{proof}
The paragraphs below give the classifier, exact enumeration and signed enclosure. The founder tie there is exact by exchange of the two initial labels.
\end{proof}
The \proofsrc{round2/proofs/fable/newground/q3-hub-probability/REPORT.md}{early-label proof} used a signed rooted-history calculation. Here is the edge-theorem-free variant of the \proofsrc{round2/proofs/fable/newground/q3-hub-probability/lead-hub_order_exact.py}{independent rational checker}, with \texttt{NOEDGE=1} and argument \texttt{9}. For a core root with $l$ leaves and $c$ core neighbours, compare its stock below by rates $j,j/(j+c)$ until zero: at equality the true stock births are at least $j$ and its losses at most $j/(j+c)$ since the true core degree never decreases. The finite exit recursion has consecutive increments in ratio $1/(c+j)$, giving probability $\sum_{i=0}^{l-1}a_i/\sum_{i=0}^{H-1}a_i$ of reaching $H$ before zero, where $a_0=1$, $a_i=a_{i-1}/(c+i)$. The comparison and the leaf-hit estimate in the founder proof allow $H\to\infty$ to give the same ratio with infinite denominator as a hub lower. Enclose the denominator by keeping $i=0,\ldots,39$ and adding at most $a_{40}(c+41)/(c+40)$ by the geometric bound on later ratios. Other roots may use lower zero, without the edge theorem.

For a leaf root, the tracked parent stock $k$ includes the root; its core degree is at least the initial $c$. Use the stopped upper with terminal $F(41)=1/40$:
\[
 [k(k+c)+k]F(k)=k(k+c)F(k+1)+(k-1)F(k-1)+1 .
\]
The cutoff bound follows from $f(k)=1/\max(1,k-1)$ with success payoff one. For $k\ge3$ its birth drift is at most $-1/(k-1)$ and the combined promotion contribution is at most $1/(k-1)$: its numerator before division by parent degree is $1-1/(k-1)+1/(k-2)\le k/(k-1)$. For $k\le2$ the barrier is already one. The checker solves the stopped recurrence by tridiagonal elimination and verifies $F$ decreases in stock and lies in $[0,1]$. Extra births reduce the value and extra core neighbours slow the nonnegative aggregate promotion contribution, proving the physical upper. A core root with leaves can use upper one; at a bare root union-bound a first gain over all neighbours. For a specified core target use the same $f$ at sender stock $k$, now with $k$ unmarked losses and target selection at rate $1/D$, where $D\ge k+1$. For $k\ge3$ their combined term is $(k^2-3k+4)/[(k-1)(k-2)D]\le1/(k-1)$, since $k-2\ge1$ and $(k-2)(k+1)-(k^2-3k+4)=2(k-3)\ge0$. Track until that first send; bounded stopping proves each neighbour's contribution.

Enumerate all 40,320 histories through size nine, multiplying attached-site weights $r_u/n$ exactly. Canonicalize a marked rooted shape by recursively sorting the children's encodings (deleting the edge to the preceding parent); induction on order proves that equality of encodings is exactly rooted isomorphism. Aggregate the weights $q_i(s)$ of each of the 286 rooted shapes at label $i$. For signed coefficients $q_i(s)-q_{i+1}(s)$ multiply a nonnegative coefficient by the lower classifier and a negative one by the upper. This lower bounds the true gap, since the continuation from a common rooted shape has the same hub probability. Exact sums give positive lower bounds for $i=2,\ldots,7$, respectively greater than $2413,375,175,87,41,11$ in units of $10^{-4}$ (the last is $0.0011$). Founder masses tie exactly. The earlier finite computation and this separate classifier do not assert any later-label or positive-threshold inequality.

\part{Further material}

\sectag{FOR}
\section{A conditional obstruction for the low-degree forest}\label{sec:for}
This section proves an implication. Its count and log-volume premises are not proved in this paper or this supplement; Section~\ref{sec:further} records where they are argued and how far that argument was checked.
Root the original tree at label 1, delete vertices of original degree $>N$, and count each remaining component containing an original core once, including singleton surviving original cores. Root each at its highest core, and omit the exterior parent. The breakdown test uses only volumes, so the choice of initial root cannot affect it. Let $M_{n,N}$ be the count and $\nu_{n,N}$ the empirical probability law on rooted shapes. Write $V$ for ordinary vertex count, not core count. The reference $\varrho_N$ is the ordinary rooted-tree law of a newly promoted, planted vertex with one fresh leaf and an inert exterior parent slot, at an independent age uniform on $[0,\log N]$. The slot counts in the root degree but sends nothing and is omitted from the observed shape. Arrows into it produce no ordinary birth. This is a particular young-tree approximation, not a claim about a uniformly chosen surviving vertex.

\tagnext{FOR}
\begin{theorem}[Conditional breakdown of the young-tree approximation]\label{th:forest-tv-obstruction}
Let the deterministic integer cutoff $N=N(n)$ satisfy $N\to\infty$, $N/n\to0$, and put $\alpha=\liminf_n\log N/\log n$. As external inputs for this sequence, assume, with $\mu$ the core exponent and some random $K$ with $0<K<\infty$ almost surely, that almost surely
\begin{equation}\label{eq:forest-tv-input}
 M_{n,N}\sim K(n/N)^\mu\log N,\qquad
 \sup_{0\le x\le1}\left|\nu_{n,N}\!\left(\frac{\log V}{\log N}\le x\right)-x\right|\longrightarrow0.
\end{equation}
If $\alpha>\mu/(1+\mu)$ then almost surely
\[
 \liminf_n\|\nu_{n,N}-\varrho_N\|_{\mathrm{TV}}
 \ge 1-\mu(\alpha^{-1}-1)>0.
\]
The same lower bound holds for the laws projected onto volume. In particular both distances tend to one when $\alpha=1$, for example at $N=\lfloor n/\log n\rfloor$.
\end{theorem}
\begin{proof}
The count and uniform log-volume CDF are premises, not conclusions of this theorem. The separate \proofsrc{round2/proofs/neon2/q5/README.md}{simultaneous count and log-volume theorem} obtains them for counted components using the sequential forest law and the stated growth and spectral estimates; its review is narrower than that of this implication. In the reference process, at size $v$ the ordinary birth rate is $v-1/D_o\ge v-1/2$, since only the root-to-slot arrow is omitted and the root degree is at least two. The size increases by one at every intrinsic jump, has no explosion by Yule domination, and visits each size $v\ge2$ once. Thus, conditional on entrance at size $v$, the expected residence time is at most $(v-1/2)^{-1}$ (no state-changing jump leaves the volume fixed). Consequently $\varrho_N(V=v)\le[(v-1/2)\log N]^{-1}$ by integrating the occupation probability up to $\log N$.

Fix a path satisfying \eqref{eq:forest-tv-input}. The observed set $S_n$ of volumes has cardinality at most $M_{n,N}$. For large $n$ put $v_n=1+\lceil M_{n,N}\log\log N/\log N\rceil$. The test event $E_n=\{V\in S_n,\ V>v_n\}$ depends only on volume. Its reference probability is at most $M_{n,N}/[(v_n-1/2)\log N]\le1/\log\log N$. Its empirical probability is $\nu_{n,N}(V>v_n)$. By the count premise, $\log v_n/\log N=\mu(\log n/\log N-1)+o(1)$, since $(n/N)^\mu\log\log N\to\infty$. The limsup of this ratio is $\mu(\alpha^{-1}-1)<1$. Apply the uniform CDF premise at the moving $x=\log v_n/\log N\in[0,1]$ eventually. Subtraction gives the claimed lower bound, valid for TV on volumes and hence also on shapes. At $\alpha=1$ the bound is one. This argument gives no TV convergence regime and no sharp transition; in particular a weak limit for normalized log-volume can coexist with TV distance tending to one.
\end{proof}

\sectag{NOTES}
\section{Results with narrower checking}\label{sec:further}
The results of this section are recorded with their review scopes. They are not theorems of the paper, and nothing in the paper or in Parts~I and~II depends on them.
This section records further results whose scope or checking is narrower than that of the numbered theorems. Each paragraph states what was checked and by whom; Section~\ref{sec:provenance} explains the terms. The \proofsrc{round2/LEDGER.md}{claim ledger} records the proofs, dependencies and review scopes. A review of a statement applies to the stated conclusion, not to its inputs, and a receipt check is not a new check of the proof that produced it.

\subsection{Tails and density of the amplitude}
Given the strong limit and the cited uniform mean and one-leaf influence estimates, the \proofsrc{round2/proofs/fable/newground/law-of-X/REPORT.md}{clock argument} with its \proofsrc{round2/proofs/fable/newground/law-of-X/referee/many-site-repair.md}{many-site repair} gives $\Prob(X\le x)\sim\kappa x^{1/\mu}$ at zero and two-sided order $-\log\Prob(X>x)\asymp x^{1/(1-\mu)}$ at infinity. In particular $\E X^{-q}$ is finite exactly for $0<q<1/\mu$ among positive $q$. Gemini checked these implications, taking the estimates as inputs.

The separate \proofsrc{round2/proofs/astra-late/source/round2/open/limit-law/right/README.md}{finite-block proof} gives the stronger equivalent $-\log\Prob(X>x)\sim I x^{1/(1-\mu)}$ with $0<I<\infty$. It uses the embedded amplitude martingale, its uniform conditional concentration, the finite-time free-tree one-leaf gradient, and two-sided right-tail order, not the left-tail prefactor. It extracts rare blocks of high eigenfunction density and compares their repetition with actual attachment histories. Gemini audited that argument, not its inherited gradient estimate. Neither tail constant is numerically certified.

With additional limit, tail and clock inputs, a repaired argument gives a bounded density with $f(x)\sim(\kappa/\mu)x^{1/\mu-1}$ at zero. Claude checked the density steps but not the detailed clock script, which GPT checked. See the \proofsrc{round2/proofs/fable/newground/law-of-X/density/REPORT.md}{density proof} and its \proofsrc{round2/proofs/fable/newground/law-of-X/density/bare-triple-repair.md}{clock repair}. The smoothness, sharp H\"older order and derivative equivalent for the repaired density have only author-family audits.

\subsection{Low-degree forests}
Delete vertices whose original degree exceeds $N$. At fixed $N$, the component sampled by choosing a surviving original non-leaf uniformly has an almost-sure total-variation limit and an exponential bound on its limiting core radius, not uniform in $N$. A uniformly chosen original vertex is an isolated survivor with probability tending to one. A different law counts each component containing an original non-leaf once, including singleton surviving cores. Let $n\to\infty$ at fixed $N$ and then $N\to\infty$. The limiting count divided by $C_n$ is asymptotic to $AN^{-\mu}\log N$, and the limiting law for counted rooted shapes approaches $\varrho_N$ in total variation. Here $A$ is the constant of Theorem~\ref{th:core-means} and $\varrho_N$ is the inert newborn at uniform age on $[0,\log N]$, as defined for Theorem~\ref{th:forest-tv-obstruction}. Gemini audited the sequential results given the spatial inputs.

Given those inputs, the sequential law, and the cited growth and spectral estimates, for each deterministic integer sequence $N=N(n)\to\infty$ with $N/n\to0$ the count is almost surely asymptotic to $AX(n/N)^\mu\log N$, and the empirical CDF of $\log V/\log N$ converges uniformly to the uniform CDF. Almost-sure rooted-shape convergence in total variation to $\varrho_N$ is proved for divergent cutoffs $N\le c_0\log n$, where $c_0$ depends on a spectral gap and has no certified numerical value. GPT audited these estimates. Theorem~\ref{th:forest-tv-obstruction} uses the count and CDF as inputs and gives breakdown when $\liminf \log N/\log n>\mu/(1+\mu)$, even for volume alone. None of these statements gives a sharp transition or the law of the entire deleted forest. See the \proofsrc{round2/proofs/astra-late/source/round2/open/forest-age/README.md}{sequential law} and \proofsrc{round2/proofs/neon2/q5/README.md}{simultaneous laws}.

The \proofsrc{round2/proofs/astra-late/source/round2/open/low-forest/marked-repair/README.md}{marked forest repair} compares with moving stationary laws with exact integer marks, not with the unmarked age mixture. It covers counted components for divergent $N\le(\log n)^\gamma$, $0<\gamma<1$, and sampling by surviving cores under an envelope $N_*\ge N$ with $N_*^2\log(N_*+1)=o(\log\log n)$. It does not restore the original surviving-core range. The \proofsrc{round2/proofs/fable/newground/q5-low-degree-forest/longest-path/REPORT.md}{longest degree-two path and fixed-cutoff maxima} estimates were not fully checked outside GPT. Neither has a full outside-family check: the marked forest repair has only author-family proof audits, and the longest-path estimates have only partial outside checks.

\subsection{Degree counts at growing degree}\label{sec:further-degrees}
Theorems~\ref{th:core-amplitude} and~\ref{th:core-means} let $n\to\infty$ before $k\to\infty$. Two supplementary results remove that order in weaker senses, and they should not be confused with each other or with the fixed-degree law.

\paragraph{Cumulative counts, almost surely.} The \proofsrc{round2/proofs/astra/source/round2/fringe/hub-spectrum.md}{supplementary hub-spectrum proof}, not reproduced here, gives an event of probability one on which $N_{\ge k(n)}(n)\sim AX(n/k(n))^\mu$ along every integer sequence $k(n)\to\infty$ with $k(n)/n\to0$. It uses the fixed-degree, amplitude and hub-mass results. A Hunyuan-family reviewer audited the ranked and intermediate cumulative limits. The Claude-family check covered finite statistics, not the asymptotic proof. This is a joint limit for cumulative counts, not for exact-degree counts.

\paragraph{Exact-degree counts, in mean.} From one edge, uniformly as $k\to\infty$ and $k/n\to0$,
\[
 \frac{\E N_{=k}(n)}{\mu A c_2\,n^\mu k^{-(1+\mu)}}\longrightarrow1,
 \qquad
 \frac{\E N_{\ge k}(n)}{A c_2\,n^\mu k^{-\mu}}\longrightarrow1,
\]
with $A$ and $c_2=\E X$ the constants of Theorems~\ref{th:core-means} and~\ref{th:core-amplitude}. The same holds from each fixed reachable seed with at least one core vertex, with $\E X$ replaced by its mean from that seed. The inputs are Theorems~\ref{th:exist}, \ref{th:core-amplitude} and~\ref{th:core-means} together with two elementary bounds from the author's earlier preprint \cite{UzayFirst}, which the review rederives: a first-moment bound on the number of cores of degree at least $d$, and an exponential moment of the core degree. The hub-spectrum notes are not used. The argument starts at fixed degrees, where Theorem~\ref{th:core-means} applies, and transfers the equivalent between scales along an iterated-logarithm tower. Each transfer compares the degree of every vertex in a band with a neutral two-colour P\'olya urn through a stopped likelihood ratio, and needs a tail bound for that urn. The bound printed in the frozen proof was false: it estimated a Beta ratio in the wrong direction, and the audit recorded an explicit counterexample. The erratum and the review replace it by a Gamma representation with relative tolerance $\eta/5$, which gives exponential tail bounds in both directions, and also correct two literal rounding inequalities in the transfer. Gemini wrote the argument and the erratum. GPT's audit records the counterexample, checks the repaired proof for every fixed reachable seed and adds exact rational controls. This is a statement about means only. It gives no almost-sure, in-probability or fluctuation law at growing degrees. See the \ordersrc{round3/q2-growing-degrees/CONTINUATION_DISCRETE_MEAN.md}{frozen argument}, \ordersrc{round3/q2-growing-degrees/ERRATUM_DISCRETE_MEAN.md}{erratum}, \ordersrc{round3/q2-growing-degrees/REVIEW.md}{review} and \ordersrc{round3/q2-growing-degrees/audit_exact.py}{exact controls}.

\subsection{Other comparisons by birth label}\label{sec:further-order}
Theorem~\ref{th:stoch-order} contains the following earlier conclusions, which are kept as an independent route for the first labels and are not used as inputs. Transport certificates compare the marginal survival functions strictly at every $0\le x<1$ for the adjacent pairs $(2,3)$ through $(19,20)$. They use ordered copies, not the two labels on a common tree. Gemini independently reconstructed the smaller flows and checked the larger pairs only in the scopes recorded in the \proofsrc{round2/reviews/attacks/neon-on-fable-5.md}{transport audit}. The \proofsrc{round2/proofs/astra-late/source/round2/open/age-order/README.md}{positive-threshold comparisons}, including the common band $[1/10,1)$ with early transport inputs and eventual order at every threshold for sufficiently late adjacent labels, do not have a complete outside audit.

Conditioning on a particular reachable tree \emph{after} both births reverses a hub-probability comparison; Claude checked this \proofsrc{round2/reviews/attacks/fable-on-astra-18.md}{reversal}, which does not contradict the comparison before both placements.

\subsection{Arithmetic reductions and numerical evidence}
Claude checked the \proofsrc{round2/proofs/astra-late/source/round2/occupation-quintet-gold2/rationals/repaired/README.md}{stationary reduction, finite potential ceiling and Farey arithmetic}, taking the Gold2 certificate and four further physical reward floors as premises. With those floors this route also gives $99/175<\mu<142/251$. Only the author family audited the derivation of the floors, and a check of the conditional reduction does not check them; the narrower corrected Gold2 repricing has the same limitation. Theorem~\ref{th:gold2} obtains the same exclusion by a different route and does not use this one.

In the same draft that proved constant-factor mean growth, Long introduced the signed renewal germ $B_s(z)=1-1/F_s(z)$ at $z=0$ and proved its critical boundary value $\lim_{z\uparrow1}B_\mu(z)=1$ along the real axis \cite{LongMean}. The comparison here evaluates his germ at $z=1$ for real $s>\mu$, by continuation along $[0,1]$. Here $F_s(1)$ is the sum of the discounted generation counts from the inert newborn. The birth times partitioned by generation include the initial core at time zero. By Tonelli, each birth at time $u$ contributes its discounted occupation $e^{-su}/s$ to the total expected core occupation $M_s(J)$, so $F_s(1)=sM_s(J)$. The \proofsrc{round2/proofs/astra/source/round2/influence/zero/critical.md}{signed resolvent identity} gives $sM_s(J)(1-U_s(J))=1$ above criticality. Thus $B_s(1)=U_s(J)$ there, without identifying the two operators or asserting convergence of the signed Taylor series at one. Using the linked resolvent identity as input, a separate GPT-authored \proofsrc{round2/proofs/neon2/obstruction/REPORT.md}{prime-denominator argument}, checked by Claude, shows that this function is transcendental over $\mathbb C(s)$. It tests denominator primes in the expansion at large $s$, not arithmetic specialization at the critical discount. It does not show that the number $\mu$ is irrational or transcendental, and it is a statement about a function of the discount, not about the sequence $\E C_n$ of Theorem~\ref{M-th:nonholonomic}. Separate arguments for transcendence of the physical trigger and of each marked renewal entry over $\mathbb C(s)$ use their own resolvents and have only author-family proof audits.

The Claude-led simulations in this project, separate from Sabba's estimate, gave two estimates near $0.56572366$. They used 19,000 realizations for the time window $[16,24]$ and 2,840 for $[24,28]$, starting from the three-vertex path at time zero. These are fixed-time simulations with random final sizes, not fits at a prescribed maximum size. GPT independently reduced the raw runs. Sampling standard errors apply to the finite-window ratios, not to $\mu$. The \proofsrc{round2/proofs/fable/CLAIM.md}{public replay} uses packaged moments, not the raw runs. Inverse-symbolic searches found no relation distinguishable from chance in the searched families and heights; neither they nor the transcendence of individual entries excludes any arithmetic class for $\mu$. The finite-window bias bounds assume unproved contractions, and an eight-digit interpretation was withdrawn. An exact root count for a solvable variant gives no pole-free strip for the friend tree. The published finite-window fit $0.569$ \cite[Fig.~6]{HKRZ} lies outside the certified interval; an independent-size fit experiment shows how that number could arise but does not reconstruct the original runs. See the \proofsrc{round2/proofs/fable/README.md}{searches and numerical controls}.

\sectag{PROV}
\section{Prior work and provenance}\label{sec:provenance}\label{sec:prov}
\subsection{Prior work}
The one-step rule belongs to the random-walk attachment models studied by Saram\"aki and Kaski \cite{SK}. Krapivsky and Redner studied it as isotropic redirection. In their terminology the non-leaves form the nucleus. Their fitted slopes ranged from $0.5652$ to $0.5673$, and they reported $0.566$ \cite[Sec.~5.1]{KR}. In a 2022 post relaying Redner's account, Strogatz wrote, ``It's an open problem to prove that $17/30$ is really the exponent'' \cite{Strogatz}. That fraction is within the range of the fitted slopes. The three-vertex certificate in Theorem~\ref{th:exist} already places the exponent below it. Hartle, Krapivsky, Redner and Zhang later used the word nucleus for a smaller set. They called all non-leaves the core and obtained a finite-size fit of $0.569$ for their count under the same rule \cite[Fig.~6]{HKRZ}. Cannings and Jordan had shown that the fraction of leaves tends to one \cite{CJ}. Addario-Berry, Briend, Devroye, Donderwinkel, Kerriou and Lugosi named the random friend tree, proved bounds for core growth and proved existence of the fixed-label degree limits \cite{AB}. Their Section~7 asks whether the gap between the growth bounds closes and whether the normalized counts have a nontrivial random limit (Question~1), for a degree law (Question~2), about birth order (Question~3), about conservation of hub mass (Question~4) and about the shapes left after deleting high degrees (Question~5).

Christopher D. Long first characterized the exponent by the critical crossing of discounted generation sums for an inert newborn. He gave a terminating two-sided approximation and a certified upper bound of $0.5664485422$, also below $17/30$ \cite{LongKernel}. Long subsequently proved $\E C_n\asymp n^\mu$ \cite{LongMean}. From one edge, his signed hypergeometric product has expectation one at every tree size, and its linear coefficient is a martingale with variance of order $n^\mu$ \cite{LongProduct}. His notes also study an exact beta--Mellin series whose rational coefficients are the means $\E C_n$, and state that no finite scalar recurrence or finite hypergeometric equation for it is claimed \cite{LongRepo}. Theorem~\ref{th:core-kernel} uses a compact operator on newborn marks. Mohamed Sabba reported the interval $[0.565615,0.565809]$ and a numerical estimate of $0.5657237$ \cite{Sabba}.

\subsection{Who wrote and who checked each proof}
Can Uzay chose the problem, directed the work and set the evidence criteria; he is responsible for the manuscript. Generative AI models from three families, GPT (OpenAI), Gemini (Google) and Claude (Anthropic), developed the arguments, wrote the programs and the text, and performed the reviews. No human referee has reviewed any of the proofs, and the reviews described here are not peer review. A result \emph{checked by another family} was audited, in a separate session, by a model of a family other than its author's. The \proofsrc{round2/who-did-what.md}{credit record} and the \proofsrc{round2/reviews/README.md}{review archive} distinguish the authors of proofs, programs and reviews.

GPT wrote the main renewal, strong-limit, degree, hub-mass and all-label proofs, and Gemini and Claude reviewed them within the ledger's scopes; for example, Gemini independently coded a rational two-birth enumeration for the hub-order proof, and Claude recomputed the first fourteen signed-response coefficients with its own exact program (\proofsrc{round2/reviews/attacks/neon-on-astra-26.md}{hub-order audit}, \proofsrc{round2/reviews/attacks/fable-on-neon2-4.md}{response audit}). Gemini wrote the five-vertex certificate, the original Beta-split density proof and the simultaneous forest arguments. GPT refined the density proof, and Claude independently rederived its split and side-limit argument. GPT and Claude replayed the authored five-vertex verifier; neither independently implemented its full row scan. Both checked the conditional forest breakdown; outside Gemini, GPT alone audited the full simultaneous count and shape-convergence arguments. GPT wrote the supplementary tail, bounded-density and regularity proofs, under more than one model lead. Gemini wrote the proof of Theorem~\ref{th:stoch-order}, its two integer certificates and their controls. GPT audited the full argument and, separately, the exact certificate bytes, with its own assertions on the certificate tables, independent rational comparisons of the comparison laws, checked-overflow runs and direct signed sums over every unlabelled old tree with at most twelve vertices. That audit found two false auxiliary statements and an omitted convention, all local and repaired as recorded in Section~\ref{app:stoch-order}. The same reviewer had earlier developed a different route to the same claim, which its review discloses.

Claude wrote the proof of Theorem~\ref{M-th:nonholonomic} and Lemma~\ref{lem:classwise}, and built the refined lower table with the Gemini-written builder. A GPT referee checked the frozen arguments and found that both hold within their stated scopes, taking the five-vertex reduction and the soundness of its verifier as inputs for the refined endpoint: for Theorem~\ref{M-th:nonholonomic} the classification, the grouped transitions, the Wilson step, the residue bookkeeping and the recurrence obstruction; for the refined endpoint the classwise lemma, the host allocation, the tail classes, the exact endpoint and the Farey bound. A separate Claude review of the classwise accounting also passed. The witness was replayed in full, once normally and once with the undefined-behaviour sanitizer in non-recovering mode, and the builder regenerated its bytes exactly. Thus Theorems~\ref{th:stoch-order} and~\ref{M-th:nonholonomic} and the refined lower endpoint have each been checked by one model family other than the author's. Claude prepared this manuscript from the frozen proofs and reviews without changing their mathematics, certificates or constants; the sources are pinned in Sections~\ref{app:stoch-order} and~\ref{app:replay}.

\sectag{VER}
\section{Verification packages and sources}\label{app:replay}\label{sec:replay}
All paths are relative to the repository root. The light checks need Python~3.10 or newer with its standard library and \texttt{g++} with C++17; the full replay also needs OpenMP and the undefined-behaviour sanitizer.

\paragraph{Light checks.} \texttt{bash verify/compact\_replay.sh} runs, in order: the exact checks of Appendix~\ref{M-app:nonholonomic} of the paper (\texttt{verify/mean-arithmetic/check\_window.py}); the receipt arithmetic of Section~\ref{app:refined} with planted defects (\texttt{verify/refined-lower/check\_ratio.py -{}-self-test}); and both stochastic-order certificates of Section~\ref{app:stoch-order} with their primitive audits and planted-defect controls. The second script authenticates the receipt, the per-class file and the Gold2 receipt by SHA-256, requires an accepted lower-mode receipt for the pinned witness digest, parses every class row $c=1,\dots,4095$, recomputes $\kappa$ and $a_-$ exactly, and checks the Farey arithmetic and the outward decimals. Eight planted defects, among them a changed class cap and a missing or duplicated row, must each be rejected. It does not rerun the drift verifier.

\paragraph{Full replay of the refined lower table.} The witness is a single file of 13{,}046{,}351{,}000 bytes with SHA-256 \texttt{80e3b80a9f962b458854828abab5b00e2806a6c1ce709835b5e884b0afd4ad28}. It is not stored in the repository. The witness was regenerated byte for byte by running the repository's five-vertex builder with its \texttt{gold3} lower preset at eight threads, using the thread-count patch in \texttt{verify/refined-lower} (the only change), and then setting the integer header numerator to $302884234$, giving $\alpha=302884234/2^{30}$. \texttt{bash verify/refined-lower/replay\_full.sh WITNESS NEW\_DIR} checks the witness digest and header, patches a copy of the verifier with the four instrumentation lines and checks the patched digest, and compiles with the undefined-behaviour sanitizer in non-recovering mode after a planted overflow control. It then requires the instrumented run to exit cleanly with an identical verification record and a per-class file byte-identical to the stored one, and requires the unmodified checker on the unmodified verifier to reproduce the stored lower receipt exactly. It needs about 14 GB of memory.

\paragraph{Sources.} The results of Appendix~\ref{M-app:nonholonomic} of the paper and Section~\ref{app:refined} were first written as research notes, which are not part of the repository. The reviews in Section~\ref{sec:provenance} apply to these exact bytes:
\begin{center}\footnotesize
\begin{tabular}{@{}ll@{}}
\toprule
Item & SHA-256\\
\midrule
window law and near-star chain & \texttt{619969f1b929ac1c8d75fc8764954f6dc6f42988ab59f767d458e1b0b4764df2}\\
near-star residues, $p\equiv2\ (3)$ & \texttt{23a8b1fc537bab71f96abe3132c0cf7442d24406e986ae422aa7f3d33aac0138}\\
alternative residue derivation & \texttt{0fc654a9c2fcd1771b5698d7929150da6baa7af5c963c255501d2c8b872bcbf6}\\
classwise aggregation & \texttt{9e08e866a13156df08909faa37277b1c0b99415d870b1c191955c6e2ae169f70}\\
refined lower witness & \texttt{80e3b80a9f962b458854828abab5b00e2806a6c1ce709835b5e884b0afd4ad28}\\
lower receipt & \texttt{f46b5f8dd7b73bfec3d91fd80c83b9fa4e1d96f326a7f4581175ed41fb44dc17}\\
per-class data & \texttt{278c423e76650feff8766405c4ab766152fa1df2aed18cffb93d5bab6db6a612}\\
instrumented verifier & \texttt{f85a172ab45195d0c9d49d809f0c903c146f0bbc570cb170b5887b493954f067}\\
\bottomrule
\end{tabular}
\end{center}
Appendix~\ref{M-app:nonholonomic} of the paper and Section~\ref{app:refined} are new write-ups of those notes; they keep the stated scope and add no new mathematical step. One remark in the alternative derivation calls a constant $5269/129600$ a unit for every prime in range, although $5269=11\cdot479$; that derivation uses only its integrality, and the proof given here does not use it.

\begin{thebibliography}{99}
\raggedright
\bibitem{SK} J.~Saram\"aki and K.~Kaski, Scale-free networks generated by random walkers, \emph{Physica A} \textbf{341} (2004), 80--86. \href{https://doi.org/10.1016/j.physa.2004.04.110}{doi:10.1016/j.physa.2004.04.110}.
\bibitem{CJ} C.~Cannings and J.~Jordan, Random walk attachment graphs,
\emph{Electronic Communications in Probability} \textbf{18} (2013),
paper 77, 1--5. \href{https://doi.org/10.1214/ECP.v18-2518}{doi:10.1214/ECP.v18-2518}.
\bibitem{KR} P.~L.~Krapivsky and S.~Redner, Emergent network modularity,
\emph{Journal of Statistical Mechanics: Theory and Experiment} (2017),
073405. \href{https://doi.org/10.1088/1742-5468/aa7a3f}{doi:10.1088/1742-5468/aa7a3f}. \href{https://arxiv.org/abs/1706.01514}{arXiv:1706.01514}.
\bibitem{Strogatz} S.~Strogatz, post on the exponent in isotropic
redirection, 20 September 2022.
\url{https://x.com/stevenstrogatz/status/1572287647715934210}.
\bibitem{AB} L.~Addario-Berry, S.~Briend, L.~Devroye, S.~Donderwinkel,
C.~Kerriou and G.~Lugosi, Random friend trees, \emph{Random Structures \& Algorithms} \textbf{69} (2026), e70082. \href{https://doi.org/10.1002/rsa.70082}{doi:10.1002/rsa.70082}.
\href{https://arxiv.org/abs/2403.20185v1}{arXiv:2403.20185v1} (numbering used here).
\bibitem{HKRZ} H.~Hartle, P.~L.~Krapivsky, S.~Redner and Y.~Zhang,
Anomalous scaling in redirection networks, \emph{Physical Review E}
\textbf{114} (2026), 034314.
\href{https://doi.org/10.1103/9lmp-pydp}{doi:10.1103/9lmp-pydp}. \href{https://arxiv.org/abs/2604.01540}{arXiv:2604.01540}.
\bibitem{LongKernel} C.~D.~Long, \emph{Effective computation and boundary robustness of the random friend-tree exponent}, 23 November 2026. \href{https://github.com/long-mathematics/friend-tree-growth/blob/090bbbe5faec83ec4124bd427eb5fda32c79bef1/friend_tree_growth.pdf}{PDF on GitHub}, commit 090bbbe.
\bibitem{LongMean} C.~D.~Long, \emph{Transcendence targets for the friend-tree exponent: exact first-promotion reduction and arithmetic obstructions}, 24 November 2026. \href{https://github.com/long-mathematics/friend-tree-growth/blob/790491725d9cf445789e14c24c277532f5d2e5c7/friend_tree_transcendence_progress.pdf}{PDF on GitHub}, commit 7904917. \href{https://github.com/long-mathematics/friend-tree-growth}{Long's research repository}.
\bibitem{LongRepo} C.~D.~Long, \emph{Friend-tree growth}, research repository at commit 5bd25a7, November 2026, section ``Arithmetic and exact root representations''. \href{https://github.com/long-mathematics/friend-tree-growth/blob/5bd25a7d506cb64a8444b86334ab531d295f6c78/README.md#arithmetic-and-exact-root-representations}{README on GitHub}.
\bibitem{LongProduct} C.~D.~Long, \href{https://github.com/long-mathematics/friend-tree-growth/blob/5bd25a7d506cb64a8444b86334ab531d295f6c78/README.md#signed-hypergeometric-products-for-the-original-tree}{\emph{Signed hypergeometric products for the original tree}}, in the \emph{Friend-tree growth} repository at commit 5bd25a7, November 2026.
\bibitem{Stanley} R.~P.~Stanley, Differentiably finite power series, \emph{European Journal of Combinatorics} \textbf{1} (1980), 175--188. \href{https://doi.org/10.1016/S0195-6698(80)80051-5}{doi:10.1016/S0195-6698(80)80051-5}.
\bibitem{EC2} R.~P.~Stanley, \emph{Enumerative Combinatorics}, Volume~2, Cambridge Studies in Advanced Mathematics~62, Cambridge University Press, 1999.
\bibitem{Main} C.~Uzay, \emph{The non-leaf growth exponent of the random friend tree}, 25 November 2026. Main paper \href{friend-tree-exponent.pdf}{\texttt{friend-tree-exponent.pdf}}, distributed with this supplement.
\bibitem{UzayFirst} C.~Uzay, \emph{The non-leaf growth exponent of the random friend tree}, preprint, 23 November 2026. \proofsrc{paper/friend-tree-exponent.pdf}{archived PDF}, commit a097208.
\bibitem{Sabba} M.~Sabba, \href{https://x.com/ma_sabba/status/2124768206141206762}{Interval for the friend-tree exponent} and \href{https://x.com/ma_sabba/status/2124791407667716208}{numerical estimate}, replies in S.~Strogatz's thread, November 2026.
\end{thebibliography}
\end{document}
